% ============================================================
% MyQCD/PyQCD 格点 QCD 物理理论审计报告
% 编译：xelatex -interaction=nonstopmode -halt-on-error 本文件（至少两遍）
% ============================================================
\documentclass[UTF8,11pt,oneside]{ctexbook}

\usepackage[a4paper,margin=2.45cm,headheight=15pt]{geometry}
\usepackage{amsmath,amssymb,mathtools}
\usepackage{booktabs,longtable,tabularx,array}
\usepackage{multirow}
\usepackage{enumitem}
\usepackage{xcolor}
\usepackage{listings}
\usepackage{url}
\usepackage[hidelinks]{hyperref}
\usepackage{fancyhdr}
\usepackage{microtype}

\linespread{1.08}
\setlength{\parskip}{0.32\baselineskip}
\setlength{\parindent}{2em}
\setlength{\emergencystretch}{3em}
\setlist{nosep,leftmargin=2.2em}
\setlist[description]{style=nextline,leftmargin=3.4em,labelindent=0em}
\newcolumntype{Y}{>{\raggedright\arraybackslash}X}
\newcolumntype{P}[1]{>{\raggedright\arraybackslash}p{#1}}

\definecolor{primary}{HTML}{17365D}
\definecolor{evidence}{HTML}{1E8449}
\definecolor{inferred}{HTML}{1F77B4}
\definecolor{warning}{HTML}{C55A11}
\definecolor{muted}{HTML}{5B6573}
\definecolor{panel}{HTML}{F4F7FA}
\definecolor{rulegray}{HTML}{CBD5E1}

\pagestyle{fancy}
\fancyhf{}
\fancyhead[L]{MyQCD/PyQCD 格点 QCD 物理理论审计}
\fancyhead[R]{\leftmark}
\fancyfoot[C]{\thepage}
\hypersetup{pdftitle={MyQCD/PyQCD 格点 QCD 物理理论审计报告},
  pdfauthor={Codex},colorlinks=true,linkcolor=primary,urlcolor=primary}

% 按普通文本排 CJK 路径，并仅在路径分隔符处插入可选断行点。
% 不使用 verbatim URL 字体，以免中文参考名丢字。
\ExplSyntaxOn
\NewDocumentCommand{\fileline}{m}
  {
    \group_begin:
    \normalfont
    \str_map_inline:nn {#1}
      {
        ##1
        \str_case:nn {##1}
          {{/}{\allowbreak}{_}{\allowbreak}{.}{\allowbreak}{-}{\allowbreak}}
      }
    \group_end:
  }
\ExplSyntaxOff
\ExplSyntaxOn
\NewDocumentCommand{\codeid}{m}
  {
    \group_begin:
    \ttfamily
    \str_map_inline:nn {#1}
      {##1\str_case:nn {##1}{{_}{\allowbreak}{/}{\allowbreak}{.}{\allowbreak}}}
    \group_end:
  }
\ExplSyntaxOff
\newcommand{\src}[1]{\par\begingroup\scriptsize\raggedright\normalfont 来源：\fileline{#1}\par\endgroup}
\newcommand{\verified}{\textcolor{evidence}{\textbf{确证}}}
\newcommand{\structural}{\textcolor{inferred}{\textbf{结构性}}}
\newcommand{\inferred}{\textcolor{inferred}{\textbf{推断}}}
\newcommand{\unverified}{\textcolor{warning}{\textbf{未验证}}}
\newcommand{\Tr}{\operatorname{Tr}}
\newcommand{\tr}{\operatorname{tr}}
\newcommand{\dd}{\mathrm{d}}
\newcommand{\e}{\mathrm{e}}
\newcommand{\ii}{\mathrm{i}}
\newcommand{\Nc}{N_{\mathrm c}}
\newcommand{\Ns}{N_{\mathrm s}}
\newcommand{\Oflow}{\mathcal O_{\mathrm{flow}}}
\newcommand{\Obare}{\mathcal O_{\mathrm{bare}}}
\newcommand{\Oren}{\mathcal O_{\mathrm R}}
\newcommand{\qtilde}{\widetilde q}
\newcommand{\gtilde}{\widetilde g}
\newcommand{\Ttilde}{\widetilde T}

\lstset{
  basicstyle=\ttfamily\small,
  backgroundcolor=\color{panel},
  frame=single,
  rulecolor=\color{rulegray},
  breaklines=true,
  breakatwhitespace=false,
  keepspaces=true,
  columns=fullflexible,
  showstringspaces=false
}

\begin{document}

\frontmatter
\begin{titlepage}
  \centering
  \vspace*{1.5cm}
  {\Huge\bfseries MyQCD/PyQCD 格点 QCD 物理理论审计报告\par}
  \vspace{0.9cm}
  {\LARGE 从规范组态、蒸馏与 Wick 收缩到 OPE、梯度流、准 PDF 与 TMD\par}
  \vspace{1.2cm}
  {\large 代码—公式—物理对象—证据边界的交叉复现\par}
  \vfill
  \begin{tabular}{rl}
    工作目录：& \fileline{/Users/zhangxin/MyQCD}\\
    对照项目：& \fileline{/Users/zhangxin/PyQCD}\\
    默认数据目录：& \fileline{/Users/zhangxin/MyQCD/data}\\
    报告日期：& 2026 年 9 月 6 日\\
    复现工具：& Python、SymPy、XeLaTeX
  \end{tabular}
  \vfill
  {\small
  本报告将有限维代数验证、源码接口审阅、文献公式和真实系综验证严格分层。
  \par}
\end{titlepage}

\chapter*{摘要}
\addcontentsline{toc}{chapter}{摘要}
本报告审计工作目录 \fileline{MyQCD} 与对照项目 \fileline{PyQCD} 中格点
QCD 物理理论的连续链条：
\begin{align*}
 &\text{规范组态}\longrightarrow\text{链接涂抹}
 \longrightarrow\text{Laplacian 本征矢/蒸馏}
 \longrightarrow\text{VdV/VVV}\\
 &\longrightarrow\text{perambulator 与 Wick 收缩}
 \longrightarrow C_2,C_3,\mathrm{OPE}
 \longrightarrow\text{统计与比值}
 \longrightarrow \Obare\longrightarrow \Oren\\
 &\longrightarrow\text{Fourier、matching、PDF/TMD}.
\end{align*}
主结论是：对照项目 PyQCD 的源码和文档足以固定一组接口、轴序与算符约定，
工作项目 MyQCD 则对其中可在有限维或形式代数层闭合的关系进行了局部复现。
这支持“公式结构和局部数值管线可审计”的判断；它不支持把当前工作树
直接称为“某个真实系综上的完整 TMD-PDF 结果”。特别是三点函数、比值、
裸矩阵元、完整重整化/matching、soft/rapidity 因子和最终物理外推，仍须
绑定真实数据、方案、误差和共同几何的验证门。

\begin{table}[htbp]
\centering
\caption{本次交付的证据分级}
\small
\begin{tabularx}{0.96\textwidth}{@{}P{2.4cm}Y Y@{}}
\toprule
状态 & 可以支持的结论 & 不能越界声称的结论\\
\midrule
\verified & SymPy 精确代数、明确有限维模型、源码接口和本次命令输出直接支持。
 & 不自动覆盖真实系综、GPU 精度、统计相关性或非微扰动力学。\\
\structural & 公式链、参数接口和代数拼接闭合，物理对象关系明确。
 & 不等于方案系数、圈积分、拟合结果或实验可比量已完成。\\
\inferred & 由多个源码和公式证据合理推出的物理解释。
 & 需要额外数值/文献条件才能升级为确证。\\
\unverified & 数据、完整积分、软因子、rapidity 因子或生产配置尚未闭合。
 & 不得用接口存在、示意数据或单元测试替代该验证。\\
\bottomrule
\end{tabularx}
\end{table}

\chapter*{如何阅读与复现}
\addcontentsline{toc}{chapter}{如何阅读与复现}
正文采用“物理图像 $\to$ 定义 $\to$ 方程 $\to$ 算法 $\to$ 极限/量纲检查
$\to$ 源码与文献”的顺序。文件引用使用真实相对路径和行号；其中
\fileline{../PyQCD/...} 是相对于本报告所在 \fileline{MyQCD/docs/} 的对照路径。
论文优先引用 \fileline{MyQCD/refer/papers/} 的中文/英文转排源，教材引用
\fileline{MyQCD/refer/books/}；PyQCD 的中文文档和源码只承担实现接口与约定
的对照证据。

核心复现命令为：
\begin{lstlisting}
python -m myqcd
python -m myqcd.pyqcd_theory_audit \
  --output data/pyqcd_theory_audit/theory_audit.json
pytest -q
\end{lstlisting}
第一条运行既有 MyQCD 核心推导和 50 篇论文显示公式索引；第二条运行本报告
新增的九个独立理论审计并落盘 JSON。由于本任务性质为分析与汇报，本报告
不修改 PyQCD 核心源代码，也不默认读取工作目录之外的数据。

\tableofcontents
\mainmatter

\chapter{任务、证据和十五步交叉分析框架}

\section{任务边界与基本判断}
本任务的工作对象是 \fileline{MyQCD} 内已有的完整公式链、SymPy 推导和
报告，以及用户明确引入的对照项目 \fileline{PyQCD} 的 \fileline{pyqcd/}、
\fileline{docs/}、\fileline{refer/papers/} 与 \fileline{refer/books/}。
这里的“充分理解”不应简化为列出函数名，而要回答四个问题：
\begin{enumerate}
  \item 输入张量和物理状态是什么，轴序、单位、边界与颜色/Dirac 指标如何约定；
  \item 中间量对应哪一个连续或格点物理对象，哪些是辅助变量、哪些是可观测量；
  \item 算法由哪条物理方程、定理或对称性推出，停止条件和误差项是什么；
  \item 当前证据究竟到达“接口、代数、方案、真实数据”哪一层。
\end{enumerate}
因此，报告将“代码能调用”与“物理结论已验证”分开。尤其要区分
\[
\text{局部代数恒等式}\ne\text{非微扰 QCD 定理的完整证明}
\ne\text{真实格点数据的数值验证}.
\]

\section{十五步交叉分析框架}
下表是本报告实际采用的分析轴。每一步都能独立指向源码、公式或验证记录，
同时由前一步的输出约束后一步的物理解释。

\begin{longtable}{@{}P{0.08\textwidth}P{0.22\textwidth}P{0.34\textwidth}P{0.25\textwidth}@{}}
\caption{十五步交叉分析框架：从输入约定到证据闭环}\label{tab:fifteen-steps}\\
\toprule
编号 & 分析轴 & 需要回答的物理/工程问题 & 主要证据\\
\midrule
\endfirsthead
\caption[]{十五步交叉分析框架（续）}\\
\toprule
编号 & 分析轴 & 需要回答的物理/工程问题 & 主要证据\\
\midrule
\endhead
\bottomrule
\endfoot
1 & 任务定义 & 目标量是强子谱、矩阵元、PDF 还是 TMD？观测量的方案边界是什么？ & 用户规格、报告范围\\
2 & 约定固定 & \((N_t,N_z,N_y,N_x,4,3,3)\)、方向、自然单位、精度和边界是否唯一？ & \fileline{PyQCD/pyqcd/lattice/_constants.py}、\fileline{pipeline/_config.py}\\
3 & 物理公理 & 局域规范不变性、Euclidean 反射正性、Grassmann 高斯积分、平移不变性等假设在哪里使用？ & 第 \ref{sec:axioms} 节\\
4 & 连续—格点映射 & \(A_\mu,F_{\mu\nu},D_\mu\) 怎样离散为链接、plaquette 和 Clover？ & 第 \ref{sec:lattice-map} 节\\
5 & UV 预处理 & HYP、Stout、动量涂抹和蒸馏分别抑制什么高频结构，是否改变目标算符？ & 第 \ref{sec:smearing} 节\\
6 & 低秩表示 & Laplacian 本征矢与 \(\Box=VV^\dagger\) 如何把空间传播约化到蒸馏子空间？ & 第 \ref{sec:distillation} 节\\
7 & 算符元素 & VdV/VVV、Wilson 线和场强插入的颜色、动量、端点怎样闭合？ & 第 \ref{sec:vertices}、\ref{sec:ope} 节\\
8 & 费米子收缩 & perambulator、\(\gamma_5\)-Hermiticity、Wick 拓扑与费米符号怎样生成 \(C_2,C_3\)？ & 第 \ref{sec:wick} 节\\
9 & 谱学提取 & 传递矩阵谱分解、有效质量和 GEVP 如何隔离基态？ & 第 \ref{sec:spectrum} 节\\
10 & 统计推断 & Jackknife/bootstrap、协方差、自相关和拟合窗口如何进入误差？ & 第 \ref{sec:statistics} 节\\
11 & 矩阵元归一 & ratio 与裸矩阵元的 \(2E\)、投影、真空扣除和方向平均是什么？ & 第 \ref{sec:matrix} 节\\
12 & 重整化 & 局域/非局域线性发散、\(Z_R\)、RI/MOM、hybrid 与 flow 方案分别消除什么？ & 第 \ref{sec:renorm} 节\\
13 & 因子化匹配 & OPE、LaMET、Fourier、matching 核的尺度层级和卷积测度是什么？ & 第 \ref{sec:matching} 节\\
14 & TMD 几何 & \(b_T\)、finite staple、soft factor、rapidity 与 Collins--Soper 核是否全部闭合？ & 第 \ref{sec:tmd} 节\\
15 & 证据闸门 & 哪些是确证、结构性、推断、未验证？下一道可验收测试是什么？ & JSON 审计、\ref{sec:boundaries} 节\\
\end{longtable}

\section{证据层级、源码行号与数据边界}
本报告的代码引用不是把 docstring 当作物理定理，而是先检查实现的输入、
中间量和输出，再用独立公式审计其局部不变量。例如：
\begin{itemize}
  \item \fileline{vertex/_vertex.py:134--234} 明确给出 VdV/VVV 的张量收缩；
  \item \fileline{analysis/_analyse.py:91--162} 给出 leave-one-out Jackknife 和
  协方差因子；
  \item \fileline{operator/_gluon_ope.py:492--569} 给出 Clover 和对偶场强；
  \item \fileline{renorm/_tmd.py:71--195} 给出 staple 与双场强颜色迹；
  \item \fileline{renorm/_tmdextract.py:30--367} 给出 Fourier、CS、hybrid
  TMD 和 SFTX 的接口；
  \item \fileline{pipeline/_config.py:32--38} 的输入路径是 \fileline{/public/...}
  集群路径，因此本地没有这些文件时不能宣称完整生产管线可直接执行。
\end{itemize}
新增审计 JSON 还会逐条列出公式、假设、检查项、代码引用、文献引用和
缺失来源。JSON 中的 \texttt{unverified} 是证据状态，不是 Python 异常。

\chapter{物理公理、群论约定与连续—格点映射}

\section{使用的物理公理与工作假设}
\label{sec:axioms}
格点 QCD 的数值链不是从一段程序自动“推出”所有物理结论，而是在以下
公理/假设下把连续理论变成有限数组问题：
\begin{description}
  \item[局域规范对称性] 对任意 \(G(x)\in SU(3)\)，物质场和连接按局部变换
  \(\psi(x)\mapsto G(x)\psi(x)\)、\(A_\mu\mapsto G A_\mu G^\dagger+\cdots\)。
  这要求可观测量颜色指标最终闭合。
  \item[Euclidean 路径积分] Wick 旋转后以 \(Z^{-1}\exp(-S_E)\) 对组态平均；
  费米子是 Grassmann 变量，积分给出 determinant 与 Wick 收缩。
  \item[反射正性/传递矩阵] 在适当局域、正性和边界条件下，Euclidean 时间相关函数
  具有非负能谱分解，长时间极限由最低能态支配。
  \item[平移和周期边界] 空间求和可做动量投影，时间方向通常周期或反周期；
  \(\mathrm{roll}\) 的数值实现只有在边界约定与物理路径一致时才是正确的。
  \item[连续极限/重整化] \(a\to0\) 时通过调节裸参数和乘法/混合重整化因子，
  让物理矩阵元在固定方案和尺度下有限。
  \item[因子化层级] LaMET、OPE 和 TMD 公式需要 \(P^z\)、短距离、\(b_T\)、
  流时间和 \(\Lambda_{\rm QCD}\) 之间存在可控窗口；没有窗口时只能报告原型。
\end{description}
这些假设分别服务于不同步骤，不能把“采用了公理”误写成“已用代码证明公理”。

\section{QCD Euclidean 作用量与颜色归一}
取 \(T^a\) 为 Hermitian 的 \(SU(\Nc)\) 生成元，
\[
 \operatorname{tr}(T^aT^b)=\frac12\delta^{ab},\qquad
 [T^a,T^b]=\ii f^{abc}T^c,
\]
连续 Euclidean QCD 的形式作用量可写为
\begin{align}
 S_E[A,\bar\psi,\psi]
 &=\int\dd^4x\left[
 \frac14F^a_{\mu\nu}F^a_{\mu\nu}
 +\bar\psi(\gamma_\mu D_\mu+m)\psi\right],\\
 D_\mu&=\partial_\mu+\ii g A_\mu,\qquad
 F_{\mu\nu}=\partial_\mu A_\nu-\partial_\nu A_\mu
 +\ii g[A_\mu,A_\nu].
\end{align}
在自然单位中 \([x]=-1\)、\([A_\mu]=1\)、\([F_{\mu\nu}]=2\)、
\([m]=[P_\mu]=1\)。这给出后文的量纲检查：
\[
 [t]=-2,\qquad [E(t)]=4,\qquad [t^2E(t)]=0,\qquad
 [zP^z]=0.
\]
MyQCD 的 \fileline{derive_generating_functional} 在一维有限维高斯模型中
验证了源导数关系；这正是 Grassmann/玻色路径积分的一种可计算代理，
不是一般相互作用 QCD 路径积分的闭式评价。

\section{链接变量、plaquette 和 Wilson 作用量}
\label{sec:lattice-map}
把中点连接 \(x\to x+a\hat\mu\) 编码成群元素：
\begin{align}
 U_\mu(x)&=\mathcal P\exp\!\left[\ii g\int_x^{x+a\hat\mu}
 \dd s\,A_\mu(s)\right]
 \simeq \exp\!\left[\ii ag A_\mu(x+\tfrac a2\hat\mu)\right],\\
 U_\mu(x)&\mapsto G(x)U_\mu(x)G^\dagger(x+a\hat\mu),\\
 P_{\mu\nu}(x)&=U_\mu(x)U_\nu(x+\hat\mu)
 U_\mu^\dagger(x+\hat\nu)U_\nu^\dagger(x).
\end{align}
这里省略了每一个链接上的 \(a\) 以符合代码的整数格点偏移写法。闭合
plaquette 的端点变换相消，\(\Re\Tr P_{\mu\nu}\) 规范不变；Wilson 作用量
\[
 S_W[U]=\beta\sum_x\sum_{\mu<\nu}
 \left(1-\frac1{\Nc}\Re\Tr P_{\mu\nu}(x)\right),
 \qquad \beta=\frac{2\Nc}{g_0^2}.
\]
弱场展开 \(U_\mu=1+\ii agA_\mu-\frac12a^2g^2A_\mu^2+\cdots\) 后，
plaquette 的一阶项给 \(F_{\mu\nu}\)，从而恢复 \(\frac14F^2\)。这同时是
“量纲—对称性—连续极限”的第一道闭环。

\section{PyQCD 的张量约定}
\begin{table}[htbp]
\centering
\caption{工作对象的关键轴序、输入和输出约定}
\small
\begin{tabularx}{0.96\textwidth}{@{}P{3.1cm}P{4.2cm}Y@{}}
\toprule
对象 & 数组约定 & 物理含义/注意事项\\
\midrule
规范场 \(U\) & \((N_t,N_z,N_y,N_x,4,3,3)\) & 方向标签 \(0=x,1=y,2=z,3=t\)，坐标轴是 \(3-\mu\)。\\
本征矢 \(\phi\) & \((N_{\rm ev},N_z,N_y,N_x,N_c)\) & 每个时间片的空间 Laplacian 低模；颜色维最后。\\
perambulator & \((t,d_{\rm sink},d_{\rm src},e_{\rm sink},e_{\rm src})\) & \(\tau=V^\dagger D^{-1}V\) 的蒸馏子空间表示。\\
VdV & \((N_p,N_{\rm ev},N_{\rm ev})\) & 介子/双线性顶点，含 Fourier 相位。\\
VVV & \((N_p,N_{\rm ev},N_{\rm ev},N_{\rm ev})\) & 重子颜色 singlet 顶点，含 \(\epsilon_{abc}\)。\\
legacy OPE & \((\Delta z,N_t)\) & 直 Wilson 线、Clover \(F\widetilde F\) 组合；默认 finite-\(a\)、legacy、空间求和。\\
TMD OPE & \((N_z,N_b,N_t)\) 或逐时间片结果 & finite staple、非零 \(b_\perp\) 和无迹颜色场的双场强。\\
\bottomrule
\end{tabularx}
\end{table}

\section{本章可验证的极限与不变量}
\begin{enumerate}
  \item \(G(x)=1\) 时链接变换退化为恒等，开链保持原值、闭链为规范标量；
  \item \(a\to0\) 的链接展开中一阶为 \(\ii agA_\mu\)，二阶为
  \(-a^2g^2A_\mu^2/2\)，量纲为零；
  \item \(U_\mu^\dagger U_\mu=1\)、\(\det U_\mu=1\) 是 SU(3) 组态的输入不变量；
  \item 所有后续颜色指标必须通过迹、\(\delta_{ab}\) 或 \(\epsilon_{abc}\) 闭合；
\item 本次 SymPy 审计的 \texttt{wilson\_link\_endpoint} 对 U(1) 交换代理和
  SU(3) 非交换矩阵符号都验证了中间端点抵消。
\end{enumerate}


\chapter{涂抹、动量增强与蒸馏低秩表示}
\label{sec:smearing}

\section{为什么先处理 UV 模式}
格点链接含有 \(O(a^{-1})\) 附近的紫外涨落；局域算符、Wilson 线和
本征模都可能被这些涨落污染。涂抹的物理目的不是“把数据变平滑”这么
简单，而是在保持规范协变的前提下改变插值算符或链接的短程重叠：
\[
 U\ \mapsto\ U^{\rm smear},\qquad
 \psi\ \mapsto\psi^{\rm smear},\qquad
 \langle\mathcal O\rangle\ \text{的连续极限必须由重整化/匹配恢复}.
\]
HYP、Stout、Gaussian/Jacobi 和 momentum smearing 作用对象不同，不能
在报告或代码中只用一个“涂抹”名称混合它们。

\section{HYP 链接涂抹：三级排除方向算法}
\label{sec:hyp}
四维 HYP 的核心是每一级只使用没有被排除的方向。令
\(\mathcal P_{SU(3)}\) 表示由极分解和行列式归一得到的群投影，则
\begin{align}
 \bar V_{\mu;\nu\rho}(x)
 &=\mathcal P_{SU(3)}
 \left[(1-\alpha_3)U_\mu(x)
 +\frac{\alpha_3}{2}\sum_{\pm\eta\notin\{\mu,\nu,\rho\}}
 S_\eta[U_\eta,U_\mu](x)\right],\\
 \widetilde V_{\mu;\nu}(x)
 &=\mathcal P_{SU(3)}
 \left[(1-\alpha_2)U_\mu(x)
 +\frac{\alpha_2}{4}\sum_{\pm\rho\ne\mu,\nu}
 S_\rho[\bar V](x)\right],\\
 V_\mu(x)
 &=\mathcal P_{SU(3)}
 \left[(1-\alpha_1)U_\mu(x)
 +\frac{\alpha_1}{6}\sum_{\pm\nu\ne\mu}
 S_\nu[\widetilde V](x)\right].
\end{align}
\(\alpha_1,\alpha_2,\alpha_3=(0.75,0.6,0.3)\) 是本实现的默认标准参数。
代码 \fileline{smear/_hyp.py:24--34} 的投影是 SVD 极分解
\(A=U\Sigma V^\dagger\mapsto UV^\dagger\)，再将 determinant 归一为
\(1\)，而不是简单对每个矩阵元取模。这一点决定了群值不变量。

\begin{table}[htbp]
\centering
\caption{HYP 的三级单列伪代码：输入、排除方向、投影与输出}
\small
\begin{tabular}{@{}P{0.94\textwidth}@{}}
\toprule
\(\text{Input: }U_\mu(x)\in SU(3),\;(\alpha_1,\alpha_2,\alpha_3)\)\\
\(\text{for each }\mu,\nu,\rho\text{ with distinct directions: construct }S_\eta[U](x)\)\\
\(\quad \bar V_{\mu;\nu\rho}\gets\mathcal P_{SU(3)}[(1-\alpha_3)U_\mu+
 \alpha_3(\text{two remaining staples})/2]\)\\
\(\quad \widetilde V_{\mu;\nu}\gets\mathcal P_{SU(3)}[(1-\alpha_2)U_\mu+
 \alpha_2(\text{four staples built from }\bar V)/4]\)\\
\(\quad V_\mu\gets\mathcal P_{SU(3)}[(1-\alpha_1)U_\mu+
 \alpha_1(\text{six staples built from }\widetilde V)/6]\)\\
\(\quad \text{if }\alpha_1=0\text{ then return an independent copy of }U\)\\
\(\text{Check: }V_\mu^\dagger V_\mu=I,\;\det V_\mu=1,\;
 \text{axis order remains }(N_t,N_z,N_y,N_x,4,3,3)\)\\
\(\text{Output: }V_\mu(x)\text{ for all four link directions}\)\\
\bottomrule
\end{tabular}
\end{table}
\src{../PyQCD/pyqcd/smear/_hyp.py:61--132；refer/papers/Analytic_Smearing_of_SU3_Link_Variables_latex/}

\paragraph{对称性与极限检查。}
若 \(\alpha_i=0\)，对应层退化为原始链接；若所有 \(U=I\)，所有 staple
和投影均给 \(I\)。由于每一级 staple 具有相同端点变换，线性组合在局部
变换下协变，投影满足 \(\mathcal P(GAG^\dagger)=G\mathcal P(A)G^\dagger\)
（在无奇异极分解的条件下）。这解释了“涂抹后仍可放入规范不变量”的
物理依据；但投影处在近奇异矩阵上时的数值稳定性仍须用真实组态检查。

\section{Stout 涂抹：解析指数保持群值}
\label{sec:stout}
Stout 与 HYP 的根本差别是：Stout 先构造 Lie 代数元素 \(Q_\mu\)，再
指数回到群，而非对线性和做极分解。PyQCD 只涂抹空间链接，时间链接
保持：
\begin{align}
 C_\mu(x)&=\sum_{\substack{\nu\ne\mu\\ \nu=0,1,2}}
 \left[U_\nu(x)U_\mu(x+\hat\nu)U_\nu^\dagger(x+\hat\mu)
 +U_\nu^\dagger(x-\hat\nu)U_\mu(x-\hat\nu)
 U_\nu(x-\hat\nu+\hat\mu)\right],\\
 \Omega_\mu(x)&=\rho C_\mu(x)U_\mu^\dagger(x),\\
 Q_\mu(x)&=\frac{\ii}{2}\bigl(\Omega_\mu^\dagger-\Omega_\mu\bigr)
 -\frac{\ii}{2\Nc}\Tr\bigl(\Omega_\mu^\dagger-\Omega_\mu\bigr)I,\\
 U_\mu^{(n+1)}(x)&=\exp\!\bigl(\ii Q_\mu^{(n)}(x)\bigr)U_\mu^{(n)}(x),
 \qquad \mu\in\{x,y,z\}.
\end{align}
\(Q_\mu\) 是 Hermitian、无迹矩阵，故 \(\exp(\ii Q_\mu)\in SU(3)\)。
实现用 Cayley--Hamilton 形式
\[
 \exp(\ii Q)=f_0 I+f_1Q+f_2Q^2,
\]
其中 \(f_i\) 由 \(c_0=\Tr(Q^3)/3\)、\(c_1=\Tr(Q^2)/2\) 和小角
sinc 展开计算，\(\rho=0.12,n_{\rm step}=20\) 是已有数据命名对应的默认
约定。该解析指数减少了每一步重新做一般矩阵指数的成本，但小 \(Q\)
分支和有限精度仍需检查。

\begin{table}[htbp]
\centering
\caption{Stout 空间链接算法（与 HYP 分开记录）}
\small
\begin{tabular}{@{}P{0.94\textwidth}@{}}
\toprule
\(\text{Input: }U^{(0)}_\mu,\;\rho,\;n_{\rm step};\quad \mu=0,1,2\text{，时间 }\mu=3\text{ 保持}\)\\
\(\text{for }n=0,\ldots,n_{\rm step}-1\text{:}\)\\
\(\quad C_\mu\gets\text{前向与反向空间 staple 的和}\)\\
\(\quad \Omega_\mu\gets\rho C_\mu U_\mu^\dagger\)\\
\(\quad Q_\mu\gets\frac{\ii}{2}(\Omega_\mu^\dagger-\Omega_\mu)
-\frac{\ii}{2\Nc}\Tr(\Omega_\mu^\dagger-\Omega_\mu)I\)\\
\(\quad (c_0,c_1)\gets(\Tr Q_\mu^3/3,\Tr Q_\mu^2/2)\)\\
\(\quad (f_0,f_1,f_2)\gets\text{Cayley--Hamilton 系数（小角 sinc 分支）}\)\\
\(\quad U_\mu^{(n+1)}\gets(f_0I+f_1Q_\mu+f_2Q_\mu^2)U_\mu^{(n)}\)\\
\(\quad U_3^{(n+1)}\gets U_3^{(n)}\)\\
\(\text{Check: }Q_\mu^\dagger=Q_\mu,\;\Tr Q_\mu=0,\;
 U_\mu^{(n+1)\dagger}U_\mu^{(n+1)}=I,\;\det U_\mu^{(n+1)}=1\)\\
\(\text{Output: }U^{(n_{\rm step})}\)\\
\bottomrule
\end{tabular}
\end{table}
\src{../PyQCD/pyqcd/smear/_stout.py:27--165；refer/papers/Analytic_Smearing_of_SU3_Link_Variables_latex/}

\section{Gaussian/Jacobi 夸克涂抹与 momentum smearing}
对夸克场，协变三维 Laplacian \(\nabla^2\) 生成 Gaussian smearing：
\begin{align}
 J_{\sigma,n_\sigma}
 &=\left(1+\frac{\sigma\nabla^2}{n_\sigma}\right)^{n_\sigma}
 \xrightarrow[n_\sigma\to\infty]{}\exp(\sigma\nabla^2),\\
 \left(\nabla^2\psi\right)(x)
 &=\sum_{j=1}^3\left[
 U_j(x)\psi(x+\hat j)+U_j^\dagger(x-\hat j)\psi(x-\hat j)
 -2\psi(x)\right].
\end{align}
在自由场平面波 \(\psi_p(x)=e^{ipx}\) 上，\(\nabla^2\) 的本征值为
\[
 \lambda(p)=2\sum_{j=1}^3\bigl[\cos(ap_j)-1\bigr]
 =-4\sum_{j=1}^3\sin^2(ap_j/2),
\]
所以 Gaussian 因子 \(\exp(\sigma\lambda(p))\) 抑制大 \(p\)，而
momentum smearing 以空间相位把峰值从 \(p=0\) 移到目标动量：
\[
 \psi^{(p)}(x)=e^{-\ii\boldsymbol k\cdot\boldsymbol x}\psi(x),\qquad
 \Box_p=e^{-\ii\boldsymbol k\cdot x}\Box e^{+\ii\boldsymbol k\cdot x}.
\]
PyQCD 的 \fileline{vertex/_vertex.py:643--733} 实现的是对 eigvec 逐点乘
\[
 \exp\!\left[-2\pi\ii\left(
 \frac{p_z z}{L_z}+\frac{p_y y}{L_y}+\frac{p_x x}{L_x}\right)\right],
\]
并明确区分它与 sink momentum projection、VdV/VVV contraction。

\begin{table}[htbp]
\centering
\caption{Gaussian/Jacobi 与 momentum smearing 的单列算法}
\small
\begin{tabular}{@{}P{0.94\textwidth}@{}}
\toprule
\(\text{Input: }\psi,\;U_j,\;\sigma,\;n_\sigma,\;\boldsymbol p=(p_z,p_y,p_x)\)\\
\(\text{Build: }(\nabla^2\psi)(x)\text{，每一项均用前向/反向链接平行传输}\)\\
\(\text{Iterate: }\psi_{k+1}\gets\psi_k+\sigma\nabla^2\psi_k/n_\sigma\)\\
\(\quad\text{after }n_\sigma\text{ steps, }\psi_{n_\sigma}\to J_{\sigma,n_\sigma}\psi\)\\
\(\text{Momentum option: }\psi_{n_\sigma}(x)\gets e^{-\ii\boldsymbol p\cdot\boldsymbol x}
\psi_{n_\sigma}(x)\)\\
\(\text{Check: }U=I\text{ 时恢复自由平面波因子；}\boldsymbol p=0\text{ 时相位为 1}\)\\
\(\text{Output: }\psi^{\rm smear}_{\boldsymbol p}\text{，其形状/后端/复数精度保持输入契约}\)\\
\bottomrule
\end{tabular}
\end{table}
\src{../PyQCD/pyqcd/vertex/_vertex.py:643--749；MyQCD \fileline{derivations.py} 中 Gaussian/Jacobi 与 momentum-shift 推导}

\section{蒸馏：空间低模投影与 perambulator}
\label{sec:distillation}
蒸馏的物理直觉是用三维协变 Laplacian 的低频子空间构造平滑、规范协变
的强子算符。固定时间 \(t\)，解
\[
 -\widetilde\Delta(t)v^{(k)}(t)=\lambda_k(t)v^{(k)}(t),\qquad
 V(t)=\bigl(v^{(1)}(t),\ldots,v^{(N_{\rm ev})}(t)\bigr),
\]
并定义低秩投影
\[
 \Box(t)=V(t)V^\dagger(t),\qquad
 \Box^2=\Box,\qquad \Box^\dagger=\Box.
\]
若 \(\{v^{(k)}\}\) 是正交本征矢，\(\Box\) 是投影而非任意平滑核。夸克传播
\[
 \tau(t,t')=V^\dagger(t)\,D^{-1}(t,t')\,V(t')
\]
称为 perambulator；空间颜色自由度被压缩到 \(N_{\rm ev}\)，Dirac 指标
仍保留。有限 \(N_{\rm ev}\) 的截断误差不是统计误差，应在 \(N_{\rm ev}\)
扫描中单独报告。

\begin{table}[htbp]
\centering
\caption{蒸馏与 perambulator 构造算法}
\small
\begin{tabular}{@{}P{0.94\textwidth}@{}}
\toprule
\(\text{Input: }U_j(t),\;D[U],\;N_{\rm ev}\text{，每个时间片}\)\\
\(\text{Construct: }-\widetilde\Delta(t)\text{，在空间颜色空间保持 Hermitian}\)\\
\(\text{Solve: }-\widetilde\Delta v^{(k)}=\lambda_kv^{(k)},\;
 \lambda_1\le\cdots\le\lambda_{N_{\rm ev}}\)\\
\(\text{Orthonormalize/check: }V^\dagger V=I,\;\Box=VV^\dagger,\;\Box^2=\Box\)\\
\(\text{Project source/sink: }\eta(t)=V^\dagger(t)\,\psi(t)\)\\
\(\text{Solve Dirac system: }D\,S=\text{distilled sources}\)\\
\(\text{Compress: }\tau(t,t')=V^\dagger(t)S(t,t')\)\\
\(\text{Check: }\gamma_5D\gamma_5=D^\dagger\text{（若采用该费米子约定）}\)\\
\(\text{Output: }V(t),\;\Box(t),\;\tau(t,t')\text{，并记录 }N_{\rm ev},D\text{ 方案与边界}\)\\
\bottomrule
\end{tabular}
\end{table}
\src{../PyQCD/docs/格点QCD蒸馏方法解析.tex；../PyQCD/pyqcd/contraction/_seqperam.py:20--59}

\section{VdV 与 VVV：算符元素的物理含义}
\label{sec:vertices}
对动量 \(p\) 的介子/双线性元素，代码实现的是
\[
 V_{mn}(p,t)=\sum_{\boldsymbol x}
 e^{-\ii\boldsymbol p\cdot\boldsymbol x}
 \phi_m^\dagger(\boldsymbol x,t)\phi_n(\boldsymbol x,t).
\]
因此输入是每个时间片的 \(\phi\) 和形状为
\((N_p,N_zN_yN_xN_c)\) 的相位，输出是 \((N_p,N_{\rm ev},N_{\rm ev})\)。
重子颜色 singlet 元素为
\[
 V_{mnl}(p,t)=\sum_{\boldsymbol x}
 e^{-\ii\boldsymbol p\cdot\boldsymbol x}
 \epsilon_{abc}\phi_m^a(\boldsymbol x,t)
 \phi_n^b(\boldsymbol x,t)\phi_l^c(\boldsymbol x,t).
\]
\(\epsilon_{abc}\) 保证交换任意两个蒸馏指标时变号；如果两个向量相同，
元素为零。VdV 则满足
\[
 V_{mn}(p,t)^*=V_{nm}(-p,t)
\]
（相位、坐标和本征矢按 Hermitian 共轭定义）。新增
\texttt{vertices\_vdv\_vvv} 审计对这些有限维指标关系逐项展开并通过。
\src{../PyQCD/pyqcd/vertex/_vertex.py:32--234；MyQCD \fileline{pyqcd_theory_audit.py:190--258}}

\section{本章的证据边界}
HYP/Stout 的群投影、Gaussian 的自由平面波因子、\(\Box^2=\Box\) 和
VdV/VVV 的指标对称性是可验证的结构。它们不能单独证明：
\begin{enumerate}
 \item 某一真实 lattice ensemble 的 eigenvector 文件与 \(U\) 来源一致；
 \item \(N_{\rm ev}=100\) 已足以覆盖目标强子或高动量态；
 \item smearing 没有改变目标矩阵元的方案定义；
 \item GPU/CuPy/Torch 后端上的 contraction 没有轴序、精度或设备隐式拷贝错误。
\end{enumerate}
这些问题需要 \fileline{data/} 中的实际文件、缓存契约、设备信息和独立
CPU/GPU 对照。当前报告只把它们列为验证门，不把结构审计升级为生产结论。


\iffalse
\chapter{Grassmann 积分、Wick 收缩与关联函数}
\label{sec:wick}

\section{从路径积分到传播子}
费米子部分在给定规范背景下是 Grassmann 高斯积分。将全部颜色、Dirac、
味和格点坐标压缩为复合指标 (i,j)，写成
\begin{align}
 Z_F[\eta,\bar\eta;U]
 &=\int\mathcal D\bar\psi\mathcal D\psi\,
 \exp\left[-\bar\psi D[U]\psi+\bar\eta\psi+\bar\psi\eta\right]
 \nonumber\\
 &=\det D[U],\exp\left(\bar\eta D^{-1}[U]\eta\right),
 \\[-2pt]
 S[U]&=D^{-1}[U].
\end{align}
这里要求 (D) 在所选边界、质量和规范背景下可逆；零模或接近零模时，
数值求逆和统计估计都要另作处理。对 Grassmann 源求导得到 Wick 定理：
\begin{align}
 \langle\psi_{i_1}\cdots\psi_{i_n}
 \bar\psi_{j_1}\cdots\bar\psi_{j_n}\rangle_U
 &=\sum_{\pi\in S_n}(-1)^{\pi}
 \prod_{k=1}^{n}S_{i_kj_{\pi(k)}}[U],
 \\[-2pt]
 \langle\psi_i\bar\psi_j\rangle_U&=S_{ij}[U].
\end{align}
费米负号不是经验系数，而是把 Grassmann 场重新排序到成对收缩所需的
置换奇偶性。这是自动 Wick 引擎可以枚举而不能随意省略的第一性原理。
它依赖 Grassmann 反对易公理和高斯测度；相互作用通过 (det D) 和规范
组态平均进入，而不是被上述有限维代数单独消除。

若费米子离散化满足
\begin{equation}
 D^\dagger=\gamma_5D\gamma_5,
 \qquad
 S^\dagger=\gamma_5S\gamma_5,
\end{equation}
则反向时间/源汇顺序的传播子可以从已有传播子重建。PyQCD 的
\texttt{seq\_peram} 实际执行
\begin{equation}
 \tau_{\rm seq}(t_{\rm sink},t_{\rm src})
 =\gamma_5\,\tau(t_{\rm src},t_{\rm sink})^\dagger\,\gamma_5,
 \label{eq:seq-peram}
\end{equation}
其中厄米共轭同时交换 Dirac 和蒸馏指标。这个等式的前提是所用 (D)、
边界和 gamma 约定确实满足 (gamma_5)-Hermiticity；代码接口本身不能
替调用者证明该前提。
\src{../PyQCD/pyqcd/contraction/_seqperam.py:20--59；refer/books/INTRODUCTION_TO_LATTICE_QCD_latex/chapters/sec17_correlators.tex}

\section{算符、味平衡和自动收缩}
一个双线性算符的典型形式为
\begin{equation}
 \mathcal O_M(x)=\bar q_a^f(x)\,\Gamma_{\alpha\beta}\,
 q_b^g(x)\,\mathcal C_{ab}^{fg},
\end{equation}
重子算符则以 (epsilon_{abc}) 闭合颜色：
\begin{equation}
 \mathcal O_B(x)=\epsilon_{abc}
 (q_1^a q_2^b)q_3^c,
\end{equation}
并由 gamma 链、味结构和自旋投影决定具体态。PyQCD 的 DSL 用 ``|``
分隔不同时间/粒子区间，用 ``gamma_5``、``gamma_7``、``gamma_mu`` 等
字符串表示 Dirac 插入，用 ``u``、``d`` 和 ``u^d`` 等标记夸克及反夸克。
分隔符不是物理场，而是解析器识别 VDV/VVV 区间和时间标签的元数据。

\begin{table}[htbp]
\centering
\caption{自动 Wick 收缩的单列算法：从算符 DSL 到可执行 contraction plan}
\small
\begin{tabular}{@{}P{0.94\textwidth}@{}}
\toprule
\(\text{Input: }\texttt{sink\_operators},\;\texttt{source\_operators},\;
\texttt{curr\_operators},\;Cpt\in\{\text{bubble},2\text{pt},3\text{pt},4\text{pt}\}\)\\
\(\text{Normalize: }\text{若首尾没有 }|\text{，自动补上；2pt 情形丢弃 current 区间}\)\\
\(\text{Parse: }\text{把味标记分为 }q\text{ 与 }q^d，\text{把 gamma 和 }|\text{ 位置分别记录}\)\\
\(\text{Expand wildcard: }q\to\{u,d,s,c,b,t\},\quad l\to\{u,d\},\quad\text{逐组合代入}\)\\
\(\text{Balance test: }N_q(f)=N_{\bar q}(f)\text{；不平衡的味组合直接跳过}\)\\
\(\text{Label fields: }(q,\bar q)\text{ 按出现次序分配 }ab,cd,ef,\ldots\text{ 收缩标签}\)\\
\(\text{Enumerate: }\text{对每种味分别生成夸克到反夸克的所有排列，再做不同味匹配的笛卡尔积}\)\\
\(\text{Fermi sign: }\mathcal s=(-1)^{N_{\rm inv}},\quad N_{\rm inv}=\#\{(i,j):i<j,\,\ell_i>\ell_j\}\)\\
\(\text{Time parse: }|\cdots|\text{ 内含两个夸克给 VDV，含三个夸克给 VVV；区间标成 }t_{\rm sink},t_{\rm cur},t_{\rm src}\)\\
\(\text{Gamma parse: }\text{每个 gamma 读取相邻夸克标签，生成 gamma registry 的输入索引}\)\\
\(\text{Peram entry: }S_{q\bar q}[t_q,t_{\bar q}]\text{ 保存 }(q\text{ 位},\bar q\text{ 位},\text{味},\text{Dirac/ev 标签})\)\\
\(\text{Vertex entry: }\text{两夸克区间绑定 }V_{mn}\text{，三夸克区间绑定 }V_{mnl}\text{，并写入时间标签}\)\\
\(\text{Index projection: }\text{重复小写字母求和，自由小写字母保留；大写动量/链接标签保留一次}\)\\
\(\text{Diagram output: }\texttt{result\_indx},\texttt{result\_name},\texttt{result\_sign},\texttt{peram},\texttt{V},\texttt{gamma\_pos}\)\\
\(\text{For each diagram: }\mathcal C_D=\mathcal s_D\,\operatorname{einsum}(V_{\rm sink},S_1,\ldots,S_n,\Gamma,V_{\rm src})\)\\
\(\text{2pt: }\mathcal C_2(t)=\sum_D\mathcal C_D(t)\)\\
\(\text{3pt: }\mathcal C_3(t,\tau)=\sum_D\mathcal C_D(t,\tau;\Gamma_{\rm cur})\)\\
\(\text{4pt/bubble: }\text{保留所有 current/区间连通分支，并按 registry 解析每个中间量}\)\\
\(\text{Check: }\text{味平衡、颜色闭合、输出自由指标与预期轴序一致；否则停止而不猜测}\)\\
\(\text{Output: }\text{一个字典（无 wildcard）或字典列表（有 wildcard），供静态绘图和动态 contraction 使用}\)\\
\bottomrule
\end{tabular}
\end{table}
\src{../PyQCD/pyqcd/contraction/_autowick.py:39--59,208--445；../PyQCD/pyqcd/contraction/_dynamic.py:262--647；../PyQCD/docs/格点QCD中的维克收缩与傅里叶变换.tex}

该算法的“枚举所有排列”在 (n_f) 个相同味夸克/反夸克时规模含有
(n_f!) 因子；这是物理拓扑的完整性与计算复杂度之间的直接冲突。实际
高阶多强子问题应先利用同味交换对称性、颜色张量和等价图合并，再由
\texttt{dynamic\_contraction} 根据 einsum 路径和设备估算执行顺序。合并
必须保持相对费米号，不能只按字符串去重。
\fi
\chapter{证据分级、参考源与后续核验}
\label{sec:boundaries}

\section{理论审计结果}
本次审计结果落盘到 \fileline{data/pyqcd_theory_audit/theory_audit.json}，
由 \fileline{myqcd/pyqcd_theory_audit.py} 生成。联合检查包含 9 项独立
结构审计和既有 MyQCD 核心 SymPy 审计；其中 7 项为 verified、1 项为
structural、1 项为 unverified，\texttt{failed\_local\_checks=0}，
\texttt{missing\_files=0}，核心审计状态为 verified。

\begin{table}[htbp]
\centering
\caption{本次理论审计的统计摘要}
\small
\begin{tabularx}{0.96\textwidth}{@{}P{4.2cm}Y@{}}
\toprule
项目 & 结果\\
\midrule
结构审计总数 & 9\\
verified / structural / unverified & 7 / 1 / 1\\
failed\_local\_checks & 0\\
\fileline{source_inventory.missing_files} & 0\\
\fileline{existing_myqcd_core.status} & verified\\
\bottomrule
\end{tabularx}
\end{table}

未验证边界保持如下：
\begin{itemize}
  \item \texttt{matching\_tmd\_boundary} 仍是 \unverified，因为完整 TMD 需要
  soft factor、rapidity renormalization、Collins--Soper 核和真实卷积闭环。
  \item finite-\(a\) 的 straight OPE 与 finite-staple TMD bilocal 不是同一对象。
  \item \texttt{quasi\_pdf\_gluon} 的 sin 通道与 \texttt{quasi\_tmd\_pdf} 的
  Hermitian 半轴 cos/sin 通道不能互换。
  \item \texttt{hR\_PDF} 对精确 \(x=0\) 节点直接拒绝；这里不能用任意小数替代
  principal-value 处方。
  \item Wilson flow 只做 UV 平滑，不替代 soft subtraction 或 rapidity renorm。
  \item 有限维 SymPy 只能证明局部代数恒等式，不能把真实系综结果升级成定理。
\end{itemize}
\src{data/pyqcd_theory_audit/theory_audit.json；../myqcd/pyqcd_theory_audit.py}

\section{参考源汇总}
\begin{table}[htbp]
\centering
\caption{本报告优先使用的参考源与它们对应的物理语义}
\small
\begin{tabularx}{0.96\textwidth}{@{}P{3.1cm}P{6.4cm}Y@{}}
\toprule
类别 & 代表路径 & 主要承担的物理角色\\
\midrule
教材/公理 &
\fileline{refer/books/格点QCD导论_latex/chapters/sec06_formulation_gauge.tex}；
\fileline{refer/books/格点QCD导论_latex/chapters/sec17_correlators.tex}；
\fileline{refer/books/格点QCD导论_latex/chapters/sec18_hadron_spectrum.tex}；
\fileline{refer/books/格点QCD导论_latex/chapters/sec16_errors.tex}
& 规范场、相关函数、谱学和统计误差的基础定义。\\
OPE/PDF/LaMET &
\fileline{refer/papers/准部分子分布单圈匹配_latex/}；
\fileline{refer/papers/Quasi_Pseudo_PDF_Radyushkin_latex/}；
\fileline{refer/papers/Complete_NP_Renorm_QuasiPDF_latex/}；
\fileline{refer/papers/Accessing_Gluon_PDF_LaMET_latex/}；
\fileline{refer/papers/PDF_Moments_Any_Order_latex/}
& 匹配核、准 PDF、非微扰重整化与矩展开。\\
TMD/soft/CS &
\fileline{refer/papers/从准TMD波函数确定CS核_latex/}；
\fileline{refer/papers/TMD_Soft_Function_LaMET_latex/}；
\fileline{refer/papers/LaMET计算TMD软函数_latex/}；
\fileline{refer/papers/Locally_Smeared_OPE_latex/}
& staple 几何、软因子、rapidity 方案与 CS 核。\\
流与 SFTX &
\fileline{refer/papers/Properties_Uses_Wilson_Flow_latex/}；
\fileline{refer/papers/Quasi_PDF_Gradient_Flow_latex/}；
\fileline{refer/papers/准部分子分布与梯度流_latex/}
& Wilson flow、小流时间展开和 flow-based renormalization。\\
实现对照 &
\fileline{../PyQCD/docs/格点QCD中的关联函数.tex}；
\fileline{../PyQCD/docs/格点QCD中的OPE算符.tex}；
\fileline{../PyQCD/docs/格点QCD中的TMD_PDF.tex}；
\fileline{../PyQCD/pyqcd/...}
& 接口、轴序、实现契约和源码级对照。\\
\bottomrule
\end{tabularx}
\end{table}

\chapter{欧氏 QCD 的构造与连续极限：从公理到离散作用量}
\label{sec:deep-lattice}

前面的接口审计已经固定了数组的轴序，但数组只有在一个明确的场论构造中
才有物理意义。本章把这一构造完整写出。连续 Euclidean QCD、格点链接、
plaquette、费米子核、泛函测度和连续极限不是彼此独立的背景知识：它们共同
决定了后来可以对哪个对象做涂抹、求逆、收缩和外推。局部规范协变性是格点
上可以逐项保持的精确性质；洛伦兹对称、手征对称和连续谱则一般只在
\(a\to0\) 或经过改进后恢复。

\section{局域规范原理如何强制产生链接}

取 \(T^a\) 为 Hermitian 的 \(SU(N_c)\) 基本表示生成元，满足
\begin{equation}
 \operatorname{tr}(T^aT^b)=\frac12\delta^{ab},
 \qquad [T^a,T^b]=\mathrm{i}f^{abc}T^c,
 \qquad A_\mu=A_\mu^aT^a .
 \label{eq:deep-suN-convention}
\end{equation}
采用 \(D_\mu=\partial_\mu+\mathrm{i}gA_\mu\) 的约定。若
\(\psi^G(x)=G(x)\psi(x)\)，要求协变导数满足
\(D_\mu^G\psi^G=G(D_\mu\psi)\)，逐项比较导数项得到
\begin{align}
 A_\mu^G(x)
 &=G(x)A_\mu(x)G^\dagger(x)
   +\frac{\mathrm{i}}{g}(\partial_\mu G(x))G^\dagger(x),
 \label{eq:deep-gauge-transform}\\
 F_{\mu\nu}
 &=\frac{1}{\mathrm{i}g}[D_\mu,D_\nu]
 =\partial_\mu A_\nu-\partial_\nu A_\mu
  +\mathrm{i}g[A_\mu,A_\nu],
 \qquad F_{\mu\nu}^G=GF_{\mu\nu}G^\dagger .
 \label{eq:deep-field-strength}
\end{align}
这里第二式由协变导数的交换子定义出来。因而
\(\operatorname{tr}F_{\mu\nu}F_{\mu\nu}\) 是局域规范标量，连续作用量
\begin{equation}
 S_E=\int\dd^4x\left[
 \frac12\operatorname{tr}(F_{\mu\nu}F_{\mu\nu})
 +\sum_{f=1}^{N_f}\bar\psi_f(\gamma_\mu D_\mu+m_f)\psi_f\right]
 \label{eq:deep-continuum-action}
\end{equation}
在 Euclidean gamma 矩阵满足
\(\{\gamma_\mu,\gamma_\nu\}=2\delta_{\mu\nu}\) 时，规范场部分为实数；
Grassmann 双线性本身不是普通复数，费米子积分后的 determinant 才进入
玻色化权重。由于
\(F^a_{\mu\nu}F^a_{\mu\nu}/4=\frac12\operatorname{tr}F_{\mu\nu}^2\)，式
\eqref{eq:deep-continuum-action} 与通常的分量形式完全相同。

从 \(x\) 到 \(x+a\hat\mu\) 的有限平行输运要求输运后的场仍按端点变换，
唯一自然的对象是群值路径有序指数
\begin{equation}
 U_\mu(x)=\mathcal P_{x\leftarrow x+a\hat\mu}\exp\left[\mathrm{i}g
 \int_0^a\dd s\,A_\mu(x+s\hat\mu)\right]
 \in SU(N_c),
 \qquad
 U_\mu^G(x)=G(x)U_\mu(x)G^\dagger(x+a\hat\mu).
 \label{eq:deep-link}
\end{equation}
把相邻链接相乘时，中间端点的 \(G^\dagger G\) 自动抵消。因此一条从
\(x_0\) 到 \(x_n\) 的开链
\begin{equation}
 W(x_0,x_n)=U_{\mu_1}(x_0)U_{\mu_2}(x_1)\cdots
 U_{\mu_n}(x_{n-1})
 \label{eq:deep-open-line}
\end{equation}
满足
\begin{equation}
 W^G(x_0,x_n)=G(x_0)W(x_0,x_n)G^\dagger(x_n).
 \label{eq:deep-open-line-transform}
\end{equation}
所以只有闭合路径的迹，或由端点费米场补帽后的颜色 singlet，才能成为
规范不变量。这个端点定理正是 PyQCD 中 straight OPE、staple TMD、
\texttt{VdV\_sink\_t\_link} 和 \texttt{gluon\_tmd\_operator}
共用的物理底层；它不依赖某个具体数组库。

\section{plaquette 展开与 Wilson 规范作用量}

最小闭合回路是
\begin{equation}
 P_{\mu\nu}(x)=U_\mu(x)U_\nu(x+a\hat\mu)
 U_\mu^\dagger(x+a\hat\nu)U_\nu^\dagger(x),
 \qquad P_{\mu\nu}^G=G(x)P_{\mu\nu}(x)G^\dagger(x).
 \label{eq:deep-plaquette}
\end{equation}
将每条链接放在几何中点并使用 Baker--Campbell--Hausdorff 展开，可以得到
\begin{equation}
 P_{\mu\nu}(x)
 =\exp\left[\mathrm{i}g a^2F_{\mu\nu}(x_c)
 +O(a^3D F)+O(a^4F^2)\right],
 \qquad x_c=x+\tfrac a2(\hat\mu+\hat\nu),
 \label{eq:deep-plaquette-bch}
\end{equation}
为了看清作用量的系数，只需在迹中展开到 \(a^4\)：
\begin{align}
 \Re\operatorname{Tr}P_{\mu\nu}
 &=N_c-\frac{g^2a^4}{2}\operatorname{tr}(F_{\mu\nu}^2)+O(a^6)
 \label{eq:deep-plaquette-trace}\\
 &=N_c-\frac{g^2a^4}{4}F_{\mu\nu}^aF_{\mu\nu}^a+O(a^6),
 \nonumber\\
 1-\frac1{N_c}\Re\operatorname{Tr}P_{\mu\nu}
 &=\frac{g^2a^4}{4N_c}F_{\mu\nu}^aF_{\mu\nu}^a+O(a^6).
 \label{eq:deep-plaquette-real}
\end{align}
因而令 \(\beta=2N_c/g_0^2\)，Wilson 规范作用量
\begin{equation}
 S_g[U]=\beta\sum_x\sum_{\mu<\nu}
 \left(1-\frac1{N_c}\Re\operatorname{Tr}P_{\mu\nu}(x)\right)
 \label{eq:deep-wilson-gauge-action}
\end{equation}
在连续极限给出
\begin{align}
 S_g
 &=\sum_x\sum_{\mu<\nu}\frac{a^4}{2}F_{\mu\nu}^aF_{\mu\nu}^a+O(a^2)\\
 &=\int\dd^4x\,\frac14F_{\mu\nu}^aF_{\mu\nu}^a+O(a^2).
 \label{eq:deep-wilson-continuum}
\end{align}
第二行使用了 \(\sum_{\mu<\nu}F_{\mu\nu}^2=\frac12\sum_{\mu,\nu}F_{\mu\nu}^2\)。
这一步同时说明：\(U\) 必须是群元素而非任意复矩阵，取实迹保证作用量
实数，而有限 \(a\) 的误差来自更高维、保持格点对称的局域回路。

Euclidean 配分函数的定义是
\begin{equation}
 Z=\int\prod_{x,\mu}dU_\mu(x)
 \prod_{f=1}^{N_f}d\bar\psi_f d\psi_f\,
 \exp[-S_g[U]-S_f[U,\bar\psi,\psi]],
 \label{eq:deep-partition-function}
\end{equation}
其中 \(dU\) 是每条链接上的 Haar 测度。左不变性与右不变性使得
\(dU=d(GUG^\dagger)\)，因此变量变换不会改变 \(Z\)。对费米场做 Grassmann
积分后，若
\begin{equation}
 S_f=\sum_f\bar\psi_fD_f[U]\psi_f,
 \qquad
 \int d\bar\psi d\psi\,e^{-\bar\psi D\psi}=\det D,
 \label{eq:deep-fermion-determinant}
\end{equation}
则
\begin{equation}
 Z=\int[dU]\,e^{-S_g[U]}
 \prod_{f=1}^{N_f}\det D_f[U].
 \label{eq:deep-gauge-only-partition}
\end{equation}
这解释了为什么在真实 dynamical simulation 中费米子不能只被看作“传播子
求逆”：\(\det D\) 改变规范组态的统计权重；而在给定组态上计算
\(D^{-1}\) 只负责插入算符的收缩。

\section{费米子离散、doubling 与 Wilson 项}

定义规范协变的一阶差分
\begin{align}
 \nabla_\mu\psi(x)&=\frac{U_\mu(x)\psi(x+a\hat\mu)-\psi(x)}{a},\\
 \nabla_\mu^*\psi(x)&=\frac{\psi(x)-U_\mu^\dagger(x-a\hat\mu)\psi(x-a\hat\mu)}{a},
\end{align}
并令 \(\Delta_\mu=\nabla_\mu^*\nabla_\mu\)。朴素费米子作用量与 Wilson
作用量分别是
\begin{align}
 S_N&=a^4\sum_x\bar\psi(x)
 \left[m_0+\frac12\sum_\mu\gamma_\mu(\nabla_\mu+\nabla_\mu^*)\right]\psi(x),\\
 S_W&=a^4\sum_x\bar\psi(x)
 \left[m_0+\frac12\sum_\mu\gamma_\mu(\nabla_\mu+\nabla_\mu^*)
 -\frac{ar}{2}\sum_\mu\Delta_\mu\right]\psi(x).
 \label{eq:deep-wilson-fermion-action}
\end{align}
这里 \(r\) 通常取 1。所有差分项都由前后链接连接，因此式
\eqref{eq:deep-wilson-fermion-action} 逐项规范不变。

在自由场 \(U_\mu=1\) 且采用平面波
\(\psi(x)=u(p)e^{ip\cdot x}\) 时，朴素核的动量形式为
\begin{equation}
 D_N(p)=m_0+\frac{\mathrm{i}}a\sum_\mu\gamma_\mu\sin(ap_\mu),
 \qquad
 D_N^\dagger(p)D_N(p)=m_0^2+\frac1{a^2}\sum_\mu\sin^2(ap_\mu).
 \label{eq:deep-naive-free-kernel}
\end{equation}
除了 \(p_\mu=0\) 外，任何分量 \(p_\mu=\pi/a\) 也使正弦为零，故在布里渊
区角点共有 \(2^4=16\) 个轻模。这个结论是动量空间的零点计数，不是
数值离散的偶然误差。

Wilson 项把核改为
\begin{equation}
 D_W(p)=m_0+\frac{\mathrm{i}}a\sum_\mu\gamma_\mu\sin(ap_\mu)
 +\frac{r}{a}\sum_\mu[1-\cos(ap_\mu)].
 \label{eq:deep-wilson-free-kernel}
\end{equation}
若有 \(N_\pi\) 个分量取 \(p_\mu=\pi/a\)，则该模额外得到
\begin{equation}
 m_{\rm eff}=m_0+\frac{2rN_\pi}{a}.
 \label{eq:deep-doubler-mass}
\end{equation}
因此除物理角点外的 doublers 在 \(a\to0\) 时质量发散而解耦。代价是
Wilson 项显式破坏有限格距上的手征对称，导致临界质量发生加性重整化；
这与后续 unpolarized quasi-PDF 的 vector/scalar 混合有同一根源：有限格距
对称性比连续 \(O(4)\times SU(N_f)_L\times SU(N_f)_R\) 更小。

在 \(\gamma_5\)-Hermiticity 成立时，
\begin{equation}
 D_W^\dagger=\gamma_5D_W\gamma_5,
 \qquad Q=\gamma_5D_W=Q^\dagger .
 \label{eq:deep-gamma5-hermiticity}
\end{equation}
两种简并味的 determinant 可写成
\begin{equation}
 \det(D_W)^2=\det(D_W^\dagger D_W)=\det(Q^2)\ge0,
 \label{eq:deep-det-positivity}
\end{equation}
从而可以引入复伪费米子
\begin{equation}
 \det(Q^2)\propto\int d\phi^\dagger d\phi
 \exp[-\phi^\dagger(Q^2)^{-1}\phi].
 \label{eq:deep-pseudofermion}
\end{equation}
这一步是 HMC 轨道可以处理 dynamical fermion 的概率基础。单个非简并
味、有限化学势或不满足正定条件的核则不能未经说明地套用这一式。

\section{Symanzik 展开与改进的物理含义}

格点理论在尺度 \(\pi/a\) 截断连续场论。只要保持局域性、规范对称性和
格点离散对称，连续理论可用 Symanzik 有效作用量表示为
\begin{equation}
 \mathcal L_{\rm lat}
 =\mathcal L_{\rm QCD}
 +a\sum_i c_i^{(5)}\mathcal O_i^{(5)}
 +a^2\sum_i c_i^{(6)}\mathcal O_i^{(6)}+\cdots .
 \label{eq:deep-symanzik-action}
\end{equation}
对一个格点算符也有
\begin{equation}
 \mathcal O_{\rm lat}
 =\sum_{j,\,d_j\le d_\mathcal O}
 Z_{\mathcal O j}(a\mu)\,a^{d_j-d_\mathcal O}\mathcal O_j^R(\mu)
 +a\sum_k d_k(a\mu)\mathcal O_k^{(d_\mathcal O+1)}(\mu)+\cdots .
 \label{eq:deep-symanzik-operator}
\end{equation}
如果同一离散表示中存在低维 \(\mathcal O_j\)，其系数可能按
\(a^{d_j-d_\mathcal O}\) 发散；这就是非局域 Wilson 线线性发散和有限格距
费米子 mixing 必须单独账本化的原因。改进作用量只能取消已识别的高维项，
不能凭“采用了 smearing”自动取消低维混合。

以 Wilson 费米子为例，Sheikholeslami--Wohlert 项在连续展开中含有
\begin{equation}
 \delta S_{\rm SW}
 =-a^5\frac{r g_0c_{\rm SW}}4\sum_x\bar\psi(x)
 \sigma_{\mu\nu}F_{\mu\nu}(x)\psi(x),
 \qquad
 \sigma_{\mu\nu}=\frac{\mathrm{i}}2[\gamma_\mu,\gamma_\nu].
 \label{eq:deep-clover-improvement}
\end{equation}
在本文的 \(D_\mu=\partial_\mu+\mathrm{i}g_0A_\mu\) 和上述
\(\sigma_{\mu\nu}\) 约定下，直接计算得到
\begin{equation}
 (\gamma_\mu D_\mu)^2=D^2+\frac{g_0}{2}\sigma_{\mu\nu}F_{\mu\nu}.
 \label{eq:deep-dirac-square}
\end{equation}
因此 Wilson 二阶项在壳上诱导同一维的 Pauli 结构；选择树级
\(c_{\rm SW}=1\)（量子层面按方案计算）可以在壳上消去一阶 \(O(a)\) 误差，但仍要同时改进
观测算符和重整化常数。对相关函数 \(C_a\) 的连续外推应写成
\begin{equation}
 C_a=C_0+c_1(a\Lambda)^p+c_2(a\Lambda)^{p+1}+\cdots,
 \label{eq:deep-continuum-fit}
\end{equation}
其中 \(p\) 由作用量、算符和所保留对称共同决定。若只用一个 \(a\)，
式\eqref{eq:deep-continuum-fit} 的系数没有可辨识性；若不同 \(a\) 使用了
不同的 renormalization scheme，也不能把数据直接放进同一个拟合。

\section{从作用量到 HMC：辅助算法的完整链}

下面的算法以无零模、成对质量简并的费米子为对象，表中 \(f\) 标记一个
双味行列式块 \(\det(D_f^\dagger D_f)\)，而非每个单味。
它说明如何从已定义的 Euclidean 权重产生组态，以及如何在每条
轨道中求解费米子线性系统。它不是对 QCD 非微扰积分的解析解；其物理
正确性来自 Haar 测度、Hamiltonian 体积保持、可逆性和 Metropolis 接受步骤。
\begin{longtable}{@{}P{0.94\textwidth}@{}}
\caption{由 Wilson 作用量产生 dynamical 组态的 HMC 单列算法}\label{tab:deep-hmc}\\
\toprule
\textbf{输入、扩展相空间、可逆演化、接受率与验证}\\
\midrule
\endfirsthead
\caption[]{由 Wilson 作用量产生 dynamical 组态的 HMC 单列算法（续）}\\
\toprule
\textbf{输入、扩展相空间、可逆演化、接受率与验证}\\
\midrule
\endhead
\bottomrule
\endfoot
\(\text{Input: }U\in SU(N_c)^{4V},\;S_g[U],\;D_f[U],\;\Delta\tau,\;N_{\rm step},\;N_{\rm traj},\;\text{边界与质量参数}\)\\
\(\text{Momentum: }\Pi_\mu(x)=\Pi_\mu^a(x)T^a\text{，}\;S_\Pi=\sum_{x,\mu}\operatorname{tr}\Pi_\mu^2=\frac12\sum_{x,\mu,a}(\Pi_\mu^a)^2\)\\
\(\text{Pseudofermion: }\phi_f\gets Q_f\,\eta_f,\;\eta_f\sim\exp(-\eta_f^\dagger\eta_f),\;Q_f=\gamma_5D_f\text{（正定简并味情形）}\)\\
\(\text{Extended Hamiltonian: }H[U,\Pi,\phi]=S_g[U]+S_\Pi[\Pi]+\sum_f\phi_f^\dagger(Q_f^2)^{-1}\phi_f\)\\
\(\text{Force: }\dot U_\mu=\mathrm{i}\Pi_\mu U_\mu,\quad\dot\Pi_\mu^a=-\frac{\partial S_g}{\partial\epsilon_\mu^a}-\frac{\partial S_{\rm pf}}{\partial\epsilon_\mu^a},\;U\mapsto e^{\mathrm{i}\epsilon T^a}U\)\\
\(\text{Leapfrog: }\Pi_{1/2}=\Pi_0-\frac{\Delta\tau}{2}\nabla S(U_0),\;U_1=e^{\mathrm{i}\Delta\tau\Pi_{1/2}}U_0,\;\Pi_1=\Pi_{1/2}-\frac{\Delta\tau}{2}\nabla S(U_1)\)\\
\(\text{Integrate: }\Phi_{\Delta\tau}^{-1}=\Phi_{-\Delta\tau}\text{，每一步用群指数保持 }U\in SU(N_c)\text{，且 }\det\partial\Phi=1\)\\
\(\text{Linear solve: }Q_f^2X_f=\phi_f\text{ 用 CG；残差 }r_f=\phi_f-Q_f^2X_f\text{ 满足 }\|r_f\|/\|\phi_f\|<\varepsilon_{\rm solve}\)\\
\(\text{Proposal: }(U',\Pi')=\Phi_{\Delta\tau}^{N_{\rm step}}(U,\Pi)\text{，记录 }\Delta H=H[U',\Pi',\phi]-H[U,\Pi,\phi]\)\\
\(\text{Accept: }U\gets U'\text{ with probability }P_{\rm acc}=\min(1,\exp[-\Delta H]);\text{ 否则保留旧 }U\text{ 并反转动量}\)\\
\(\mathrm{DB:}\;H(X')\text{ 与 }H(X)\text{ 按 Metropolis 详细平衡配对}\)\\
\(\text{Thermalization: }\text{丢弃初始轨道，按测得的 }\tau_{\rm int}\text{ 留出足够间隔；不能由轨道数本身假定独立}\)\\
\(\text{Ensemble checks: }\langle\Delta H\rangle\text{、接受率、plaquette、拓扑荷、残差与边界条件均落盘}\)\\
\(\quad U^\dagger U=I,\qquad \det U=1\text{；不满足群值检查则拒绝该组态}\)\\
\(\text{Output: }\{U^{(n)}\}\)，附作用量、残差、接受率和自相关元数据，供后续观测量计算使用。\\
\end{longtable}
轨道接受步骤的详细平衡关系必须写成带动量反演的形式：
\begin{equation}
 \mathcal P(X\to X')e^{-H(X)}
 =\mathcal P(\mathcal R X'\to\mathcal R X)e^{-H(X')},
 \qquad X=(U,\Pi,\phi),\quad \mathcal R(U,\Pi,\phi)=(U,-\Pi,\phi).
 \label{eq:deep-hmc-detailed-balance}
\end{equation}
对称、可逆且体积保持的积分器只保证提议核的这一结构；Metropolis 因子
才把离散积分误差从目标分布中纠正回来。

\begin{longtable}{@{}P{0.94\textwidth}@{}}
\caption{正定费米子核的 CG 线性求解单列算法}\label{tab:deep-cg}\\
\toprule
\textbf{输入、Krylov 方向、残差递推与收敛检查}\\
\midrule
\endfirsthead
\caption[]{正定费米子核的 CG 线性求解单列算法（续）}\\
\toprule
\textbf{输入、Krylov 方向、残差递推与收敛检查}\\
\midrule
\endhead
\bottomrule
\endfoot
\(\text{Input: }A=A^\dagger>0,\;b,\;x^{(0)},\;\varepsilon,\;N_{\max}\)\\
\(\text{Initialize: }r^{(0)}=b-Ax^{(0)},\;p^{(0)}=r^{(0)},\;\rho_0=r^{(0)\dagger}r^{(0)}\)\\
\(\text{Initial gate: }\|r^{(0)}\|\le\varepsilon_{\rm abs}+\varepsilon\|b\|\text{ 时直接输出 }x^{(0)}\text{；}b=0\text{ 时使用绝对残差}\)\\
\(\text{for }k=0,1,\ldots,N_{\max}-1\text{ iterate:}\)\\
\(\quad q^{(k)}=Ap^{(k)},\;\delta_k=p^{(k)\dagger}q^{(k)},\;\text{若 }\delta_k\le0\text{ 或非有限则停止并报错}\)\\
\(\quad \alpha_k=\rho_k/\delta_k,\;x^{(k+1)}=x^{(k)}+\alpha_kp^{(k)},\;r^{(k+1)}=r^{(k)}-\alpha_kq^{(k)}\)\\
\(\quad \text{若 }\|r^{(k+1)}\|_2\le\varepsilon_{\rm abs}+\varepsilon\|b\|_2\text{，输出 }x^{(k+1)}\text{ 并停止}\)\\
\(\quad \rho_{k+1}=r^{(k+1)\dagger}r^{(k+1)},\;\beta_k=\rho_{k+1}/\rho_k,\;p^{(k+1)}=r^{(k+1)}+\beta_kp^{(k)}\)\\
\(\text{end for: }\text{达到 }N_{\max}\text{ 仍未满足残差门槛则输出 }\texttt{not\_converged}\)\\
\(\text{Output: }x,\;r,\;\|r\|/\|b\|,\;k,\;\delta_k\text{ 与迭代元数据}\)\\
\end{longtable}
CG 的三个前提 \(A=A^\dagger>0\)、精确的算子作用和一致的残差定义缺一
不可；在浮点数中应同时抽查显式残差 \(b-Ax\)，因为递推残差可能因舍入
而失真。

\begin{longtable}{@{}P{0.94\textwidth}@{}}
\caption{右预条件 \(AM^{-1}\) 的 Bi--CGStab 单列算法}\label{tab:deep-bicgstab}\\
\toprule
\textbf{输入、影子残差、预条件方向、双正交更新与 breakdown 分支}\\
\midrule
\endfirsthead
\caption[]{右预条件 \(AM^{-1}\) 的 Bi--CGStab 单列算法（续）}\\
\toprule
\textbf{输入、影子残差、预条件方向、双正交更新与 breakdown 分支}\\
\midrule
\endhead
\bottomrule
\endfoot
\(\text{Input: }Ax=b,\;x^{(0)},\;\widetilde r,\;M,\;\varepsilon,\;N_{\max};\;M=I\text{ 时为无预条件版本}\)\\
\(\text{Initialize: }r^{(0)}=b-Ax^{(0)},\;\rho_0=\alpha_0=\omega_0=1,\;v^{(0)}=p^{(0)}=0\)\\
\(\text{Initial gate: }r^{(0)}=0\text{ 时返回；否则取 }\widetilde r=r^{(0)},\;M\text{ 可逆，}b=0\text{ 使用绝对残差}\)\\
\(\text{for }k=1,2,\ldots,N_{\max}\text{ iterate:}\)\\
\(\quad \rho_k=\widetilde r^\dagger r^{(k-1)};\;\rho_k=0\text{ 时为 breakdown，停止并报告}\)\\
\(\quad \beta_{k-1}=\frac{\rho_k}{\rho_{k-1}}\frac{\alpha_{k-1}}{\omega_{k-1}},\;p^{(k)}=r^{(k-1)}+\beta_{k-1}(p^{(k-1)}-\omega_{k-1}v^{(k-1)})\)\\
\(\quad M\widehat p^{(k)}=p^{(k)},\;v^{(k)}=A\widehat p^{(k)},\;d_k=\widetilde r^\dagger v^{(k)};\;d_k=0\text{ 时报告 breakdown；否则 }\alpha_k=\rho_k/d_k\)\\
\(\quad s^{(k)}=r^{(k-1)}-\alpha_kv^{(k)};\;\|s^{(k)}\|\le\varepsilon_{\rm abs}+\varepsilon\|b\|\text{ 时 }x^{(k)}=x^{(k-1)}+\alpha_k\widehat p^{(k)}\text{ 后停止}\)\\
\(\quad M\widehat s^{(k)}=s^{(k)},\;t^{(k)}=A\widehat s^{(k)},\;\omega_k=\frac{t^{(k)\dagger}s^{(k)}}{t^{(k)\dagger}t^{(k)}}\)\\
\(\quad t^{(k)\dagger}t^{(k)}=0\text{ 或 }\omega_k=0\text{ 时为 breakdown；否则 }r^{(k)}=s^{(k)}-\omega_kt^{(k)}\)\\
\(\quad x^{(k)}=x^{(k-1)}+\alpha_k\widehat p^{(k)}+\omega_k\widehat s^{(k)}\)\\
\(\quad \|r^{(k)}\|\le\varepsilon_{\rm abs}+\varepsilon\|b\|\text{ 时停止；否则保留 }\rho_k,\alpha_k,\omega_k\text{ 并进入下一轮}\)\\
\(\text{Output: }x,\;r,\;\text{显式残差、breakdown 状态、迭代次数与预条件器元数据}\)\\
\end{longtable}
该算法不要求 \(A\) Hermitian 正定，但要求影子残差与非零分母保持双正交
递推。此处先解 \(M\widehat p=p\)，再作用 \(A\)，故迭代算子是
\(AM^{-1}\)，残差仍为原系统的 \(b-Ax\)；左预条件 \(M^{-1}A\) 的残差
定义不同，不能混用。因而 HMC 中的 \(Q^2\)
若确为正定应优先使用 CG；将 Bi--CGStab 的“收敛”直接等同于 HMC
哈密顿量力的可逆性仍是不成立的物理—数值推断。

其中 CG 只对 Hermitian 正定核有保证；Hermitian 不定核可考虑 MINRES，
一般非 Hermitian 核可考虑 BiCGStab。多重位移 CG 则要求各系统恰为
\((A+\sigma_i I)x_i=b\)，不能把不同 Wilson 质量核无条件当作这一形式。
RHMC 是分数次行列式的有理近似构造，不是独立的线性求解器；
其位移系统与接受步骤必须重新核对。对现有 PyQCD，
\texttt{pipeline} 主要接收或处理已有组态，不能据此声称它已经生成了
满足式\eqref{eq:deep-gauge-only-partition} 的 dynamical ensemble。
\src{refer/books/Quantum_Chromodynamics_on_the_Lattice_latex/chapters/chapter08.tex；refer/books/INTRODUCTION_TO_LATTICE_QCD_latex/chapters/sec06_formulation_gauge.tex；../PyQCD/pyqcd/pipeline/}

\section{本章的可复核推理与极限}

把本章的推理压缩成一条可检查的链，可得到
\begin{align}
 D_\mu^G&=GD_\mu G^\dagger
 &&\Longrightarrow&& F_{\mu\nu}^G=GF_{\mu\nu}G^\dagger \\
 & &&\Longrightarrow&& U_\mu^G(x)=G(x)U_\mu(x)G^\dagger(x+a\hat\mu) \\
 & &&\Longrightarrow&& \operatorname{Tr}P_{\mu\nu}\text{ 与 }\bar\psi W\psi\text{ 规范不变} \\
 & &&\Longrightarrow&& S_g,\;Z\text{ 和规范 singlet 关联函数可在有限 }a\text{ 定义}.
 \label{eq:deep-gauge-chain}
\end{align}
在自由极限 \(g_0\to0\) 时，式\eqref{eq:deep-wilson-gauge-action} 的涨落
变成 Gaussian，式\eqref{eq:deep-wilson-free-kernel} 给出的 doubler 质量
计数必须恢复；在 \(a\to0\) 且保持物理长度 \(x=na\) 不变时，plaquette
展开必须恢复连续作用量；在 \(U_\mu=I\) 时，所有 staple 和 Wilson 线变成
单位矩阵。这三个极限同时通过，才足以把数组上的群值检查解释成格点场论
的一部分。它们仍不证明有限体积、有限质量和真实 Monte Carlo 采样已达到
连续 QCD。
\src{refer/books/INTRODUCTION_TO_LATTICE_QCD_latex/chapters/sec06_formulation_gauge.tex；refer/books/INTRODUCTION_TO_LATTICE_QCD_latex/chapters/sec08_continuum_limit.tex；refer/books/INTRODUCTION_TO_LATTICE_QCD_latex/chapters/sec09_fermion_doubling.tex；refer/books/Quantum_Chromodynamics_on_the_Lattice_latex/chapters/chapter09.tex}

\chapter{传递矩阵、反射正性与关联函数的谱理论}
\label{sec:deep-spectrum}

\section{Euclidean 时间为何能给出能谱}

Wick 旋转把 Minkowski 的时间演化 \(e^{-\mathrm{i}Ht}\) 变成 Euclidean
演化 \(e^{-Ht_E}\)。在一个时间层局部、满足适当边界条件的格点规范理论中，
传递矩阵定义为
\begin{equation}
 T=\exp(-aH),\qquad T^\dagger=T,\qquad T\ge0 .
 \label{eq:deep-transfer-matrix}
\end{equation}
这里正性不是形式上把指数写出来就自动成立，而是由
Osterwalder--Schrader 反射正性保证。若 \(\Theta\) 把时间 \(t>0\) 的场
反射到 \(t<0\)，则要求对所有只依赖正半时间区的泛函 \(F\)
\begin{equation}
 \langle(\Theta F)F\rangle\ge0 .
 \label{eq:deep-reflection-positivity}
\end{equation}
在局部规范作用量、Haar 测度、合适的时间反射和无化学势的条件下，这条
不等式允许从 Euclidean 相关函数构造正定 Hilbert 空间；若有化学势、
非局部时间作用量或不合适的边界，反射正性可能失效，此时不能直接把拟合
指数解释成物理能量。这就是谱拟合公式前面必须列出公理的原因。

在周期时间长度 \(\beta_T=N_ta\) 下，配分函数和两点函数可写为
\begin{align}
 Z&=\operatorname{Tr}T^{N_t},\\
 C_2(t,\boldsymbol p)
 &=\frac1Z\operatorname{Tr}\left[
 T^{N_t-t/a}\widehat O_{\boldsymbol p}
 T^{t/a}\widehat O_{\boldsymbol p}^\dagger\right].
 \label{eq:deep-trace-correlator}
\end{align}
插入能量本征态后，严格的有限温度式不是简单的“正向项加反向项”，而是
\begin{equation}
 C_2(t)=\frac1Z\sum_{m,n}e^{-E_m(\beta_T-t)}e^{-E_nt}
 \left|\langle m|\widehat O|n\rangle\right|^2 .
 \label{eq:deep-finite-temperature-trace}
\end{equation}
在低温、真空支配且算符具有确定的周期性符号时，才可用相对论归一化
\(\langle n,\boldsymbol p|n,\boldsymbol p\rangle=2E_nV\) 将它约化为
\(0<t<\beta_T/2\) 下的真空谱式
\begin{equation}
 C_2(t,\boldsymbol p)=\sum_n
 \frac{|Z_n(\boldsymbol p)|^2}{2E_n(\boldsymbol p)}
 \left[e^{-E_nt}+\eta_Oe^{-E_n(\beta_T-t)}\right],
 \qquad Z_n=\langle0|\widehat O|n,\boldsymbol p\rangle .
 \label{eq:deep-finite-time-spectrum}
\end{equation}
\(\eta_O=+1\) 或 \(-1\) 只在相应低温/反向传播通道可由一个符号概括时成立；
有限温度的热激发已被式\eqref{eq:deep-finite-temperature-trace}中的 \(m\ne0\)
项保留。无限时间近似只是丢掉热绕行和反向传播项后的特殊情况，不能对靠近
\(N_t/2\) 的数据无条件使用。

\section{有效质量不是经验平台}

若单态、周期边界且
\(C(t)=A[e^{-Et}+e^{-E(\beta_T-t)}]\)，则
\begin{align}
 C(t+a)+C(t-a)
 &=A e^{-Et}(e^{-aE}+e^{aE})
  +A e^{-E(\beta_T-t)}(e^{-aE}+e^{aE}) \nonumber\\
 &=2\cosh(aE)C(t),
 \label{eq:deep-cosh-recursion}\\
 aE_{\rm eff}^{\cosh}(t)
 &=\operatorname{arccosh}\left[
 \frac{C(t+a)+C(t-a)}{2C(t)}\right].
 \label{eq:deep-cosh-identity}
\end{align}
若反向传播项可忽略，同一谱分解退化为
\begin{equation}
 aE_{\rm eff}^{\log}(t)=\log\frac{C(t)}{C(t+a)} .
 \label{eq:deep-log-effective-energy}
\end{equation}
多态情况下写成
\(C(t)=A_0e^{-E_0t}(1+\sum_{n>0}r_ne^{-\Delta E_nt})\)，展开对数可得
\begin{equation}
 aE_{\rm eff}^{\log}(t)=aE_0+
 \sum_{n>0}r_n(1-e^{-a\Delta E_n})e^{-\Delta E_nt}
 +O(r^2e^{-2\Delta E t}) .
 \label{eq:deep-effective-contamination}
\end{equation}
因此平台表示指数修正已经小于统计误差和目标系统误差，而不是曲线肉眼
变平。若代码在 \(\operatorname{arccosh}\) 自变量小于 1 时置零，那是
数值掩码；物理报告必须记录掩码比例与窗口，不能把零解释成 \(E=0\)。

\section{GEVP 的变分推导}

取 \(N\) 个插值算符 \(O_i\)，相关矩阵为
\begin{equation}
 C_{ij}(t)=\sum_n\frac{Z_i^{(n)}Z_j^{(n)*}}{2E_n}e^{-E_nt},
 \qquad Z_i^{(n)}=\langle0|O_i|n\rangle .
 \label{eq:deep-correlation-matrix}
\end{equation}
在截断到 \(N\) 个态时令 \(Z_{in}=Z_i^{(n)}\)，
\(D(t)=\operatorname{diag}(e^{-E_nt}/(2E_n))\)，则
\begin{equation}
 C(t)=Z D(t)Z^\dagger,\qquad C(t_0)=ZD(t_0)Z^\dagger .
 \label{eq:deep-correlation-factorization}
\end{equation}
若 \(Z\) 满秩且 \(C(t_0)\) 正定，广义特征值方程
\begin{equation}
 C(t)v_n=\lambda_n(t,t_0)C(t_0)v_n
 \label{eq:deep-gevp}
\end{equation}
在约束 \(v^\dagger C(t_0)v=1\) 下使 \(v^\dagger C(t)v\) 取驻值，
拉格朗日函数给出完整的变分步骤
\begin{equation}
 \mathcal L=v^\dagger C(t)v-\lambda(v^\dagger C(t_0)v-1)
 \Longrightarrow\frac{\partial\mathcal L}{\partial v^\dagger}=0
 \Longrightarrow C(t)v=\lambda C(t_0)v .
 \label{eq:deep-gevp-variation}
\end{equation}
可左乘 \(Z^{-1}\) 并令 \(w_n=Z^\dagger v_n\)，得到
\begin{equation}
 D(t)w_n=\lambda_nD(t_0)w_n,\qquad
 \lambda_n(t,t_0)=e^{-E_n(t-t_0)} .
 \label{eq:deep-gevp-diagonal}
\end{equation}
有限基底且 (Z) 为方阵时上式是精确等式；有限基底而真实谱含高态时，
变分估计只能写成
\begin{equation}
 \lambda_n(t,t_0)=e^{-E_n(t-t_0)}
 \left[1+\epsilon_n(t,t_0)\right],
 \qquad \epsilon_n\xrightarrow[t,t_0\to\infty]{}0,
 \label{eq:deep-gevp-correction}
\end{equation}
其衰减阶数依赖 \(t,t_0\) 的联合取极限以及算符子空间，不能在没有说明
取极限的情况下把一个固定的 \(\Delta E_n\) 写成普适定律。这个推导的
必要条件是 \(C(t_0)\) 的正定性、算符重叠矩阵满秩和传递矩阵谱的正性；
将 \(C\) 替换为 \((C+C^\dagger)/2\) 只能去掉复数舍入的反 Hermitian
部分，不能修复真正的负本征值。

\section{三点函数、矩阵元和比值归一化}

设源在 0、插入在 \(\tau\)、汇在 \(t\)，初末动量可不同。插入两次完备
态得到
\begin{equation}
 C_3(t,\tau)=\sum_{m,n}
 \frac{Z_m^f Z_n^{i*}}{4E_m^fE_n^i}
 e^{-E_m^f(t-\tau)}e^{-E_n^i\tau}
 \langle m,f|\mathcal O|n,i\rangle .
 \label{eq:deep-three-point-spectrum}
\end{equation}
基态项和两点基态项分别是
\begin{align}
 C_3^{00}(t,\tau)&=
 \frac{Z_0^fZ_0^{i*}}{4E_fE_i}
 e^{-E_f(t-\tau)}e^{-E_i\tau}\mathcal M_{fi},
 \qquad \mathcal M_{fi}=\langle f|\mathcal O|i\rangle,\\
 C_2^f(s)&=\frac{|Z_0^f|^2}{2E_f}e^{-E_fs},\qquad
 C_2^i(s)=\frac{|Z_0^i|^2}{2E_i}e^{-E_is}.
 \label{eq:deep-two-three-ground}
\end{align}
将这些式子代入对称比值
\begin{equation}
 R(t,\tau)=\frac{C_3(t,\tau)}{C_2^f(t)}
 \sqrt{\frac{C_2^i(t-\tau)C_2^f(\tau)C_2^f(t)}
 {C_2^f(t-\tau)C_2^i(\tau)C_2^i(t)}} ,
 \label{eq:deep-symmetric-ratio}
\end{equation}
逐项消去 \(Z_i,Z_f\) 和指数，得到
\begin{equation}
 R(t,\tau)\xrightarrow[\tau,\,t-\tau\to\infty]{}
 \frac{\mathcal M_{fi}}{2\sqrt{E_fE_i}} .
 \label{eq:deep-ratio-limit}
\end{equation}
若只保留最低激发态，有限时间结构为
\begin{align}
 R(t,\tau)&=R_0+A e^{-\Delta E_i\tau}
 +B e^{-\Delta E_f(t-\tau)}
 +C e^{-\Delta E_i\tau-\Delta E_f(t-\tau)}+\cdots .
 \label{eq:deep-ratio-excited}
\end{align}
比值消灭了基态重叠和运动学指数，却没有消灭激发态。\(t_{\rm sep}\)
扫描、两态拟合和 summation method 仍是矩阵元可信度的必要证据。

对 summation method，把插入时间的可用格点求和：
\begin{align}
 S(t)&=a\sum_{\tau=a}^{t-a}R(t,\tau)\\
 &=c+t\,\frac{\mathcal M_{fi}}{2\sqrt{E_fE_i}}
 +O\left(t e^{-\Delta E t}\right).
 \label{eq:deep-summation-method}
\end{align}
因而相邻 \(t\) 的斜率给矩阵元，而常数 \(c\) 吸收端点效应。若斜率
随最小 \(t\) 漂移，说明高态项仍然与目标误差同量级。

\section{有限体积、动量量子化和多强子能级}

空间周期边界把平移算符的本征值量子化为
\begin{equation}
 p_i=\frac{2\pi n_i}{L_i},\qquad n_i\in\mathbb Z,\qquad
 \boldsymbol p^2=\sum_i p_i^2 .
 \label{eq:deep-momentum-quantization}
\end{equation}
单粒子自由色散关系在连续理论是 \(E^2=M^2+\boldsymbol p^2\)，在格点上
则含有
\begin{equation}
 4\sinh^2\frac{aE}{2}
 =4\sinh^2\frac{aM}{2}
 +\sum_i4\sin^2\frac{ap_i}{2}
 \label{eq:deep-lattice-dispersion}
\end{equation}
这样的格距修正；采用何种离散动量是拟合模型的一部分。有限体积的单粒子
指数修正一般由绕盒传播的最低质量态控制，量级为
\begin{equation}
 \delta Q(L)\sim\sum_{\boldsymbol n\ne0}
 \frac{e^{-m_\pi|\boldsymbol n|L}}{(m_\pi|\boldsymbol n|L)^\alpha},
 \label{eq:deep-finite-volume}
\end{equation}
其中幂指数 \(\alpha\) 取决于观测量和维数。对两粒子态，有限体积不只是
指数小修正：相互作用会把离散能级移动，能级与散射相移通过 Lüscher
型量子化条件联系。以 \(s\) 波为例，
\begin{equation}
 p\cot\delta_0(p)=\frac{2}{\sqrt{\pi}L}
 Z_{00}\left(1;\frac{p^2L^2}{4\pi^2}\right),
 \qquad E_2=2\sqrt{m^2+p^2}.
 \label{eq:deep-luscher-condition}
\end{equation}
其前提是短程相互作用、周期盒和能级远离不受控的多粒子阈值。把
\(C_2\) 的高激发态简单叫作噪声会丢掉这一有限体积动力学信息。

\section{统计相关性与谱量可辨识性}

对于 Markov 链上的观测 \(X_s\)，定义
\begin{align}
 \Gamma_X(\Delta)&=\langle(X_s-\bar X)(X_{s+\Delta}-\bar X)\rangle,\\
 \rho_X(\Delta)&=\frac{\Gamma_X(\Delta)}{\Gamma_X(0)},\\
 \tau_{\rm int,X}&=\frac12+\sum_{\Delta=1}^{W}\rho_X(\Delta),
 \qquad N_{\rm eff}\simeq\frac{N}{2\tau_{\rm int,X}} .
 \label{eq:deep-autocorrelation}
\end{align}
窗口 \(W\) 的选择本身带系统误差；把每个配置当作独立样本会低估
拓扑荷、长 Wilson 线和 disconnected loop 的误差。删除一项 Jackknife 的
无偏协方差为
\begin{equation}
 \widehat{\operatorname{Cov}}_J
 =\frac{N-1}{N}\sum_{s=1}^{N}
 (\boldsymbol x_{(-s)}-\bar{\boldsymbol x})
 (\boldsymbol x_{(-s)}-\bar{\boldsymbol x})^\dagger .
 \label{eq:deep-jackknife-covariance}
\end{equation}
相关拟合则最小化
\begin{equation}
 \chi^2(\theta)=
 [\boldsymbol y-\boldsymbol f(\theta)]^\dagger
 C^{-1}[\boldsymbol y-\boldsymbol f(\theta)] .
 \label{eq:deep-correlated-fit}
\end{equation}
若样本数小于数据点数，\(C\) 的经验秩最多为 \(N-1\)，其逆并非普通矩阵
意义上的稳定逆。SVD 截断、shrinkage 或 Bayesian prior 都是不同的统计
模型；加入 prior 后，参数受数据约束的程度必须和 prior 约束分开报告。

\begin{longtable}{@{}P{0.94\textwidth}@{}}
\caption{从逐组态相关函数到基态矩阵元的谱提取单列算法}
\label{tab:deep-spectrum-algorithm}\\
\toprule
\textbf{谱分解、重采样、激发态控制和输出}\\
\midrule
\endfirsthead
\caption[]{从逐组态相关函数到基态矩阵元的谱提取单列算法（续）}\\
\toprule
\textbf{谱分解、重采样、激发态控制和输出}\\
\midrule
\endhead
\bottomrule
\endfoot
\(\text{Input: }C_2^{(s)}(t,p),\;C_3^{(s)}(t,\tau,p_f,p_i),\;a,\;N_t,\;\text{周期/反周期边界，算符投影和拟合窗口}\)\\
\(\text{Index audit: }\text{确认配置轴、源时轴、汇时轴、插入轴和动量轴各自唯一，且 }t,\tau,t-\tau\text{ 的单位一致}\)\\
\(\text{Boundary map: }\Delta t=(t_{\rm sink}-t_{\rm src})\bmod N_t\text{；跨边界时按费米子/玻色子条件施加符号}\)\\
\(\text{Resample: }C^{(s)}\to C_{(-s)}\text{ 或 bootstrap 样本；对 }C_2,C_3,O\text{ 使用同一索引集合，保留样本相关性}\)\\
\(\text{Spectrum: }C_2(t)\to E_{\rm eff}^{\log}\text{ 或 }E_{\rm eff}^{\cosh}\text{；若矩阵基底可辨识则同时形成 }C(t)v=\lambda C(t_0)v\)\\
\(\text{GEVP gate: }C(t_0)=C(t_0)^\dagger>0\text{，本征值排序和 }t_0\text{ 约定固定；否则输出 }\texttt{unidentifiable}\text{ 并停止升级}\)\\
\(\text{Two-state model: }C_2=A_0e^{-E_0t}+A_1e^{-E_1t}\text{；三点函数保留 }(m,n)=(0,0),(1,0),(0,1),(1,1)\text{ 项}\)\\
\(\text{Ratio: }R\gets C_3/C_2\times\text{式 }\eqref{eq:deep-symmetric-ratio}\text{ 的运动学因子；逐样本计算而不是先独立取商}\)\\
\(\text{Plateau test: }\text{扫描 }t_{\rm sep},\tau_{\rm min},\tau_{\rm max}\text{，比较平台、两态拟合和求和斜率}\)\\
\(\text{Covariance gate: }\operatorname{rank}(C),\;\kappa(C),\;N_{\rm data},\;N_{\rm par}\text{ 全部记录；SVD cut 作为系统扫描变量}\)\\
传播误差：在每个重采样样本上重做能量、GEVP、ratio 和拟合；保留联合样本及非对角协方差。\\
\(\text{Checks: }\text{自由场色散、时间反射、动量反演、局域 }z=0\text{ 通道和不同拟合窗的一致性逐项检查}\)\\
\(\text{Output: }E_n,Z_n,\mathcal M_{fi},\text{协方差、有效秩、激发态模型、窗口、边界、投影与重采样元数据}\)\\
\end{longtable}
\src{refer/books/INTRODUCTION_TO_LATTICE_QCD_latex/chapters/sec05_minkowski_euclidean.tex；refer/books/INTRODUCTION_TO_LATTICE_QCD_latex/chapters/sec16_errors.tex；refer/books/INTRODUCTION_TO_LATTICE_QCD_latex/chapters/sec17_correlators.tex；refer/books/INTRODUCTION_TO_LATTICE_QCD_latex/chapters/sec18_hadron_spectrum.tex；refer/books/Quantum_Chromodynamics_on_the_Lattice_latex/chapters/chapter04.tex；../PyQCD/pyqcd/analysis/_analyse.py；../PyQCD/pyqcd/analysis/_fitter.py}

\section{定理、假设与可验证边界的对应}

\begin{table}[htbp]
\centering
\caption{谱理论中公理到结论的依赖关系}
\small
\begin{tabularx}{0.98\textwidth}{@{}P{3.0cm}P{4.1cm}Y@{}}
\toprule
前提 & 推导结论 & 不能由有限维代数单独证明的部分\\
\midrule
局部作用量、Haar 测度、反射正性 & \(T\) 正、自伴，Euclidean 指数可解释为 \(E_n\) & 给定实现是否满足反射正性以及边界是否正确。\\
平移不变性、周期边界 & 动量投影、源时平均和有限时间双向传播 & 实际配置的时间轴是否被正确滚移。\\
\(C(t_0)>0\)、\(Z\) 满秩 & GEVP 本征值呈指数、变分基态提取 & 噪声、截断基底和有限统计下的有效秩。\\
Grassmann 高斯积分、\(D\) 可逆 & Wick 收缩和费米置换号 & 数值线性求解误差与近零模的统计放大。\\
平稳 Markov 链、有限 \(\tau_{\rm int}\) & \(N_{\rm eff}\) 与协方差的自相关修正 & 拓扑冻结或慢模是否已被测量窗口捕获。\\
\bottomrule
\end{tabularx}
\end{table}

因此 SymPy 对式\eqref{eq:deep-cosh-identity}、\eqref{eq:deep-ratio-limit}
和协方差恒等式的化简只能确证局部代数；真实 \(C_2,C_3\)、自相关时间、
GEVP 基底稳定性和 plateau 必须继续以逐组态数据和窗口扫描为证据。
\src{refer/books/INTRODUCTION_TO_LATTICE_QCD_latex/chapters/sec05_minkowski_euclidean.tex；refer/books/INTRODUCTION_TO_LATTICE_QCD_latex/chapters/sec16_errors.tex；refer/books/INTRODUCTION_TO_LATTICE_QCD_latex/chapters/sec17_correlators.tex；refer/books/INTRODUCTION_TO_LATTICE_QCD_latex/chapters/sec18_hadron_spectrum.tex}

\chapter{非局域算符、重整化、OPE 与 TMD 因子化的理论闭环}
\label{sec:deep-nonlocal}

\section{Wilson 线自能为何产生线性发散}

直线双局域算符的规范不变量定义为
\begin{equation}
 \mathcal O_\Gamma(z)=\bar\psi(0)\Gamma W(0,z\hat z)\psi(z\hat z),
 \qquad
 W(0,z\hat z)=\mathcal P\exp\left[
 \mathrm{i}g\int_0^z\dd\lambda\,A_z(\lambda\hat z)\right].
 \label{eq:deep-bilocal}
\end{equation}
它的 UV 结构不能仅由两个端点场的 \(Z_\psi\) 描述。一个清楚的论证是引入
沿路径传播的辅助静态场 \(Q(\lambda)\)，使 Wilson 线成为一维传播子：
\begin{equation}
 \mathcal L_Q=\bar Q(\lambda)
 [\partial_\lambda+\delta m+\mathrm{i}gA_z(\lambda)]Q(\lambda),
 \qquad
 W(0,z)\sim\langle Q(z)\bar Q(0)\rangle_A .
 \label{eq:deep-auxiliary-line}
\end{equation}
一维静态场的质量反项乘以传播长度；若将辅助线方向选成与
\(W(0,z)\) 相反的传播方向，则其背景传播子满足
\begin{equation}
 G_Q(z,0;A)\equiv\langle Q(z)\bar Q(0)\rangle_A
 =\theta(z)\,e^{-\delta m z}\,W_Q(z,0;A),
 \qquad W_Q(z,0;A)=W(0,z)\ \text{或其定向共轭}.
 \label{eq:deep-auxiliary-propagator}
\end{equation}
因此裸算符和重整化算符之间必含
\begin{equation}
 \mathcal O_\Gamma^{\rm bare}(z,a)
 =e^{-\delta m(a)|z|}
 Z_\Gamma^{\log}(z\mu,a\mu)
 \sum_{\Gamma'}\mathcal Z_{\Gamma\Gamma'}(z,a,\mu)
 \mathcal O_{\Gamma'}^{\rm R}(z,\mu).
 \label{eq:deep-nonlocal-renorm}
\end{equation}
格点硬截断使 \(\delta m(a)\sim c(g_0)/a\)，所以固定非零物理长度时
裸矩阵元会有 \(\exp[-c|z|/a]\) 的剧烈变化。维数正规化中这类 power
divergence 可被设为零，并不意味着格点正则化下它不存在。

若路径有一个 cusp，局部切向量从 \(v_1\) 跳到 \(v_2\)，一圈 Wilson 线图
还产生只依赖夹角的 cusp 反常维数。抽象地写作
\begin{equation}
 \mu\frac{\dd}{\dd\mu}\ln Z_{\rm cusp}(\theta,\mu)
 =-\Gamma_{\rm cusp}(\alpha_s,\theta),
 \label{eq:deep-cusp-rge}
\end{equation}
而 finite staple 至少含两个端点和若干 cusp。故直线 OPE 的 \(Z_R(z)\)、
staple TMD 的 \(Z_R(z,b_T,\ell)\) 和 soft factor 的 \(S(b_T)\) 不能用一个
没有几何元数据的标量 \(Z\) 混代。端点发散、线性反项、cusp 反项和普通
场波函数反项的乘法分解，只有在选定 regulator 与算符基底后才是有定义的
陈述。

\section{RI/MOM 与混合矩阵的定义}

令 \(\{\mathcal O_i\}\) 是有限格距对称允许的算符基底。对外部夸克动量
\(p\) 的 amputated vertex 定义
\begin{equation}
 \Lambda_i(p,z)=S^{-1}(p)G_i(p,z)S^{-1}(p),
 \qquad
 \mathcal M_{ij}(p,z)=\operatorname{Tr}[P_j\Lambda_i(p,z)],
 \label{eq:deep-ri-vertex}
\end{equation}
其中 \(P_j\) 是满足树级归一的投影。非微扰 RI/MOM 条件是
\begin{equation}
 Z_q^{-1}(p)\,Z_{ik}^{\rm RI}(z,p)
 \mathcal M_{kj}(p,z)\big|_{p^2=\mu_R^2}
 =\mathcal M_{ij}^{\rm tree}(p,z).
 \label{eq:deep-ri-matrix-condition}
\end{equation}
在树级投影已取为单位矩阵时，直接写成
\begin{equation}
 Z^{\rm RI}(z,\mu_R)=Z_q(\mu_R)
 [\mathcal M(z,\mu_R)]^{-1}.
 \label{eq:deep-ri-inverse-mixing}
\end{equation}
这里矩阵求逆的方向很重要：若把 \(\mathcal M^{-1}\) 乘在错误一侧，
会把算符指标和投影指标互换，导致 off-diagonal mixing 的符号和位置
同时错误。

对于 Wilson 费米子、\(\Gamma=\gamma_z\) 的 unpolarized quasi-PDF，
有限格距手征破缺允许 scalar 通道进入同一重整化问题：
\begin{equation}
 \begin{pmatrix}O_V^R\\O_S^R\end{pmatrix}
 =\begin{pmatrix}Z_{VV}&Z_{VS}\\Z_{SV}&Z_{SS}\end{pmatrix}
 \begin{pmatrix}O_V^{\rm bare}\\O_S^{\rm bare}\end{pmatrix}.
 \label{eq:deep-vector-scalar-mixing}
\end{equation}
因此只计算 \(h_V^{\rm bare}\) 而忽略 \(h_S^{\rm bare}\) 并不是一个小的
统计近似；它相当于把一个允许的算符方向从重整化基底中删掉。若采用
\(\gamma_t\) 或具有更强手征对称的离散化，混合结构可以改变，但必须由
对称性和独立 vertex 条件证明，而不能靠命名推断。

RI/MOM 到 \(\overline{\rm MS}\) 的转换写成
\begin{equation}
 \mathcal O_i^{\overline{\rm MS}}(\mu)
 =C_{ij}^{\overline{\rm MS}\leftarrow\rm RI}(\mu,\mu_R)
 \mathcal O_j^{\rm RI}(\mu_R),
 \qquad
 C=I+\frac{\alpha_s}{4\pi}C^{(1)}+O(\alpha_s^2).
 \label{eq:deep-scheme-conversion}
\end{equation}
\(C^{(1)}\) 取决于 projector、外部动量方向、规范固定、Wilson 线长度和
算符基底。将不同 RI momentum orbit 的数据直接平均，会把离散化误差误当
成方案独立性。实际的窗口要求同时满足
\begin{equation}
 \Lambda_{\rm QCD}^2\ll\mu_R^2\ll a^{-2},
 \qquad
 \mu_R|z|\ \text{不能过大且离散方向效应可控}.
 \label{eq:deep-ri-window}
\end{equation}
左边保证外部态的非微扰污染较小，右边保证格距伪影不主导；不存在重叠
窗口时，RI/MOM 数值可以定义，但不能被解释为可靠的连续方案转换。

\section{协变 Taylor 展开与局域 OPE}

straight Wilson 线的关键恒等式是协变输运后的场可以在一个端点展开：
\begin{equation}
 W(0,z)\psi(z\hat z)
 =\sum_{n=0}^{\infty}\frac{(\mathrm{i}z)^n}{n!}
 \bigl[(-\mathrm{i}D_z)^n\psi\bigr](0)
 \quad\text{（形式上，路径有序意义下）}.
 \label{eq:deep-covariant-taylor}
\end{equation}
代回式\eqref{eq:deep-bilocal} 得到
\begin{equation}
 \mathcal O_\Gamma(z)
 =\sum_{n=0}^{\infty}\frac{(\mathrm{i}z)^n}{n!}
 \bar\psi(0)\Gamma(-\mathrm{i}D_z)^n\psi(0).
 \label{eq:deep-bilocal-local-expansion}
\end{equation}
量子修正把裸局域算符换成按 twist 分类的重整化基底。对矢量通道，
令 \(O_{\Gamma,n}^{(2)}\) 为对 \(n+1\) 个 Lorentz 指标对称并去迹的
twist-2 算符（维数 \(3+n\)），则算符级展开的量纲一致形式为
\begin{equation}
 \mathcal O_\Gamma^R(z,\mu)
 =\sum_{n=0}^{\infty}\frac{(\mathrm{i}z)^n}{n!}
 C_n^{(2)}\,O_{\Gamma,n}^{(2),R}
 +z^2\sum_{n,k}\frac{(\mathrm{i}z)^n}{n!}
 C_{n,k}^{(4)}\,O_{\Gamma,n,k}^{(4),R}+\cdots .
 \label{eq:deep-local-ope}
\end{equation}
这里所有系数依赖 \(z^2\mu^2,\alpha_s\)，高 twist 算符及迹项的独立基底
用 \(k\) 区分；第二组算符的维数是 \(5+n\)。所以
\(z^2\Lambda_{\rm QCD}^2\) 只能描述剥离运动学因子后的相对幂修正，不能
直接加到维数为 3 的算符右端。具体地，矢量矩阵元分解为
\begin{align}
 \langle P|\bar\psi(0)\gamma^\mu W(0,z)\psi(z)|P\rangle_R
 &=2P^\mu\mathcal M(\nu,z^2,\mu)+2z^\mu\mathcal N(\nu,z^2,\mu),\\
 \langle P|O_{n}^{\{\mu\mu_1\cdots\mu_n\}}|P\rangle
 &=2a_n(\mu)\,[P^\mu P^{\mu_1}\cdots P^{\mu_n}-\text{迹项}].
\end{align}
选取满足 \(z^\mu=0\) 的 Lorentz 分量或显式分离第二个不变量，并除去 \(2P^\mu\) 后，
定义 \(\nu=zP_z\)，得到无量纲的 pseudo-ITD：
\begin{align}
 \mathcal M(\nu,z^2,\mu)
 &=\sum_{n=0}^{\infty}\frac{(\mathrm{i}\nu)^n}{n!}
 C_n(z^2\mu^2,\alpha_s)a_n(\mu)
 +O(z^2\Lambda_{\rm QCD}^2,z^2M^2).
 \label{eq:deep-itd-moment-expansion}
\end{align}
最后的 \(z^2M^2\) 来自非光锥分离的运动学迹项，原则上可单独重求和；
它与动力学的高 twist 修正不同。此处 \(z^2\) 表示正的空间分离平方，
Minkowski 写法相应使用 \(-z_{\rm M}^2\)。
其中 \(a_n\) 是 twist-2 分布的 Mellin 矩。若光锥 PDF 以 \(f(x,\mu)\)
定义，则
\begin{equation}
 a_n(\mu)=\int_{-1}^{1}\dd x\,x^n f(x,\mu).
 \label{eq:deep-mellin-moment}
\end{equation}
负 \(x\) 区域承载反夸克约定；把积分范围强行改成 \([0,1]\) 会丢失
偶/奇矩与电荷共轭信息。\(n=0\) 的归一化还受到局域流守恒约束，
它是检查 Fourier、matching 和端点 delta 项是否相容的重要不变量。

\section{LaMET 匹配和有限 Fourier 区间}

对沿 \(z\) 方向有大动量 \(P_z\) 的等时矩阵元，准分布定义为
\begin{equation}
 \widetilde q(x,P_z,\mu)=\int_{-\infty}^{\infty}
 \frac{\dd z}{2\pi}\,e^{-\mathrm{i}xzP_z}h_R(z,P_z,\mu).
 \label{eq:deep-quasi-fourier}
\end{equation}
LaMET 的 leading-power 因子化不是 \(\widetilde q=q\)，而是
\begin{equation}
 \widetilde q(x,P_z,\mu)=\int_{-1}^{1}\frac{\dd y}{|y|}
 C\left(\frac{x}{y},\frac{\mu}{|y|P_z}\right)q(y,\mu)
 +\frac{\Lambda_{\rm QCD}^2}{P_z^2}H_2(x)+O(P_z^{-4}).
 \label{eq:deep-lamet-factorization}
\end{equation}
绝对值来自负 \(y\) 反夸克分支的 Jacobian。匹配核的 plus 分布必须先定义
其作用：对在 \([0,1]\) 上可积的测试函数 \(\varphi\)，
\begin{equation}
 \int_0^1\dd\xi\,[f(\xi)]_+\varphi(\xi)
 =\int_0^1\dd\xi\,f(\xi)[\varphi(\xi)-\varphi(1)].
 \label{eq:deep-plus-definition}
\end{equation}
若还含端点 delta 项，必须把虚拟图贡献一起写入，才能检查守恒矩
\begin{equation}
 \int_{-\infty}^{\infty}\dd x\,\widetilde q(x,P_z,\mu)
 =\int_{-1}^{1}\dd y\,q(y,\mu)
 \label{eq:deep-number-sum-rule}
\end{equation}
在相应非 singlet 矢量流约定下成立。对于胶子分布，守恒约束一般由
动量 sum rule 与 quark singlet--gluon mixing 共同表达，不能机械套用
非 singlet 的零阶归一化。

格点上只能得到 \(z\in[-z_{\max},z_{\max}]\)。截断 Fourier 的精确表达是
\begin{align}
 \widetilde q_{z_{\max}}(x)
 &=\int_{-z_{\max}}^{z_{\max}}\frac{\dd z}{2\pi}
 e^{-\mathrm{i}xzP_z}h_R(z)\\
 &=\int\dd y\,\widetilde q_\infty(y)
 \frac{\sin[(x-y)P_zz_{\max}]}{\pi(x-y)}.
 \label{eq:deep-truncated-fourier-kernel}
\end{align}
第二行是 Dirichlet 核卷积，说明有限长度会造成端点振铃与
\(\Delta x\sim1/(P_zz_{\max})\) 的分辨率。零填充可以细化输出网格，
不能提高物理信息量；改变窗函数会改变卷积核，必须作为系统误差而非
美化步骤报告。

\section{finite-staple 胶子 TMD 与 soft factor}

胶子场强的 gauge invariant 双局域结构需要把两个 field strength 用同一路径
连接。令 \(y=x+z\hat z+b_T\hat b\)，staple 路径由纵向臂、横向臂和回程
臂组成，记为
\begin{equation}
 W_\sqcup(x,y;\ell)=W_z(x,x+\ell\hat z)
 W_\perp(x+\ell\hat z,x+\ell\hat z+b_T\hat b)
 W_z(x+\ell\hat z+b_T\hat b,y),
 \qquad y=x+z\hat z+b_T\hat b .
 \label{eq:deep-staple-geometry}
\end{equation}
相应规范不变算符的强子矩阵元为
\begin{equation}
 M^{\mu\nu;\rho\sigma}(z,b_T;\ell)
 =\left\langle P\left|
 \operatorname{Tr}\left[F^{\mu\nu}(x)W_\sqcup(x,y)
 F^{\rho\sigma}(y)W_\sqcup^\dagger(x,y)\right]
 \right|P\right\rangle .
 \label{eq:deep-gluon-bilocal}
\end{equation}
自旋平均、非极化通道需要固定投影。例如在一组选定横向方向 \(i,j\) 下，
可以定义与当前对照实现一致的组合
\begin{equation}
 \mathcal O_g^{\rm unpol}
 =M^{ti;ti}+M^{tj;tj}-2M^{ij;ij}.
 \label{eq:deep-gluon-unpolarized-projector}
\end{equation}
系数和整体归一化依赖场强、Euclidean projector、link class 和 \(2E\)
约定，不能脱离这些元数据复用。研究目标 \(f_1^{g[-,-]}\) 时，未来向/
过去向 link class、核子自旋平均和投影必须在定义层固定。

为使“非极化胶子分布”不依赖口头命名，先给出光锥侧的定义约定。取
\(n^2=\bar n^2=0,\;n\cdot\bar n=1\)，横向指标 \((i,j=1,2)\)，并把
\([-,-]\) 的两条过去向 Wilson 线写成伴随表示
\(\mathcal U_{ab}^{[-,-]}=2\operatorname{Tr}[T^aU^{[-,-]}T^bU^{[-,-]\dagger}]\)。
则自旋平均的未减除相关张量定义为
\begin{align}
 \Phi_{g,\rm unsub}^{ij}(x,b_T;\delta)
 &=\frac{1}{xP^+}\int\frac{\dd\xi^-}{2\pi}
 e^{-\mathrm{i}xP^+\xi^-}\frac12\sum_s
 \langle P,s|F_a^{+i}(\xi^-,b_T)\,
 \mathcal U_{ab}^{[-,-]}\,F_b^{+j}(0)|P,s\rangle_\delta,\\
 \Phi_g^{ij}(x,b_T)&=\frac12\delta_T^{ij}f_1^{g[-,-]}(x,b_T)+
 \left(\frac{b_T^ib_T^j}{b_T^2}-\frac12\delta_T^{ij}\right)
 h_1^{\perp g[-,-]}(x,b_T).
 \label{eq:deep-gluon-tmd-definition}
\end{align}
这里第一式的 \(1/(xP^+)\)、颜色伴随表示、核子自旋平均和第二式的
张量分解共同构成定义；因此在本归一化下
\(f_1^{g[-,-]}=\delta_{T,ij}\Phi_g^{ij}\)。Euclidean 的
\(F^{ti},F^{ij}\) 组合只有在给定 Wick 旋转、动量投影、\(2E\) 归一化及
匹配系数后，才可与式\eqref{eq:deep-gluon-tmd-definition}对应，不能把数组
组合直接标成光锥 PDF。

TMD 的快度发散来自共线和软模式在虚度上重叠而在快度上不同。引入同一
rapidity regulator 后，未减软矩阵元 \(F_{\rm unsub}\) 与真空 soft factor
\(S\) 的物理组合可抽象写成
\begin{equation}
 F(x,b_T;\mu,\zeta)
 =Z_F(\mu,\zeta,\delta)
 \frac{F_{\rm unsub}(x,b_T;\delta)}
 {\sqrt{S(b_T;\mu,\delta)}} .
 \label{eq:deep-tmd-soft-subtraction}
\end{equation}
\(\delta\) 是中间调节器，最终左侧不应依赖它；soft 因子必须使用同一 Wilson
线方向、cusp 几何和 regulator。一个由准 TMD 波函数或 form factor 构造的
intrinsic soft ratio，例如
\begin{equation}
 S_I(b_T)=\frac{R^2(b_T)}{R^2(b_{T,0})},
 \label{eq:deep-intrinsic-soft-ratio}
\end{equation}
只有在其因子化定理、参考距离和同几何条件均成立时，才可作为 soft 因子的
方案化估计；它不是任意把 \(b_T=0\) 的数组除过去。

\section{Collins--Soper 方程的可积性}

软减除后的 TMD 同时依赖 \(\mu\) 和 \(\zeta\)。采用本报告的约定：
\begin{align}
 \frac{\partial}{\partial\ln\sqrt\zeta}\ln F(x,b_T;\mu,\zeta)
 &=K(b_T,\mu),
 \label{eq:deep-cs-equation}\\
 \mu\frac{\partial}{\partial\mu}\ln F(x,b_T;\mu,\zeta)
 &=\gamma_F\left(\alpha_s(\mu),\frac{\zeta}{\mu^2}\right).
 \label{eq:deep-tmd-mu-rge}
\end{align}
两次微分必须可交换，故
\begin{equation}
 \mu\frac{\dd}{\dd\mu}K(b_T,\mu)
 =\frac{\partial\gamma_F}{\partial\ln\sqrt\zeta}
 =-\gamma_K(\alpha_s(\mu)),
 \qquad \gamma_K=2\Gamma_{\rm cusp}.
 \label{eq:deep-cs-integrability}
\end{equation}
最后等号的符号随 \(\gamma_F\) 定义而定，本报告固定为式
\eqref{eq:deep-cs-integrability}。在小 \(b_T\) 的微扰窗口，可写
\begin{equation}
 K(b_T,\mu)=K(b_T,\mu_b)
 -\int_{\mu_b}^{\mu}\frac{\dd\mu'}{\mu'}
 \gamma_K(\alpha_s(\mu')),
 \qquad \mu_b\sim b_T^{-1}.
 \label{eq:deep-cs-solution}
\end{equation}
而 \(b_T\gtrsim\Lambda_{\rm QCD}^{-1}\) 时初值本身是非微扰函数；在一圈
固定耦合近似下，
\begin{equation}
 \Gamma_{\rm cusp}^{(1)}=\frac{\alpha_s C_R}{\pi},
 \qquad \gamma_K^{(1)}=2\Gamma_{\rm cusp}^{(1)},
 \qquad K^{(1)}(b_T,\mu)=-\frac{\alpha_s C_R}{\pi}
 \ln\frac{\mu^2b_T^2e^{2\gamma_E}}4+K_0 .
 \label{eq:deep-cs-one-loop}
\end{equation}
其中 \(K_0\) 是边界方案常数。两个大动量
的准 TMD 比值在 leading power 下提供
\begin{equation}
 K(b_T,\mu)=\frac{1}{\ln(P_1^z/P_2^z)}
 \ln\left[
 \frac{H(P_2^z,\mu)\,\widetilde F_R(P_1^z,b_T)}
 {H(P_1^z,\mu)\,\widetilde F_R(P_2^z,b_T)}
 \right]
 +O\left(\frac{\Lambda^2}{(P_1^z)^2},
 \frac{\Lambda^2}{(P_2^z)^2}\right).
 \label{eq:deep-two-momentum-cs}
\end{equation}
此式要求两个动量使用共同 \(b_T\)、共同 \(z\) 参考点、共同 soft 方案和
共同重整化；只改变 \(P_z\) 数值而改变 \(z\) 网格或 staple 长度，所得
对数比值不是纯 CS 演化。若采用文献中
\(2\zeta\,\partial_\zeta\) 的记号，则它与本报告的
\(\partial_{\ln\sqrt\zeta}\) 完全等价；报告中不能同时采用两种记号却
省略二者的换算。

\section{TMD—collinear 联系和流窗口}

在小横向距离窗口，soft-subtracted TMD 具有 small-\(b_T\) OPE：
\begin{equation}
 F_g(x,b_T;\mu,\zeta)
 =\sum_{j\in\{g,q,\bar q\}}
 \int_x^1\frac{\dd\xi}{\xi}
 C_{gj}\left(\frac{x}{\xi},b_T;\mu,\zeta\right)f_j(\xi,\mu)
 +O(b_T^2\Lambda_{\rm QCD}^2).
 \label{eq:deep-small-bt-ope}
\end{equation}
胶子算符一般与 quark singlet 混合，因此只用 \(g\to g\) 标量系数不能在
没有 flavor basis 说明时称作完整 matching。反过来，梯度流定义
\begin{equation}
 \partial_tB_\mu=D_\nu G_{\nu\mu},
 \qquad B_\mu(0,x)=A_\mu(x),
 \label{eq:deep-flow-equation}
\end{equation}
在线性化的 Abelian 极限给出
\begin{equation}
 B_\mu(t,p)=e^{-tp^2}A_\mu(p)+\text{规范纵向项},
 \qquad r_{\rm smear}=\sqrt{8t}.
 \label{eq:deep-flow-heat-kernel}
\end{equation}
故流能抑制 \(p^2\gg1/t\) 的 UV 模式。合法的局域 SFTE 窗口至少要求
\begin{equation}
 a^2\ll t\ll\Lambda_{\rm QCD}^{-2}.
 \label{eq:deep-flow-window}
\end{equation}
对于非局域算符，短距离不能只看一个 \(z\) 或 \(b_T\)。若 \(d_\Gamma\)
是整条 Wilson 路径以及端点到源、汇和最近时间边界的最短物理距离（扣除
smearing 支撑），则应检查
\begin{equation}
 a\ll\sqrt{8t}\ll d_\Gamma,\qquad
 a\ll |z|,b_T\ll\Lambda_{\rm QCD}^{-1},
 \label{eq:deep-nonlocal-flow-window}
\end{equation}
并另查周期绕回或开放边界污染。缺少路径支撑窗口时，flowed finite-staple
只能叫 finite-flow prototype；局域 SFTE 不能代替 cusp、线性发散、soft
和 rapidity renormalization。

\begin{longtable}{@{}P{0.94\textwidth}@{}}
\caption{从 flowed gauge field 到软减除、CS 核和胶子 TMD 的完整单列算法}
\label{tab:deep-tmd-algorithm}\\
\toprule
\textbf{几何定义、逐组态测量、共同重采样、匹配与停止门}\\
\midrule
\endfirsthead
\caption[]{从 flowed gauge field 到软减除、CS 核和胶子 TMD 的完整单列算法（续）}\\
\toprule
\textbf{几何定义、逐组态测量、共同重采样、匹配与停止门}\\
\midrule
\endhead
\bottomrule
\endfoot
\(\text{Input: }U^{(s)},a,L,N_t,P_z,\tau,z,b_T,\ell,\text{ link class }[-,-],\text{ Euclidean projector},\text{ 颜色归一化}\)\\
\(\text{Geometry gate: }\text{固定 }x,y,\hat z,\hat b,\ell\text{ 和 future/past 方向；枚举路径段并检查每一段端点相容}\)\\
\(\text{Flow: }U^{(s)}\to V^{(s)}(\tau)\text{，保存积分步长、误差容限和 }\sqrt{8\tau}\text{；若不流则显式写 }\tau=0\text{ 方案}\)\\
\(\text{Field strength: }F_{\mu\nu}^{(s)}\gets\text{Clover 四叶平均，}\;F\to F-\operatorname{Tr}(F)I/N_c\text{（若通道要求无迹投影）}\)\\
\(\text{Staple: }x\to x+\ell\hat z\to x+\ell\hat z+b_T\hat b\to y=x+z\hat z+b_T\hat b\text{，逐链接保持矩阵乘法顺序}\)\\
\(\text{Operator: }M^{\mu\nu;\rho\sigma}\gets\operatorname{Tr}[F^{\mu\nu}(x)W_\sqcup F^{\rho\sigma}(y)W_\sqcup^\dagger]\text{，按式 }\eqref{eq:deep-gluon-unpolarized-projector}\text{ 投影}\)\\
\(\text{Nucleon insertion: }\text{逐组态生成 }C_2^{(s)},\;L^{(s)}(\tau,z,b_T),\;C_2^{(s)}L^{(s)}\text{；不先对两因子分别平均}\)\\
\(\text{Disconnected subtraction: }C_{3,\rm disc}^{(s)}=C_2^{(s)}L^{(s)}-C_2^{(s)}\,\langle L\rangle_{\rm vac}\text{，真空项在同一重采样中求取}\)\\
\(\text{Ratio/ground state: }R^{(s)}\gets C_{3,\rm disc}^{(s)}/C_2^{(s)}\text{ 或完整运动学比值；扫描 }t_{\rm sep},\tau\text{ 并保留激发态模型}\)\\
\(\text{Renormalize: }h_R\gets Z_Rh_{\rm bare}\text{，选择 ratio、RI/MOM matrix、hybrid 或 flow/SFTE，完整记录 }(z,b_T,\ell,\tau,\mu)\)\\
\(\text{Soft: }S\gets S_{\rm same\ geometry}/\text{参考量，或 }S_I\text{ 因子化估计；检查 regulator、cusp、方向和流方案完全相同}\)\\
\(\text{CS: }K\gets\ln[H(P_2)h_R(P_1)/(H(P_1)h_R(P_2))]/\ln(P_1/P_2)\text{，至少三个有效 }P_z\text{ 做共同 }z,b_T\text{ 检查}\)\\
\(\text{Fourier: }h_R(z,b_T)\to x\widetilde g(x,b_T,P_z)\text{，使用 Hermitian 延拓、明确窗函数和有限 }z_{\max}\text{ 核}\)\\
\(\text{Matching: }x\widetilde g\to F_g\text{，构造 }Z_{ij}=\delta_{ij}+\alpha_sM_{ij}\text{，加入 soft、CS、胶子 singlet mixing 和卷积测度}\)\\
\(\text{Joint extrapolation: }Q(a,L,P_z,\tau,\ell)=Q_0+c_a a^p+c_L e^{-m_\pi L}+c_P/P_z^2+c_\tau\tau+\cdots\text{，共同拟合并扫描先验}\)\\
\(\text{Final gate: }a\ll\sqrt{8\tau}\ll d_\Gamma,\;m_\pi L\text{ 足够大，soft/rapidity/Z/matching/协方差均有证据；否则状态降级}\)\\
\(\text{Output: }f_1^{g[-,-]}(x,b_T;\mu,\zeta),K(b_T,\mu),S(b_T),\text{匹配矩阵、全协方差、外推和证据状态}\)\\
\end{longtable}
\src{refer/papers/从准TMD波函数确定CS核_latex/chapters/section02.tex；refer/papers/TMD_Soft_Function_LaMET_latex/chapters/section02.tex；refer/papers/LaMET计算TMD软函数_latex/；refer/papers/Locally_Smeared_OPE_latex/；refer/papers/Quasi_PDF_Gradient_Flow_latex/；../PyQCD/pyqcd/renorm/_tmd.py；../PyQCD/pyqcd/renorm/_tmdextract.py}

\section{联合极限的证据分级}

把所有操作都写进函数并不等于物理闭合。可把最终状态分成四级：
\begin{align}
 \texttt{interface}&:\text{几何、轴序和入口存在；}\\
 \texttt{structural}&:\text{有限维公式、投影、卷积骨架和边界检查闭合；}\\
 \texttt{scheme}&:\text{同几何 }Z/S/K/C\text{ 与转换核在窗口内闭合；}\\
 \texttt{ensemble}&:\text{真实系综、共同重采样、多个 }a,L,P_z,\tau,\ell\text{ 外推均通过}.
 \label{eq:deep-evidence-ladder}
\end{align}
当前 MyQCD 的 SymPy 推导主要到达第二级，并对若干 scheme 公式提供代数
代理；PyQCD 的入口提供从逐点算子到 Fourier 的实现接口。缺少同几何 soft
测量、rapidity regulator、胶子 singlet matching 或真实核子 disconnected
数据时，最终 \(f_1^{g[-,-]}\) 仍应标为 \unverified。这样的降级不是保守
措辞，而是由式\eqref{eq:deep-tmd-soft-subtraction}、
\eqref{eq:deep-cs-integrability} 和窗口式\eqref{eq:deep-nonlocal-flow-window}
共同强制出来的结论。
\src{refer/papers/准部分子分布单圈匹配_latex/；refer/papers/Quasi_Pseudo_PDF_Radyushkin_latex/；refer/papers/Complete_NP_Renorm_QuasiPDF_latex/；refer/papers/从准TMD波函数确定CS核_latex/；refer/papers/TMD_Soft_Function_LaMET_latex/；skills/lqcd-course/SKILL.md；myqcd/derivations.py；myqcd/pyqcd_theory_audit.py}

\appendix

\chapter{Grassmann 收缩、裸矩阵元与 OPE 证据闭环}
\label{sec:wick}

\section{Grassmann 积分与 Wick 逻辑}
费米子积分在固定规范背景下可写成
\begin{align}
 Z_F[\eta,\bar\eta;U]
 &=\int \mathcal D\bar\psi\,\mathcal D\psi\,
 \exp[-\bar\psi D[U]\psi+\bar\eta\psi+\bar\psi\eta] \nonumber\\
 &=\det D[U]\,
 \exp(\bar\eta D^{-1}[U]\eta),\qquad S[U]=D^{-1}[U].
\end{align}
因此 Wick 收缩只是 Grassmann 高斯积分的有限指标展开：
\begin{equation}
 \langle\psi_{i_1}\cdots\psi_{i_n}\bar\psi_{j_1}\cdots\bar\psi_{j_n}\rangle_U
 =\sum_{\pi\in S_n}(-1)^\pi \prod_{k=1}^n S_{i_k j_{\pi(k)}}[U].
\end{equation}
这也是 \texttt{seq\_peram}、Wick DSL、费米号和颜色闭合必须同置于一条链上的原因。
若离散核满足 \(\gamma_5 D\gamma_5=D^\dagger\)，则源/汇互换可由厄米共轭重建；
这正是蒸馏子空间里 \texttt{sink2src} 的数学前提。
\src{../PyQCD/pyqcd/contraction/_seqperam.py:20--59；../PyQCD/pyqcd/contraction/_autowick.py:39--445；refer/books/格点QCD导论_latex/chapters/sec17_correlators.tex}

\section{裸矩阵元、ratio 与三点/二点归一}
\label{sec:matrix}
OPE/三点分析中的公共代数核首先是方向组合
\begin{equation}
 O_{\rm comb}=-O_{30}-O_{31}+2O_{01},
\end{equation}
这正是 \texttt{ope\_combine} 与 \texttt{run\_bare\_matrix} 在三方向平均前的统一起点；
换言之，方向平均只改变统计权重，不改变该组合本身。三点/二点比值的谱极限为
\begin{equation}
 R(t,\tau)
 = \frac{C_3(t,\tau)}{C_2^f(t)}
 \sqrt{\frac{C_2^i(t-\tau)C_2^f(\tau)C_2^f(t)}
 {C_2^f(t-\tau)C_2^i(\tau)C_2^i(t)}}
 \longrightarrow \frac{\langle f|O|i\rangle}{2\sqrt{E_fE_i}}.
\end{equation}
\texttt{ratio\_3pt} 允许 1D/2D 源时轴和可选 link folding；\texttt{dis\_connect}、
\texttt{run\_disconnected\_ratio} 与 \texttt{run\_disconnected\_tmd\_ratio} 则把
\(C_3^{\rm disc}=C_3-C_2\langle O\rangle\) 写成时间滚移后的不相连构造。最终
\texttt{run\_bare\_matrix} 在 x/y/z 三个方向上分别生成 ratio，再做方向平均并拟合
\(c_0+c_1e^{-dE\Delta\tau}+c_1e^{-dE(t-\Delta\tau)}\)；此时返回的 \texttt{c0}
才是所选归一化约定下的裸矩阵元。
\src{../PyQCD/pyqcd/analysis/_analyse.py:361--820；../PyQCD/pyqcd/analysis/_ratio2pt.py:1--192；../PyQCD/pyqcd/analysis/_bare_matrix.py:1--160；../PyQCD/pyqcd/analysis/_tmd_ratio.py:1--180}

\section{胶子 OPE：直线 Wilson 线、场强插入和通道语义}
\label{sec:ope}
有限格距上的 legacy 直线 OPE 可写成
\begin{equation}
 O_{\mu\nu;\mu_2\nu_2}(\Delta z,t)
 = \sum_{\vec x}\operatorname{Tr}\!\Big[
 F_{\mu\nu}(x+\Delta z\hat z)\,W^\dagger(x+\Delta z,x)\,
 F^{(\sim)}_{\mu_2\nu_2}(x)\,W(x,x+\Delta z)
 \Big],
\end{equation}
其中 \(F^{(\sim)}\) 由 \texttt{second\_insert=F/Ftilde} 决定，\texttt{direction=+1/-1}
决定直线 Wilson 线是 \(+z\) 还是 \(-z\) 的显式构造。
\texttt{gluon\_ope\_operator\_z0} 的输入是
\texttt{gauge, mu, nu, z\_dir, delta\_z, Nt, Nx} 以及可选的
\texttt{mu2, nu2, direction, second\_insert, field\_strength\_cache, flow\_time,
sum\_kind, normalization, output\_projection, field\_projection}；输出是
\((\Delta z,N_t)\) 数组。这里 \texttt{sum\_kind=full}、
\texttt{normalization=bare\_spatial\_sum}、\texttt{field\_projection=legacy\_untraced}
是一整套 legacy 语义，不是可以任意拆开的装饰参数。
\texttt{gluon\_ope\_channel} 只是 \texttt{OPEChannelSpec} 的全显式包装器；
\texttt{gluon\_ff\_operator\_z0} 是固定规范下不含 Wilson 线的对照通道；
\texttt{get\_ope\_lorentz\_pairs} 则给出 unpol/helicity/gauge-fix 的 Lorentz 组合表。
重要边界是：finite-\(a\) 的 \texttt{legacy\_untraced} 直线 OPE 只对应 straight bilocal，
而 finite-staple TMD 还包含横向分离、软因子和 rapidity 方案，二者不能混写成同一物理对象。
\src{../PyQCD/pyqcd/operator/_gluon_ope.py:617--848；../PyQCD/docs/格点QCD中的OPE算符.tex；refer/papers/Locally_Smeared_OPE_latex/chapters/section02.tex；refer/papers/Quasi_Pseudo_PDF_Radyushkin_latex/chapters/section02.tex}

\section{代码—物理对象—输入/中间量/输出映射}
\label{sec:code-map}
\begin{longtable}{@{}P{0.94\textwidth}@{}}
\caption{PyQCD 关键接口的物理对象、输入轴序、中间量与输出门槛}\label{tab:code-map}\\
\toprule
\textbf{接口—物理对象—输入/中间量/输出}\\
\midrule
\endfirsthead
\caption[]{PyQCD 关键接口的物理对象、输入轴序、中间量与输出（续）}\\
\toprule
\textbf{接口—物理对象—输入/中间量/输出}\\
\midrule
\endhead
\bottomrule
\endfoot
\texttt{hyp\_smear}: 输入 \(U\in\mathbb C^{N_t\times N_z\times N_y\times N_x\times4\times3\times3}\) 与 \((\alpha_1,\alpha_2,\alpha_3)\)；中间量 \(V_3\to V_2\to V_1\) 的三级 staple；输出同形 SU(3) 链接，且不改写原数组。\src{../PyQCD/pyqcd/smear/_hyp.py:1--116}\\
\texttt{stout\_smear}: 输入 gauge 同上、\texttt{nstep} 与 \texttt{rho}；中间量 \(\Omega\)、\(Q\) 与 \(\exp(iQ)\)；输出空间三向涂抹后的链接，时间方向保持不变。\src{../PyQCD/pyqcd/smear/_stout.py:1--170}\\
\texttt{apply\_momentum\_smearing}: 输入 \texttt{eigvecs}（\((N_{\rm ev},N_z,N_y,N_x,N_c)\) 或展平），\texttt{momentum=(p\_z,p\_y,p\_x)} 与可选 \texttt{lattice\_shape}；中间量 \(\exp[-i2\pi\,\boldsymbol p\cdot\boldsymbol x/\boldsymbol L]\)；输出同形蒸馏本征矢相位乘法。\src{../PyQCD/pyqcd/vertex/_vertex.py:643--760}\\
\texttt{Mom\_VdV\_sink\_t}: 输入 \texttt{phase\_exp} 与 \texttt{eigvecs}；中间量把空间与颜色展平后做一次 \texttt{einsum}；输出 \(V_{mn}(p)\) 的动量分辨顶点张量。\src{../PyQCD/pyqcd/vertex/_vertex.py:134--181}\\
\texttt{Mom\_VVV\_sink\_t}: 输入 \texttt{phase\_exp} 与 \texttt{eigvecs}；中间量 \( \epsilon_{abc}\phi^a\phi^b\phi^c \) 的六项反对称收缩；输出动量分辨的 \(V_{mnl}(p)\) 重子 singlet 顶点。\src{../PyQCD/pyqcd/vertex/_vertex.py:182--240}\\
\texttt{VdV\_sink\_t\_link}: 输入 \texttt{eigvecs, phase\_exp, link\_dir, link\_max, gauge\_link}；中间量沿空间方向串联 Wilson transport；输出无链接/保守流/空间链接三种模式下的 \(V_{mn}(p,\Delta x)\)。\src{../PyQCD/pyqcd/vertex/_vertex.py:241--382}\\
\texttt{sink2src}: 输入 sink 顶点张量与 \texttt{dtype} \(\in\{\texttt{VdV},\texttt{VVV}\}\)；中间量只是厄米共轭和指标交换；输出 source 顶点，用于 \texttt{seq\_peram} 与源/汇转换。\src{../PyQCD/pyqcd/vertex/_vertex.py:383--420}\\
\texttt{Jackknife}/\texttt{Bootstrap}/\texttt{resample}: 输入配置轴、\texttt{cov\_axes}、\texttt{only\_sample} 与样本数；中间量是 leave-one-out 或有放回索引；输出样本、均值、误差和可选协方差，直接服务谱学与矩阵元拟合。\src{../PyQCD/pyqcd/analysis/_analyse.py:1--240；../PyQCD/pyqcd/analysis/_disconnected.py:146--220}\\
\texttt{meff}/\texttt{solve\_gevp}: 输入两点相关函数或相关矩阵 \(\,C(t)\,\)；中间量是 log/cosh 有效质量或 Hermitian 化后的广义特征值问题 \(C(t)v=\lambda C(t_0)v\)；输出 \(m_{\rm eff}\)、\(\lambda_n\) 与本征矢。\src{../PyQCD/pyqcd/analysis/_analyse.py:241--329；../PyQCD/pyqcd/analysis/_analyse.py:599--700}\\
\texttt{ratio\_3pt}: 输入 3pt/2pt 样本、\texttt{t\_sep} 与时间轴约定；中间量是 \(C_2^I,C_2^F\) 的开方运动学因子；输出 \(R(\tau)\) 的样本、均值与误差，支持 1D/2D 源时轴。\src{../PyQCD/pyqcd/analysis/_analyse.py:330--450}\\
\texttt{dis\_connect}/\texttt{run\_disconnected\_ratio}/\texttt{run\_disconnected\_tmd\_ratio}: 输入 2pt 与 bubble/OPE 样本；中间量是按源时滚移对齐后的 \(C_3^{\rm disc}=C_3-C_2\langle O\rangle\)；输出不相连比值、\texttt{c0/c1/dE/chi2} 和 plateau 状态。\src{../PyQCD/pyqcd/analysis/_analyse.py:744--820；../PyQCD/pyqcd/analysis/_disconnected.py:447--560；../PyQCD/pyqcd/analysis/_tmd_ratio.py:78--260}\\
\texttt{ope\_combine}/\texttt{run\_bare\_matrix}: 输入三个方向的 OPE 切片与 2pt；中间量先做 \(-O_{30}-O_{31}+2O_{01}\)，再做方向平均与多窗口拟合；输出 ratio、裸矩阵元 \texttt{c0}、\texttt{fit\_report} 与画图文件。\src{../PyQCD/pyqcd/analysis/_ratio2pt.py:1--192；../PyQCD/pyqcd/analysis/_bare_matrix.py:1--160}\\
\texttt{gluon\_ope\_operator\_z0}/\texttt{gluon\_ope\_channel}/\texttt{gluon\_ff\_operator\_z0}/\texttt{get\_ope\_lorentz\_pairs}: 输入 gauge、Lorentz 对、\texttt{z\_dir}、\texttt{delta\_z}、\texttt{direction} 与通道元数据；中间量是 \(F\)、\(F^{\sim}\) 与直线 Wilson 线；输出 \((\Delta z,N_t)\) 直线 OPE、固定规范 FF 对照和明确的 Lorentz 组合表。\src{../PyQCD/pyqcd/operator/_gluon_ope.py:617--848}\\
\texttt{hR\_PDF}/\texttt{C\_gluon\_ratio}: 输入 \(x\) 网格、\texttt{Pz}、\texttt{hR\_tilde} 与 \(\lambda_s\)；中间量是 \(\xi=x/y\) 的三分区核 \(g_1,g_2,g_3\) 与 Si 项；输出匹配到光锥 PDF 的 \(h_R\) 或 ratio 方案核，且显式拒绝精确 \(x=0\) 节点。\src{../PyQCD/pyqcd/renorm/_matching.py:1--180}\\
\texttt{quasi\_pdf\_gluon}: 输入 \(h(z)\)、\(z\_grid\) 与 \(\texttt{pz\_gev}\)；中间量是非负 \(z\) 半轴上的 \(\sin(xP_zz)\) 积分；输出 collinear 胶子准 PDF 的反对称 Fourier 通道。\src{../PyQCD/pyqcd/renorm/_tmdextract.py:1--165}\\
\texttt{quasi\_tmd\_pdf}: 输入 \(h_R(z,b_\perp)\)、\texttt{z\_grid}、\texttt{b\_perp} 与 \texttt{pz\_gev}；中间量是 Hermitian 半轴延拓与梯形积分；输出 \(x\,\widetilde g(x,b_\perp,P_z)\)，并保留 \(b_\perp\) 依赖。\src{../PyQCD/pyqcd/renorm/_tmdextract.py:1--120}\\
\texttt{gluon\_tmd\_operator}: 输入 \(U,z,b_\perp,z\_dir,b\_dir,L\)；中间量是 staple Wilson 线和三个场强 \(F_{ti},F_{tj},F_{ij}\)；输出 \(M^{ti;ti}+M^{tj;tj}-2M^{ij;ij}\) 的逐点 TMD 组合。\src{../PyQCD/pyqcd/renorm/_tmd.py:233--320}\\
\texttt{soft\_function\_intrinsic}: 输入平方比值矩阵元 \(R^2(b_\perp)\)；中间量是 \(R^2(b_\perp)/R^2(0)\) 的归一框架；输出内禀软函数 \(S_I\)，但不等于独立真空 Wilson loop 的完整软因子。\src{../PyQCD/pyqcd/renorm/_tmdextract.py:220--230}\\
\texttt{cs\_kernel\_two\_momentum}: 输入两个动量的 \(c_0\) 或 \(h_R\) 数据和 \(\texttt{z\_ref}\)；中间量是 \(\log\) 比值与 \(\log(P_{z1}/P_{z2})\)；输出 Collins--Soper 核 \(K(b_\perp)\) 及可选截断。\src{../PyQCD/pyqcd/renorm/_tmdextract.py:187--220}\\
\texttt{tmd\_matching\_hybrid}: 输入 \(x/y\) 网格、\(b_\perp\)、\(\mu\)、\(P_z\)、软因子、CS 核和 \(x\widetilde g\)；中间量是离散 \(Z_{ij}\)、\(S_I\) 和快度演化因子；输出混合方案下的匹配 TMD。\src{../PyQCD/pyqcd/renorm/_tmdextract.py:230--304}\\
\texttt{wilson\_flow}: 输入 \(U,\tau,\epsilon\) 或步数；中间量是六个 staple 的 \(\Omega_\mu\) 与 su(3) 投影的 \(Z_\mu\)；输出流场 \(V_\mu(t)\)，保持 \(SU(3)\) 幺正性。\src{../PyQCD/pyqcd/renorm/_gradient_flow.py:126--267}\\
\texttt{sftx\_gluon\_matching\_coeff}: 输入 \(\mu\) 以及 \(t_{\rm GeV^{-2}}\) 或 \((\tau,a_{\rm fm})\)；中间量是 \(\alpha_s/(4\pi)\) 与 \(2b_0[\log(2\mu^2t)+\gamma_E]\)；输出一圈 SFTX 系数。\src{../PyQCD/pyqcd/renorm/_tmdextract.py:318--381}\\
\end{longtable}

\section{端到端单列算法}
\begin{longtable}{@{}P{0.94\textwidth}@{}}
\caption{从规范组态到 PDF/TMD 的完整单列算法链}\label{tab:end-to-end}\\
\toprule
\textbf{一步只说一件事：输入、物理判断、失败门槛、输出}\\
\midrule
\endfirsthead
\caption[]{从规范组态到 PDF/TMD 的完整单列算法链（续）}\\
\toprule
\textbf{一步只说一件事：输入、物理判断、失败门槛、输出}\\
\midrule
\endhead
\bottomrule
\endfoot
\(1.\) 读入规范场 \(U\) 并先做轴序、边界、颜色维与 dtype 校验；若链接不是 \((N_t,N_z,N_y,N_x,4,3,3)\) 或方向约定不匹配，则立即停止，而不是猜测重排规则；输出是可继续处理的 gauge 组态。\src{../PyQCD/pyqcd/lattice/_constants.py；../PyQCD/pyqcd/renorm/_gradient_flow.py；../PyQCD/pyqcd/operator/_gluon_ope.py}\\
\(2.\) 选择 UV 预处理：HYP 采用三级排除方向 stapled average，Stout 采用 \(\Omega\to Q\to \exp(iQ)\) 的解析指数；若研究目标是 TMD 或 OPE，必须记录是否涂抹、涂抹步数和是否保留时间方向。\src{../PyQCD/pyqcd/smear/_hyp.py；../PyQCD/pyqcd/smear/_stout.py}\\
\(3.\) 若需要蒸馏或动量投影，则先对 Laplacian 本征矢做正交归一，再施加 \(\exp[-i2\pi\boldsymbol p\cdot\boldsymbol x/\boldsymbol L]\)；失败门槛是本征矢形状、复精度或 \texttt{lattice\_shape} 不一致。\src{../PyQCD/pyqcd/vertex/_vertex.py:643--760}\\
\(4.\) 构造 \texttt{Mom\_VdV\_sink\_t} 与 \texttt{Mom\_VVV\_sink\_t}，把空间、颜色和动量压成 \(V_{mn}(p)\) 与 \(V_{mnl}(p)\)；若颜色 singlet 或相位轴序不闭合，则不能进入收缩阶段。\src{../PyQCD/pyqcd/vertex/_vertex.py:134--240}\\
\(5.\) 若要做链路输运或保守流，调用 \texttt{VdV\_sink\_t\_link} 与 \texttt{sink2src}，显式区分 sink/source；这里的关键物理图像是厄米共轭和 Wilson 输运，而不是简单复制数组。\src{../PyQCD/pyqcd/vertex/_vertex.py:241--420}\\
\(6.\) 解析算符 DSL，检查味平衡、颜色闭合与费米号；若某个算符的夸克数与反夸克数不平衡，则该图在 Grassmann 代数下直接为零，不应继续进入数值管线。\src{../PyQCD/pyqcd/contraction/_autowick.py:39--445；../PyQCD/pyqcd/contraction/_dynamic.py:262--647}\\
\(7.\) 组装两点、三点与 bubble 结构，必要时用 \texttt{dis\_connect} 做不相连构造；若研究的是 disconnected TMD，则要用 \texttt{run\_disconnected\_tmd\_ratio} 把 \(C_3^{\rm disc}\) 与 OPE 背景一起处理。\src{../PyQCD/pyqcd/analysis/_analyse.py:744--820；../PyQCD/pyqcd/analysis/_disconnected.py:447--560；../PyQCD/pyqcd/analysis/_tmd_ratio.py:78--260}\\
\(8.\) 对相关函数做源时平移平均和周期/反周期边界修正；失败门槛是源轴与汇轴长度不一致或边界符号未明确记录。\src{../PyQCD/pyqcd/analysis/_analyse.py:720--820}\\
\(9.\) 用 \texttt{Jackknife}/\texttt{Bootstrap}/\texttt{resample} 形成样本分布，再从中算协方差、均值和误差；如果 \(N_{\rm conf}<2\)，delete-one jackknife 没有可辨识协方差，只能返回边界态。\src{../PyQCD/pyqcd/analysis/_analyse.py:91--240；../PyQCD/pyqcd/analysis/_disconnected.py:146--220}\\
\(10.\) 从两点函数抽取 \texttt{meff}，从相关矩阵求 \texttt{solve\_gevp}；若 \(C(t_0)\) 不是 Hermitian 正定，或者有效秩太低，则能谱提取只能标为不可辨识。\src{../PyQCD/pyqcd/analysis/_analyse.py:241--329；../PyQCD/pyqcd/analysis/_analyse.py:599--700}\\
\(11.\) 进入比值与裸矩阵元阶段：先用 \texttt{ratio\_3pt} 计算 \(R(\tau)\)，再用 \texttt{ope\_combine} 和 \texttt{run\_bare\_matrix} 把方向组合、方向平均和多窗口拟合串起来；输出的 \(c_0\) 才是所选归一化下的裸矩阵元。\src{../PyQCD/pyqcd/analysis/_ratio2pt.py:1--192；../PyQCD/pyqcd/analysis/_bare_matrix.py:1--160}\\
\(12.\) 若研究胶子 OPE，调用 \texttt{gluon\_ope\_operator\_z0} 或 \texttt{gluon\_ff\_operator\_z0}，把 \(F\)、\(F^{\sim}\) 和直线 Wilson 线收成 \((\Delta z,N_t)\)；若通道表不明确，先用 \texttt{get\_ope\_lorentz\_pairs} 固定 \((\mu,\nu,\mu_2,\nu_2)\) 再做数值。\src{../PyQCD/pyqcd/operator/_gluon_ope.py:617--848}\\
\(13.\) 若走准 PDF / OPE 匹配，先做 \texttt{hR\_PDF} 或 \texttt{C\_gluon\_ratio} 的离散核逆；这里必须避开精确 \(x=0\) 节点，并保留 \(\xi=x/y\) 的三分区结构和 Si 项，不可把主值点硬塞成任意小数。\src{../PyQCD/pyqcd/renorm/_matching.py:1--180}\\
\(14.\) 若走 Fourier 变换，collinear 胶子准 PDF 只能用 \texttt{quasi\_pdf\_gluon} 的 sin 通道；准 TMD 则必须用 \texttt{quasi\_tmd\_pdf} 的 Hermitian 半轴 cos/sin 通道，二者不得互换。\src{../PyQCD/pyqcd/renorm/_tmdextract.py:1--165}\\
\(15.\) 若需要 UV 平滑和 SFTX，则先 \texttt{wilson\_flow} 到流时间 \(t\)，再用 \texttt{sftx\_gluon\_matching\_coeff} 把 \(\mathcal O_{\rm flow}(t)\) 送回 \(\overline{\rm MS}\)；失败门槛是 \(a^2\ll t\ll\Lambda_{\rm QCD}^{-2}\) 窗口不存在。\src{../PyQCD/pyqcd/renorm/_gradient_flow.py:126--267；../PyQCD/pyqcd/renorm/_tmdextract.py:318--381}\\
\(16.\) 若研究 TMD，则先调用 \texttt{gluon\_tmd\_operator} 构造 staple bilocal，再用 \texttt{soft\_function\_intrinsic}、\texttt{cs\_kernel\_two\_momentum} 与 \texttt{tmd\_matching\_hybrid} 处理 soft/rapidity 演化；没有同几何 soft factor 时，最终 TMD 只能停在未验证边界。\src{../PyQCD/pyqcd/renorm/_tmd.py:233--320；../PyQCD/pyqcd/renorm/_tmdextract.py:187--304}\\
\(17.\) 所有中间量必须连同 \texttt{status}、\texttt{window}、\texttt{source\_inventory}、\texttt{fit\_report} 和 \(P_z,\mu,\tau,b_\perp\) 的元数据一起落盘；如果只剩接口、没有真实系综或完整卷积，就必须把结果写成结构性或未验证，而不是物理定论。\src{data/pyqcd_theory_audit/theory_audit.json；myqcd/pyqcd_theory_audit.py}\\
\end{longtable}

\iffalse
\section{两点和三点函数的谱分解}
设传递矩阵存在、Euclidean 反射正性成立，且算符能够与能量本征态重叠。
采用相对论归一化
\(\langle n,\boldsymbol p|n,\boldsymbol p\rangle=2E_nV\)，两点函数的无限时间近似为
\begin{align}
 C_2(t,\boldsymbol p)
 &=\sum_n\frac{|Z_n(\boldsymbol p)|^2}{2E_n(\boldsymbol p)}e^{-E_n t},
 \\
 C_2^{\rm periodic}(t)
 &=\sum_n\frac{|Z_n|^2}{2E_n}
 \left[e^{-E_nt}+e^{-E_n(N_t-t)}\right].
\end{align}
有限 \(N_t\)、反周期费米边界和算符的反向传播贡献会改变拟合模型，
不能把第一式在所有 \(t\) 上硬套。三点函数在源—插入—汇时间顺序下为
\begin{equation}
 C_3(t_{\rm sep},\tau)=
 \sum_{m,n}\frac{Z_m^f Z_n^{i*}}{4E_m^fE_n^i}
 e^{-E_m^f(t_{\rm sep}-\tau)}e^{-E_n^i\tau}
 \langle m,f|\mathcal O|n,i\rangle.
 \label{eq:c3-spectrum-clean}
\end{equation}
当 \(\tau\)、\(t_{\rm sep}-\tau\) 都大于激发态时间尺度时，基态项支配；
否则每一个省略项都带有 \(e^{-\Delta E\tau}\) 或
\(e^{-\Delta E(t_{\rm sep}-\tau)}\) 的系统误差。

对于胶子 OPE 的不相连因子化分析，工作树中的中间量是
\begin{equation}
 C_3^{\rm fact}(t,\tau,z)=C_2(t)\,O(\tau,z),\qquad
 C_3^{\rm disc}=C_3^{\rm fact}-C_2\,\langle O\rangle,\qquad
 R(z)=\left\langle\frac{C_3^{\rm disc}}{C_2}\right\rangle_{t_{\rm src}}.
 \label{eq:disconnected-ratio-clean}
\end{equation}
式中的真空扣除是 disconnected 矩阵元的定义组成部分；它既不是普通
统计均值，也不能在重采样之后随意改成两个独立均值的商。代码将相对时间
通过 \texttt{roll} 对齐源点，因此周期边界是该操作正确性的必要假设。
\src{../PyQCD/pyqcd/analysis/_analyse.py:361--520；../PyQCD/pyqcd/analysis/_ratio2pt.py:80--194；../PyQCD/pyqcd/analysis/_disconnected.py:1--13；../PyQCD/docs/格点QCD中的3pt构造.tex}

\section{本章的定理、假设与边界}
\begin{table}[htbp]
\centering
\caption{Wick 与关联函数理论链的公理—结论对应}
\small
\begin{tabularx}{0.96\textwidth}{@{}P{2.8cm}Y Y@{}}
\toprule
公理/假设 & 数学结论 & 在代码中的验证边界\\
\midrule
Grassmann 反对易、高斯费米测度、\(D\) 可逆 & Wick 展开及费米置换号 & 可由有限索引枚举；不证明真实 \(D^{-1}\) 已求精确。\\
局域规范协变性、颜色 singlet 算符 & 颜色指标闭合、关联函数规范不变 & VVV/VDV 标签结构可审计；不验证输入 eigenvector 与 gauge 同组态。\\
反射正性和传递矩阵 & Euclidean 能谱非负、长时间基态支配 & 需要正确作用量和边界；代数脚本不证明反射正性。\\
\(\gamma_5D\gamma_5=D^\dagger\) & 式 \eqref{eq:seq-peram-clean} 的反向传播子关系 & gamma 矩阵代数通过；具体费米子核仍需实测。\\
平移不变性与周期/反周期边界 & 源时平均、动量投影、\texttt{roll} 对齐 & 仅检查公式；边界符号和数据布局需用真实配置回归。\\
\bottomrule
\end{tabularx}
\end{table}

\fi

\chapter{谱学提取、重采样与相关拟合}
\label{sec:spectrum}

\section{动量、能量和有效质量}
代码中的动量列表是整数单位 \(\boldsymbol n=(n_z,n_y,n_x)\)，在空间长度
\(N_xa\)（立方空间）下
\begin{equation}
 |\boldsymbol p|=\frac{2\pi}{N_xa}|\boldsymbol n|,
 \qquad
 E(\boldsymbol p)=\sqrt{M_0^2+|\boldsymbol p|^2}.
 \label{eq:mom2gev-clean}
\end{equation}
\texttt{Mom2GeV} 的一个容易误读之处是：若 \texttt{Mom} 是标量，它被解释为
\(|\boldsymbol n|^2\)，而若是 \texttt{[pz,py,px]} 列表，函数先求平方和；
因此调用者不能把标量 \(n\) 当作向量模长而再次平方。量纲检查要求
\(a\) 从 fm 转成 \(\mathrm{GeV}^{-1}\)，使用
\(\hbar c=0.1973269804\,\mathrm{GeV\,fm}\)。

单指数关联函数给出对数有效质量
\begin{equation}
 am_{\rm eff}^{\log}(t)=\log\frac{C(t)}{C(t+1)},\qquad
 m_{\rm eff}[\mathrm{GeV}]
 =\frac{\hbar c}{a[\mathrm{fm}]}\log\frac{C(t)}{C(t+1)}.
\end{equation}
周期边界下两向传播的有效质量满足
\begin{equation}
 m_{\rm eff}^{\cosh}(t)=
 \operatorname{arccosh}\frac{C(t+2)+C(t)}{2C(t+1)},\qquad
 m[\mathrm{GeV}]=\frac{\hbar c}{a[\mathrm{fm}]}\,m_{\rm eff}^{\cosh}.
 \label{eq:cosh-meff-clean}
\end{equation}
PyQCD 在复数或非浮点输入上先取绝对值；在噪声使 arccosh 自变量小于
1 时，把结果置零。这是数值保护和信号筛选，不是物理上“质量为零”的
证明；报告结果时必须保留该掩码/截断规则。

对于算符基底 \(O_i\)，广义特征值问题为
\begin{equation}
 C(t)v_n(t,t_0)=\lambda_n(t,t_0)C(t_0)v_n(t,t_0),\qquad
 \lambda_n(t,t_0)\simeq e^{-E_n(t-t_0)}
 \left[1+O(e^{-\Delta E_n(t-t_0)})\right].
 \label{eq:gevp-clean}
\end{equation}
需要 \(C(t_0)\) Hermitian 正定、算符重叠矩阵在目标子空间满秩。代码先做
\(C\leftarrow(C+C^\dagger)/2\)，再调用广义 Hermitian 特征值求解器，并按
\(t<t_0\) 与 \(t\ge t_0\) 使用不同排序；这解决的是复数舍入和排序约定，
不替代对 \(t_0\)、基底维数和激发态系统误差的扫描。
\src{../PyQCD/pyqcd/analysis/_analyse.py:30--84,253--354,599--664；refer/books/INTRODUCTION_TO_LATTICE_QCD_latex/chapters/sec17_correlators.tex}

\section{Jackknife、bootstrap 与协方差}
\label{sec:statistics}
设有 \(N\) 个配置观测 \(x_i\)，均值为 \(\bar x=N^{-1}\sum_i x_i\)。
删除第 \(i\) 个配置的 Jackknife 样本是
\begin{equation}
 x_{(-i)}=\frac{\sum_jx_j-x_i}{N-1},\qquad
 \widehat{\operatorname{Cov}}_J(\bar x)
 =\frac{N-1}{N}\sum_{i=1}^{N}
 (x_{(-i)}-\bar x)(x_{(-i)}-\bar x)^T.
 \label{eq:jack-cov-clean}
\end{equation}
如果 \texttt{cov\_axes} 指定多个时间/空间点，代码先把非协方差轴保留，再把
协方差轴展平，执行 \(r_{ni}r_{nj}\) 的 Einstein 求和，最后恢复两份
协方差轴。复杂观测量可用实部、虚部拼成联合向量，但必须在统计说明中
明确该选择。

Bootstrap 以有放回索引 \(I_b=(i_1,\ldots,i_M)\) 生成
\begin{equation}
 x^{(b)}=\frac1M\sum_{k=1}^{M}x_{i_k},\qquad
 \widehat{\operatorname{Cov}}_B
 =\frac1{N_B}\sum_b(x^{(b)}-\bar x_B)(x^{(b)}-\bar x_B)^T.
\end{equation}
PyQCD 的默认 \(M=\max(N-5,1)\)、\(N_B=4N\)，并把全配置均值放在
第一个样本位置；随机种子、配置自相关和阻塞策略应由上层管线固定，不能
把默认参数误认为普适最优选择。删除一配置 Jackknife 在 \(N<2\) 时没有
可辨识协方差，当前实现会留下 NaN/状态标记而不强行拟合。

源时平移平均利用
\begin{equation}
 C(\Delta t)=\frac1{N_t}\sum_{t_{\rm src}}
 C(t_{\rm src},t_{\rm src}+\Delta t),\qquad
 \Delta t=(t_{\rm sink}-t_{\rm src})\bmod N_t.
\end{equation}
反周期边界在跨过时间边界时引入负号；三点函数还要根据 \(t_{\rm sep}\)
确定尾部符号。\texttt{loop\_tsrc} 采用 CPU 临时数组实现这一步，若输入来自 GPU，
返回时再送回当前后端。这个数据搬运行为应在性能报告中单独记录。

\section{相关拟合和可辨识性}
给定数据向量 \(y\)、模型 \(f(\theta)\) 和协方差 \(C\)，相关拟合的核心
目标是
\begin{equation}
 \chi^2(\theta)=[y-f(\theta)]^\dagger C^{-1}[y-f(\theta)].
 \label{eq:correlated-chi2-clean}
\end{equation}
当 \(C\) 由有限 bootstrap/Jackknife 样本估计而秩不足时，直接求逆会放大
噪声。PyQCD 的拟合器先检查 Hermitian 半正定性，再对相关矩阵做 SVD
调节并报告样本秩、有效秩和条件数。若数据自由度不超过参数数目，或有效
秩不足，拟合状态为不可辨识；有 prior 时可以继续，但状态必须写成
\texttt{prior\_constrained}，不能声称参数由数据独立决定。

\begin{table}[htbp]
\centering
\caption{谱学与统计分析的单列算法}
\small
\begin{tabular}{@{}P{0.94\textwidth}@{}}
\toprule
\(\text{Input: }C_{\rm raw}[N_{\rm conf},\ldots,t],\;a[\mathrm{fm}],\;\text{边界条件},\;\text{拟合窗口}\)\\
\(\text{Validate: }N_{\rm conf}\ge2\text{（Jackknife），时间轴和源汇轴不重复且长度相容}\)\\
\(\text{Average: }C(t_{\rm src},t_{\rm sink})\to C(\Delta t)\text{，按周期/反周期边界处理符号}\)\\
\(\text{Resample: }\{C_i\}\to\{C_{(-i)}\}\text{ 或 }\{C^{(b)}\}\text{，保留配置轴在已知位置}\)\\
\(\text{Covariance: }r_{ni}=x_{ni}-\bar x_i,\;C_{ij}\gets\sum_n r_{ni}r_{nj}\text{ 并使用相应归一化}\)\\
\(\text{Spectrum option 1: }m_{\rm eff}\gets\log[C(t)/C(t+1)]\cdot\hbar c/a\)\\
\(\text{Spectrum option 2: }m_{\rm eff}\gets\operatorname{arccosh}[(C(t+2)+C(t))/(2C(t+1))]\cdot\hbar c/a\)\\
\(\text{Spectrum option 3: }C(t)v=\lambda(t,t_0)C(t_0)v\text{，检查 }C(t_0)>0\text{ 并排序 }\lambda\)\\
\(\text{Fit preparation: }y\gets \text{选定 }(t,\tau,z,p)\text{ 点，}\;C\gets \text{由同一重采样方案得到的协方差}\)\\
\(\text{Rank gate: }\operatorname{rank}_{\rm sample},\operatorname{rank}_{\rm SVD},N_{\rm data},N_{\rm par}\text{ 均记录}\)\\
\(\text{Fit: }\min_\theta\chi^2(\theta)\text{，必要时输入 prior，但标记 prior 对不可辨识方向的约束}\)\\
\(\text{Propagate: }\text{逐样本重复 }m_{\rm eff}/\text{GEVP}/\text{fit，最终从样本分布取均值和误差}\)\\
\(\text{Scan: }t_0,\;t_{\rm min},t_{\rm max},N_{\rm ev},P_z,\text{拟合模型和 SVD cut}\)\\
\(\text{Output: }\text{谱量、误差、协方差、条件数、有效秩、fit status、窗口和全部约定元数据}\)\\
\bottomrule
\end{tabular}
\end{table}
\src{../PyQCD/pyqcd/analysis/_analyse.py:91--162,169--246,527--664；../PyQCD/pyqcd/analysis/_fitter.py:59--256；../PyQCD/docs/格点QCD中的关联函数.tex}

\section{本章的 SymPy 复现}
新增 \codeid{jackknife_covariance} 审计以二维符号观测验证式
\eqref{eq:jack-cov-clean}；核心 \codeid{three_point_ratio_spectral_limit}
审计以 \(C_2=|Z|^2e^{-Et}/(2E)\) 以及上文三点函数的谱分解基态项精确消去
overlap 和指数，得到
\begin{equation}
 R\longrightarrow\frac{\langle f|O|i\rangle}{2\sqrt{E_fE_i}}.
\end{equation}
这两个检查确证的是归一化和有限维代数；它们不测量真实 \(C_2,C_3\)，也
不排除有限 \(t_{\rm sep}\) 的激发态污染。

\chapter{非局域算符重整化：\(Z_R\)、RI/MOM 与混合方案}
\label{sec:renorm}

\section{为什么 Wilson 线带来新的发散}
准 PDF 或伪 PDF 的基本夸克双线性可以写成
\begin{equation}
 \mathcal O_\Gamma(z)=\bar\psi(z\hat n)\Gamma
 W(z\hat n,0)\psi(0),\qquad
 W(z\hat n,0)=\mathcal P\exp\left(\ii g\int_0^z\dd s\,A_n(s\hat n)\right).
 \label{eq:quasi-operator-clean}
\end{equation}
有限格距上，直线 Wilson 线的自能产生与长度成正比的线性发散，常用
结构记账为
\begin{equation}
 \mathcal O_B(z,a)\sim
 \exp[-\delta m(a)|z|]\,\mathcal O_{\rm finite}(z),
 \qquad \delta m(a)=O(a^{-1}),
 \label{eq:linear-divergence-clean}
\end{equation}
并伴有端点、对数和可能的 Dirac 结构混合。这里的 \(\delta m\) 不是
普通夸克质量；它是非局域线段的 scheme-dependent 自能参数。若按
\(h_R=h_B/Z_R\) 处理，\(Z_R\) 必须和同一 Wilson 线几何、方向、流时间、
格距以及 Dirac/颜色投影相匹配。

\section{\(Z_R\) 的参数化与数据输入}
对照项目的自重整化接口同时包含短距微扰窗口和长距非微扰节点。其核心
形式可以抽象为
\begin{align}
 \log Z_R(z,a,\mu)
 &=\frac{kz_G}{a\log(a\Lambda)}
 +\frac12\log\left(1+\frac{d}{\log(a\Lambda)}\right)^2
 +\frac{5N_c}{3b_0}
 \log\frac{\log[1/(a\Lambda)]}{\log(\mu/\Lambda)}
 \nonumber\\
 &\quad +(m_0+m_2a^2)z_G+f_1a+f_2a^2,
 \qquad z_G=\frac{z_{\rm fm}}{\hbar c}.
 \label{eq:zr-param-clean}
\end{align}
其中 \(b_0=11-2N_f/3\) 是代码使用的未除以 \(4\pi\) 的一圈系数约定。
\texttt{th\_ZR} 返回式 \eqref{eq:zr-param-clean} 的指数；\texttt{Z\_MS} 在短距部分采用
\begin{equation}
 Z_{\overline{\rm MS}}(z,\mu)=1+
 \frac{\alpha_sN_c}{4\pi}
 \left[\frac53\log\frac{z^2\mu^2}{4e^{-2\gamma_E}}+3\right].
 \label{eq:zms-clean}
\end{equation}
\texttt{th\_hB} 将 \(z\le z_1\) 的短距段写为 \(\log Z_{\overline{\rm MS}}\)
加质量项，将 \(z>z_1\) 的数据用 \(g_1,\ldots,g_{14}\) 节点描述，
并加入 \(a,a^2\) 离散化修正。该结构是可拟合的模型假设，而不是所有
非局域算符的普适定理；\(z_1\)、节点数和 \(\Lambda\) 的变化应列为系统
误差。
\src{../PyQCD/pyqcd/renorm/_zr.py:25--153,160--333；refer/papers/准光前关联函数的自重正化_latex/chapters/section02.tex；refer/papers/大动量有效理论的重正化_latex/chapters/section02.tex}

\section{RI/MOM 与 Dirac 混合}
RI/MOM 方案在外部夸克动量 \(p\) 上定义 amputated vertex
\(\Lambda_\Gamma(p,z)\)，用投影条件固定 \(Z_O\)：
\begin{align}
 Z_q^{-1}Z_O\,\frac1{N_{\rm proj}}
 \operatorname{Tr}\left[P_\Gamma\Lambda_\Gamma(p,z)\right]
 \bigg|_{p^2=\mu_R^2}&=1,\\
 Z_O(\mu_R)&=Z_q(\mu_R)\left[
 \frac1{N_{\rm proj}}\operatorname{Tr}(P\Lambda)\right]^{-1}.
\end{align}
这要求固定 gauge（通常 Landau gauge）、外部动量方向、投影矩阵、离散
动量窗口和 \(p_z\) 方向；它不是规范不变量的 hadron matrix element。
直线算符的一圈/有限格距混合可能把 \(\Gamma\) 与
\(\slash n\Gamma\) 联系起来。以代码核心推导使用的形式记账：
\begin{equation}
 \Gamma'=(1+\xi r\slash n)\Gamma(1+\xi r\slash n)
 =\Gamma+\xi r\{\slash n,\Gamma\}
 +\xi^2r^2\slash n\Gamma\slash n.
 \label{eq:gamma-mixing-clean}
\end{equation}
只有在 \(\{\slash n,\Gamma\}=0\) 或其他对称性条件下，混合项才会
简化；不能因某个数值通道看起来为零就删去。用户若需要 RI/MOM 到
\(\overline{\rm MS}\) 的转换，还必须输入转换矩阵和方案尺度，不能用
\(Z_R\) 的标量除法替代。

\section{hybrid 方案和拼接连续性}
对照实现的 hybrid 方案把短距比值重整化和长距自重整化拼接为
\begin{equation}
 h_R(z,P_z)=
 \begin{cases}
 h_B(z,P_z)/h_B(z,0),&z<z_s,\\[2pt]
 [h_B(z,P_z)/Z_R(z)]\,\eta_s,&z\ge z_s,
 \end{cases}
 \qquad
 \eta_s=\frac{Z_R(z_s)}{h_B(z_s,0)}.
 \label{eq:hybrid-clean}
\end{equation}
在 \(z_s\) 的两侧取同一节点时，
\begin{equation}
 \frac{h_B(z_s,P_z)}{h_B(z_s,0)}
 =\frac{h_B(z_s,P_z)}{Z_R(z_s)}\eta_s,
\end{equation}
所以代数上连续。连续性只说明归一化常数没有人为跳变，不证明两种
方案在有限 \(a\)、有限 \(P_z\) 下的所有短距/长距系统误差都消失。实现
还用 \(\lambda=z_G P_z\) 的尾部模型
\begin{equation}
 h_R(\lambda)\simeq l_1\lambda^{-a_1}e^{-\lambda/\lambda_0}
 \label{eq:lambda-tail-clean}
\end{equation}
外推有限 \(z\) 数据；该模型参数和替代尾部是 Fourier 系统误差的重要
来源。
\src{../PyQCD/pyqcd/renorm/_hybrid.py:1--153；../PyQCD/pyqcd/renorm/_zr.py:160--333；refer/papers/Hybrid_Renorm_Quasi_LightFront_latex/chapters/section02.tex；refer/papers/Wilson_Line_Renorm_Improved_QPD_latex/chapters/section02.tex}

\begin{table}[htbp]
\centering
\caption{非局域算符从裸数据到 hybrid 重整化的单列算法}
\small
\begin{tabular}{@{}P{0.94\textwidth}@{}}
\toprule
\(\text{Input: }h_B(z,P_z),\;h_B(z,0),\;z[\mathrm{fm}],\;a[\mathrm{fm}],\;\mu,\;\text{Wilson 线几何}\)\\
\(\text{Validate: }z\text{、动量、流时间、颜色投影和格距来自同一方案；短距/长距均有 }z_s\text{ 节点}\)\\
\(\text{Normalize: }z_G\gets z_{\rm fm}/(\hbar c),\quad \lambda\gets z_G P_z\)\\
\(\text{Short branch: }h_R^{\rm short}(z,P_z)\gets h_B(z,P_z)/h_B(z,0)\)\\
\(\text{Fit }Z_R:\ \log Z_R\gets\text{式 }\eqref{eq:zr-param-clean}\text{，联合格距/系综拟合并记录协方差}\)\\
\(\text{Set switch: }\eta_s\gets Z_R(z_s)/h_B(z_s,0)\)\\
\(\text{Long branch: }h_R^{\rm long}(z,P_z)\gets[h_B(z,P_z)/Z_R(z)]\eta_s\)\\
\(\text{Join: }h_R\gets h_R^{\rm short}\text{（}z<z_s\text{）或 }h_R^{\rm long}\text{（}z\ge z_s\text{）}\)\\
\(\text{Check: }h_R^{\rm short}(z_s)=h_R^{\rm long}(z_s)\text{；检查 }Z_R\ne0\text{、有限值和误差传播}\)\\
\(\text{Tail fit: }h_R(\lambda)\gets l_1\lambda^{-a_1}e^{-\lambda/\lambda_0}\text{（仅在受控大 }\lambda\text{ 窗口）}\)\\
\(\text{Fourier input: }\text{实部/虚部按 Hermitian 延拓规则分别处理，不将有限支撑当作无限支撑}\)\\
\(\text{Output: }h_R(z,P_z),Z_R(z),\eta_s,h_R(\lambda),\text{尾部参数、协方差和方案元数据}\)\\
\bottomrule
\end{tabular}
\end{table}

\section{本章的 SymPy 证据}
\texttt{bare\_zr\_hybrid\_fourier} 审计验证了 \(\eta_s\) 使式
\eqref{eq:hybrid-clean} 在拼接点连续，并验证代码 \(2/(2\pi)=1/\pi\)
的 Fourier 前因子。这不能验证 \(Z_R\) 拟合得到的参数、短距微扰窗口
或长距尾部；这些属于真实数据和方案系统误差层。

\chapter{Fourier、OPE 与 LaMET 匹配}
\label{sec:matching}

\section{Ioffe 时间和准分布}
空间分离 \(z\hat z\) 与核子大动量 \(P_z\) 组合成无量纲 Ioffe 时间
\begin{equation}
 \nu=zP_z,\qquad [z]=-1,\quad[P_z]=1.
 \label{eq:ioffe-clean}
\end{equation}
quasi-PDF 是等时空间相关函数的 Fourier 变换；在 \(P_z\to\infty\)
时，受目标幂次修正抑制，它才逼近光锥 PDF。对有限 \(P_z\)，一般写成
\begin{equation}
 \widetilde f(x,P_z,\mu)=
 \int_{-1}^{1}\frac{\dd y}{|y|}
 C\!\left(\frac{x}{y},\frac{\mu}{|y|P_z}\right)f(y,\mu)
 +O\!\left(\frac{\Lambda_{\rm QCD}^2}{P_z^2},
 \frac{M^2}{P_z^2},a^2P_z^2\right).
 \label{eq:lamet-factorization-clean}
\end{equation}
OPE 的局部短距展开和 LaMET 的大动量因子化有相似的卷积结构，
但它们的展开参数不同：OPE 是短距离 \(z\Lambda_{\rm QCD}\ll1\)，LaMET
是 \(P_z\gg\Lambda_{\rm QCD},M\)。不能把一圈 OPE 系数直接当作任意
有限 \(P_z\) 的完整 matching。

\section{两个 Fourier 通道必须分开}
工作树提供两个语义不同的 Fourier 接口。对直线胶子 collinear quasi-PDF，
代码使用反对称/sin 通道：
\begin{equation}
 \widetilde g(x,P_z)=\frac{2P_z}{x}
 \int_0^{z_{\max}}\dd z\,h(z,P_z)\sin(xP_zz),
 \qquad \widetilde g(0,P_z)\text{ 由实现保护为 }0.
 \label{eq:quasi-pdf-sin-clean}
\end{equation}
对准 TMD 接口，输入是非负 \(z\) 半轴的 \(h_R(z,b_T)\)，以
\(h_R(-z)=h_R(z)^*\) 延拓：
\begin{align}
 x\widetilde g(x,b_T,P_z)
 &=\frac1{\mathcal N}\int_{-z_{\max}}^{z_{\max}}
 \frac{\dd z}{2\pi P_z}e^{-\ii xzP_z}h_R(z,b_T,P_z),\\
 \mathcal N&=\frac{P_z^2}{P_t^2},\qquad
 P_t=\sqrt{M_N^2+P_z^2}\quad\text{（默认模型）}.
 \label{eq:quasi-tmd-fourier-clean}
\end{align}
把半轴写成实数积分，代码实际计算
\begin{equation}
 2\int_0^{z_{\max}}\dd z\,
 [\cos(xzP_z)\Re h_R(z)+\sin(xzP_z)\Im h_R(z)].
\end{equation}
因此 \texttt{quasi\_pdf\_gluon} 的 sin 变换和 \texttt{quasi\_tmd\_pdf} 的
cos/sin Hermitian 半轴变换不能互换。二者还分别使用 fm
\(\to\mathrm{GeV}^{-1}\) 转换、梯形权重、\(z=0\) 起始网格和有限
\(z_{\max}\) 截断。
\src{../PyQCD/pyqcd/renorm/_hybrid.py:115--153；../PyQCD/pyqcd/renorm/_tmdextract.py:30--165；refer/papers/Quasi_Pseudo_PDF_Radyushkin_latex/chapters/section02.tex；refer/papers/准部分子分布单圈匹配_latex/chapters/section02.tex}

\section{匹配核的分区与离散矩阵}
胶子一圈核按 \(\xi=x/y\) 的三个区域分段：
\begin{equation}
 \xi<0:\ g_{-}(\xi),\qquad
 0<\xi<1:\ g_{\rm mid}(\xi),\qquad
 \xi>1:\ g_{+}(\xi),
\end{equation}
并含由有限 \(\lambda_s\) 产生的 sine-integral 项
\(\operatorname{Si}[(1-\xi)|y|\lambda_s]\)。离散化后工作实现构造
\begin{align}
 Z_{ij}&=\delta_{ij}+\frac{\alpha_sN_c}{2\pi}
 \frac{x_i}{|x_i|}\Delta x
 \left[\frac{g(x_i/y_j)}{y_j}
 -\delta_{ij}y_i\sum_k\frac{g(x_i/y_k)}{y_k^2}\right],\\
 h_R^{\rm PDF}&=Z^{-1}h_R^{\rm quasi}.
 \label{eq:matching-matrix-clean}
\end{align}
第二项是离散求和规则/虚项的补偿结构，不应在改变网格后保留旧的
\(\Delta x\) 而不重算。代码的 \texttt{hR\_PDF} 明确拒绝含 \(x=0\) 节点，
因为当前离散实现没有经过验证的 principal-value 处方；这是重要的输入
边界，不能用 \(10^{-15}\) 的任意替换隐藏奇点。

\begin{table}[htbp]
\centering
\caption{从重整化坐标矩阵元到 quasi-PDF/PDF 的单列算法}
\small
\begin{tabular}{@{}P{0.94\textwidth}@{}}
\toprule
\(\text{Input: }h_B(z,P_z),h_B(z,0),Z_R(z),z_s,P_z,a,\mu,\text{有限 }z\text{ 网格}\)\\
\(\text{Renormalize: }h_B\to h_R\text{，按式 }\eqref{eq:hybrid-clean}\text{ 做短距比值/长距 }Z_R\text{ 拼接}\)\\
\(\text{Ioffe map: }z\to z_G=z/(\hbar c),\quad \lambda\gets z_GP_z,\quad \nu\gets z_GP_z\)\\
\(\text{Tail: }\lambda>\lambda_{\rm extra}\text{ 用尾部模型，扫描 }\lambda_{\rm extra}\text{ 和外推窗}\)\\
\(\text{Choose channel: }\text{collinear 胶子使用 sin；TMD 使用 Hermitian 半轴 cos/sin}\)\\
\(\text{Quadrature: }\text{把 }z\text{ 插值到均匀网格，单位统一为 }\mathrm{GeV}^{-1}\text{，使用梯形权重}\)\\
\(\text{Fourier: }\widetilde g\gets\text{式 }\eqref{eq:quasi-pdf-sin-clean}\text{ 或 }x\widetilde g\gets\text{式 }\eqref{eq:quasi-tmd-fourier-clean}\)\\
\(\text{Grid gate: }z\text{ 从 }0\text{ 开始且无重复；matching 网格拒绝精确 }x=0\)\\
\(\text{Kernel: }\xi_{ij}\gets x_i/y_j,\quad g_{ij}\gets g_-,g_{\rm mid},g_+\text{ 和 }\operatorname{Si}\text{ 项}\)\\
\(\text{Match: }Z\gets I+\alpha_s M,\quad h_R^{\rm PDF}\gets Z^{-1}h_R^{\rm quasi}\)\\
\(\text{Power corrections: }\text{保留 }M^2/P_z^2,\Lambda^2/P_z^2,a^2P_z^2\text{ 的 }P_z\text{ 外推}\)\\
\(\text{Output: }x\text{ 网格、quasi/PDF、协方差、截断/匹配参数和幂次修正系统误差}\)\\
\bottomrule
\end{tabular}
\end{table}
\src{../PyQCD/pyqcd/renorm/_matching.py:20--177；../PyQCD/pyqcd/renorm/_tmdextract.py:230--304；refer/papers/准部分子分布单圈匹配_latex/chapters/section03.tex；refer/papers/大动量有效理论获取胶子分布_latex/chapters/section02.tex}

\section{证据边界}
\texttt{matching\_tmd\_boundary} 审计只令 \(Z=I+\alpha_sM\)，验证
\((I-\alpha_sM)(I+\alpha_sM)=I+O(\alpha_s^2)\) 和树级恒等；这不能
验证实际的 \(g_1/g_2/g_3\) 圈积分、主值处方、软函数、rapidity 重整化或真实
PDF 卷积。最终 PDF 还需要多个 \(P_z\)、体积和格距的联合外推，并检查
支持区间、矩和规范化。

\chapter{Wilson/梯度流、RK3 与小流时间展开}
\label{sec:flow}

\section{连续流方程与平滑半径}
Wilson flow 把初始规范势 \(A_\mu(x)\) 演化为 \(B_\mu(t,x)\)：
\begin{equation}
 \partial_tB_\mu=D_\nu G_{\nu\mu},\qquad B_\mu(0,x)=A_\mu(x).
 \label{eq:continuum-flow-clean}
\end{equation}
其中 \(G_{\mu\nu}\) 由流场 \(B_\mu\) 构造。它是规范场空间上的抛物型
扩散方程，流时间 \(t\) 的质量维数为 \(-2\)，在 \(t>0\) 对局域场的
紫外涨落起平滑作用。树级热核给出典型平滑半径
\begin{equation}
 r_{\rm sm}=\sqrt{8t}.
 \label{eq:smoothing-radius-clean}
\end{equation}
在格点接口中传入的是 \(\tau=t/a^2\)，因此
\begin{equation}
 t[\mathrm{GeV}^{-2}]
 =\tau\left(\frac{a[\mathrm{fm}]}{\hbar c[\mathrm{GeV\,fm}]}\right)^2,
 \qquad
 r_{\rm sm}[\mathrm{fm}]=a[\mathrm{fm}]\sqrt{8\tau}.
 \label{eq:flow-conversion-clean}
\end{equation}
一个可用的窗口需要同时满足 \(a\ll r_{\rm sm}\) 以压低离散噪声，并且
\(r_{\rm sm}\) 小于有限体积、算符分离和目标短距尺度；不存在窗口时，流
时间扫描不能被解释为连续物理量的稳定平台。

\section{格点流导数与三阶 Runge--Kutta}
对链接 \(V_\mu(x)\)，六个 staple 的和记为 \(\Omega_\mu(x)\)。流导数
从 \(\Omega_\mu V_\mu^\dagger\) 的反厄米、无迹部分取出：
\begin{align}
 Z_\mu(V)&=P_{\mathfrak{su}(3)}[\Omega_\mu V_\mu^\dagger],\\
 P_{\mathfrak{su}(3)}(X)
 &=\frac{X-X^\dagger}{2}
 -\frac{\operatorname{Tr}(X-X^\dagger)}{2N_c}I,\\
 \dot V_\mu&=Z_\mu(V)V_\mu.
\end{align}
由于 \(Z_\mu^\dagger=-Z_\mu\) 且 \(\operatorname{Tr}Z_\mu=0\)，精确
\(\exp(\epsilon Z_\mu)\) 属于 \(SU(3)\)。代码用三阶 Lüscher RK 组合：
\begin{align}
 W_0&=V_t,\\
 W_1&=\exp\left(\frac{\epsilon}{4}Z_0\right)W_0,\\
 W_2&=\exp\left(\frac{8\epsilon}{9}Z_1-\frac{17\epsilon}{36}Z_0\right)W_1,\\
 V_{t+\epsilon}
 &=\exp\left(\frac{3\epsilon}{4}Z_2-\frac{8\epsilon}{9}Z_1
 +\frac{17\epsilon}{36}Z_0\right)W_2.
 \label{eq:rk3-flow-clean}
\end{align}
每次 \(\epsilon\) 由 \(\tau/n_{\rm step}\) 给出；默认 \(\epsilon=0.01\)，
也可以显式指定步数。实现的 \texttt{su3\_exp} 用四阶截断级数并做
determinant 归一，这对小步长是工程近似，和精确矩阵指数不是同一个
误差模型；应通过 \(\epsilon\) 收敛扫描确认。

Wilson 作用量密度和 Clover 场强能量密度是两个不同离散量：
\begin{align}
 s_W(x)&=\frac16\sum_{\mu<\nu}
 \left[1-\frac1{N_c}\Re\operatorname{Tr}P_{\mu\nu}(x)\right],\\
 E_{\rm cl}(t,x)&=\sum_{\mu<\nu}
 \operatorname{Tr}\left[F_{\mu\nu}^{\rm traceless}(t,x)
 F_{\mu\nu}^{\rm traceless}(t,x)\right].
\end{align}
连续流的梯度性质应优先用 \(s_W\) 判定；粗糙组态上不能要求 Clover
\(E_{\rm cl}\) 逐点严格单调。尺度设定接口以
\begin{equation}
 \tau_0^2\langle E(\tau_0)\rangle=c
\end{equation}
的离散插值定义参考流时间；这里 \(c\) 是约定参数，不能从代码默认值
推断成通用物理常数。
\src{../PyQCD/pyqcd/renorm/_gradient_flow.py:38--267；refer/papers/Properties_Uses_Wilson_Flow_latex/chapters/section02.tex；refer/papers/梯度流的微扰分析_latex/chapters/section02.tex；refer/papers/杨米尔斯梯度流提取能动量张量_latex/chapters/section02.tex}

\section{小流时间展开和 SFTX}
小流时间展开将流算符表示为同一重整化方案的局域算符：
\begin{equation}
 \mathcal O_{\rm flow}(t,x)
 =c_1(t,\mu)\mathcal O_{\overline{\rm MS}}(\mu,x)
 +t\,c_2(t,\mu)\mathcal O_6(x)+\cdots.
 \label{eq:sftx-clean}
\end{equation}
对于胶子一圈匹配接口，工作实现采用
\begin{equation}
 \mathcal O_{\overline{\rm MS}}(\mu)
 =\left[1+\frac{\alpha_s(\mu)}{4\pi}c(t,\mu)\right]\mathcal O_{\rm flow}(t),
 \qquad
 c(t,\mu)=2b_0[\log(2\mu^2t)+\gamma_E].
 \label{eq:sftx-coeff-clean}
\end{equation}
能量密度的树级量纲为
\begin{equation}
 \langle E(t)\rangle
 =\frac{3(N_c^2-1)}{16\pi^2t^2}
 [1+O(\alpha_s)+O(t\Lambda_{\rm QCD}^2)],
 \qquad [E]=4.
 \label{eq:flow-energy-clean}
\end{equation}
因此 \(t^2E(t)\) 无量纲，\(\log(2\mu^2t)\) 的参数也无量纲。SFTX
需要 \(a^2\ll t\ll\Lambda_{\rm QCD}^{-2}\) 的小流时间窗口；它消除的是
局域短距算符的 UV 定义差异，不自动消除 finite-staple TMD 的 rapidity
发散和软因子。

\begin{table}[htbp]
\centering
\caption{Wilson flow 与 SFTX 的单列算法}
\small
\begin{tabular}{@{}P{0.94\textwidth}@{}}
\toprule
\(\text{Input: }U,\;\tau=t/a^2,\;\epsilon,\;a[\mathrm{fm}],\;\mu,\;N_f,\;\text{流观测量类型}\)\\
\(\text{Validate: }\tau\ge0,\;\epsilon>0,\;a>0,\;\text{输入链接轴序为 }(t,z,y,x,\mu,a,b)\)\\
\(\text{Flow step: }\Omega_\mu\gets\text{六 staple},\quad Z_\mu\gets P_{\mathfrak{su}(3)}(\Omega_\mu V_\mu^\dagger)\)\\
\(\text{RK3: }V\gets\text{式 }\eqref{eq:rk3-flow-clean}\text{，每步用新的 }Z(W_i)\text{，并保持同一后端}\)\\
\(\text{Group check: }V_\mu^\dagger V_\mu\simeq I,\;\det V_\mu\simeq1\text{；随 }\epsilon\text{ 缩小收敛}\)\\
\(\text{Flow observable: }s_W\text{ 用 Wilson plaquette；}E_{\rm cl}\text{ 用无迹 Clover，二者不要混名}\)\\
\(\text{Scale option: }\tau_0^2\langle E(\tau_0)\rangle=c\text{，在扫描网格内插值并报告 }c\text{ 与窗口}\)\\
\(\text{Unit map: }t_{\mathrm{GeV}^{-2}}\gets\tau(a/\hbar c)^2,\quad r_{\rm sm}\gets\sqrt{8t}\)\\
\(\text{SFTX: }c(t,\mu)\gets2b_0[\log(2\mu^2t)+\gamma_E],\quad O_{\overline{\rm MS}}\gets[1+\alpha_s c/(4\pi)]O_{\rm flow}\)\\
\(\text{Window gate: }a^2\ll t\ll\Lambda_{\rm QCD}^{-2}\text{，另检查非局域路径尺度和 }b_T\text{ 几何}\)\\
\(\text{Output: }V_\tau,s_W,E(t),\tau_0,t_{\mathrm{GeV}^{-2}},r_{\rm sm},\text{SFTX 系数与误差元数据}\)\\
\bottomrule
\end{tabular}
\end{table}
\src{../PyQCD/pyqcd/renorm/_gradient_flow.py:139--267；../PyQCD/pyqcd/renorm/_tmdextract.py:318--381；refer/papers/准部分子分布与梯度流_latex/chapters/section02.tex}

\section{SymPy 与流实现的边界}
已有核心 SymPy 推导验证了流方程的五维形式、Duhamel 解、RK3 三阶全局
阶、参考尺度的无量纲性和简化格点流的作用量下降；新增
\texttt{gradient\_flow\_sftx\_dimensions} 审计验证了
\(t=\tau(a/\hbar c)^2\)、\(r_{\rm sm}=\sqrt{8t}\) 和
\([t^2E]=0\)。这些是解析/量纲证据，不是对生产组态流演化、有限步长误差
或小流时间窗口的数值替代。

\chapter{胶子 quasi-TMD、软因子与 Collins--Soper 演化}
\label{sec:tmd}

\section{finite-staple 几何和规范闭合}
TMD 与直线 quasi-PDF 的关键区别是横向分离 \(b_T\) 和 finite staple。
令纵向方向为 \(z\)，横向方向为 \(b\)，从端点 \(x\) 到
\(y=x+z\hat z+b\hat b\) 的 staple 可以写成三段：
\begin{align}
 W_{\sqcup}(x,y)
 &=U_z^\dagger((z+L)\hat z+b\hat b,b\hat b)\,
 U_b((z+L)\hat z+b\hat b,(z+L)\hat z)\nonumber\\
 &\quad\times U_z((z+L)\hat z,z\hat z).
 \label{eq:staple-clean}
\end{align}
这里省略了每一段起点 \(x\) 的平移；第一段长度 \(-L\)，第二段长度
\(b\)，第三段长度 \(L+z\)。若 \(b=0\) 或 \(L=0\)，几何可能退化，
但这不自动意味着 TMD soft subtraction 的所有极限都连续。
\(W_{\sqcup}\) 的端点变换为
\begin{equation}
 W_{\sqcup}(x,y)\mapsto G(x)W_{\sqcup}(x,y)G^\dagger(y).
\end{equation}
因此双场强闭合颜色：
\begin{equation}
 M^{\mu\lambda;\nu\rho}(z,b_T)
 =\left\langle
 \operatorname{Tr}\left[
 F^{\nu\rho}(x)W_{\sqcup}(x,y)
 F^{\mu\lambda}(y)W_{\sqcup}^\dagger(x,y)\right]\right\rangle.
 \label{eq:gluon-tmd-bilocal-clean}
\end{equation}
PyQCD 的 TMD 路径明确在每个场强上做 \(su(N_c)\) 无迹投影：
\begin{equation}
 F\longmapsto F-\frac{\operatorname{Tr}F}{N_c}I.
\end{equation}
这与第 \ref{sec:ope} 章的 finite-\(a\)
\texttt{legacy\_untraced} 直线 OPE 是两个不同语义层，不能共用一个
无元数据的数组。
\src{../PyQCD/pyqcd/renorm/_tmd.py:53--205；../PyQCD/docs/格点QCD中的TMD_PDF.tex；refer/papers/从准TMD波函数确定CS核_latex/chapters/section02.tex}

\section{胶子非极化组合及数组接口}
取纵向方向 \(k\) 之外的两个空间方向 \(i,j\)，非极化组合为
\begin{equation}
 \mathcal O_g(z,b_T)
 =M^{ti;ti}(z,b_T)+M^{tj;tj}(z,b_T)
 -2M^{ij;ij}(z,b_T).
 \label{eq:tmd-combination-clean}
\end{equation}
它的系数依赖场强和极化张量约定；不能从“看起来是三项”推断为任意
归一化。\codeid{gluon_tmd_operator} 返回逐时空点值
\((N_t,N_z,N_y,N_x)\)，\texttt{tmd\_matrix\_elements} 复用一次场强计算，
对 \(z\_list,b\_list\) 循环，空间求和后平均时间，返回
\((N_z,N_b)\)。\texttt{tmd\_matrix\_elements\_time} 保留时间轴，返回
\((N_z,N_b,N_t)\)，供不相连三点的 \(C_3=C_2O\) 因子化使用。
默认 \texttt{L=None} 时同一批次取 \(\max|z|\)，因此一个批次内的 staple
臂长必须作为元数据保存。

在共线交叉极限 \(b_T\to0\) 时，TMD 算符可与胶子 quasi-PDF 的纵向
算符比较；但真正的 TMD 定义还含软因子和 rapidity 方案，故只能将此
比较作为交叉检查，而不能当作无条件的数学等同。

\section{软因子和 Collins--Soper 核}
TMD 相关函数含有 Wilson 线 cusp 和 rapidity 发散。一个抽象的重整化
组合可写为
\begin{equation}
 \mathcal F_{\rm TMD}(b_T;\mu,\zeta)
 =\frac{\mathcal F_{\rm unsub}(b_T)}
 {\sqrt{S(b_T;\mu,\rho)}}\,R_{\rm rap}(\mu,\zeta,\rho),
 \label{eq:tmd-soft-rapidity-clean}
\end{equation}
其中 \(\rho\) 是 rapidity 调节器，最终物理量应独立于中间调节器。工作
接口 \texttt{soft\_function\_intrinsic} 提供的是归一化框架
\begin{equation}
 S_I(b_T,\mu)=\frac{R^2(b_T)}{R^2(b_T\to0)}.
 \label{eq:intrinsic-soft-clean}
\end{equation}
它需要同一 staple 几何和同一流/重整化方案的输入；代码的归一化函数
本身不等于从独立真空 Wilson loop 完整计算的 soft factor。

不同大动量下的矩阵元比值可给出 Collins--Soper 斜率：
\begin{equation}
 K(b_T)\simeq
 \frac{\log[h_R(z_0,b_T,P_{z1})/
 h_R(z_0,b_T,P_{z2})]}
 {\log(P_{z1}/P_{z2})}.
 \label{eq:cs-slope-clean}
\end{equation}
这里 \(z_0\) 必须选在信号和 scheme 受控的参考点；若 \(z=0\) 的矩阵元
恒等为一或噪声过大，接口提供 \texttt{z\_ref} 选择非零行。不同文献对
\(K\)、\(\gamma_\zeta\) 的符号及二分之一因子有约定差异，报告中必须
沿用定义式 \eqref{eq:cs-slope-clean}，不能只写“CS kernel”而省略
归一化。

\section{TMD 匹配闭环}
在一个明确方案下，准 TMD 与光锥 TMD 的因子化结构可以写为
\begin{equation}
 x\widetilde g(x,b_T;P_z)
 =\int\frac{\dd y}{|y|}
 C^{\rm hybrid}\left(\frac{x}{y},
 \lambda_s,\frac{\mu}{|y|P_z}\right)
 \frac{y\,g(y,b_T;\mu,\zeta)}{\sqrt{S_I(b_T,\mu)}}
 \exp\left[\frac12\log\frac{\zeta_z}{\zeta}K(b_T,\mu)\right].
 \label{eq:tmd-matching-clean}
\end{equation}
实现把这一卷积离散为矩阵 \(Z\)，对每个 \(b_T\) 乘上软因子和快度演化
因子，再求 \(Z^{-1}\)。这一步要求输入 \texttt{x\_tmd} 的含义已经是
\(y\,g(y,b_T)\)，而不是 \(g(y,b_T)\)；否则会多/少一个 \(y\) 权重。
\texttt{x\_grid} 不可含零，匹配矩阵的 \(y\) 网格需与 \(x\) 网格同维。

\begin{table}[htbp]
\centering
\caption{胶子 quasi-TMD 从流场到匹配输出的单列算法}
\small
\begin{tabular}{@{}P{0.94\textwidth}@{}}
\toprule
\(\text{Input: }U,\;\tau,\;\texttt{z\_list},\;\texttt{b\_list},\;z_{\rm dir},b_{\rm dir},L,\;P_z,\;\text{颜色归一化}\)\\
\(\text{Flow option: }V\gets\texttt{wilson\_flow}(U,\tau,\epsilon)\text{，或明确声明不做流}\)\\
\(\text{Geometry: }W_{\sqcup}\gets\text{式 }\eqref{eq:staple-clean}\text{，固定同一批次的 }L\text{ 和周期路径}\)\\
\(\text{Fields: }(F_{ti},F_{tj},F_{ij})\gets\text{Clover，并逐点投影到无迹色矩阵}\)\\
\(\text{Bilocal: }M^{\mu\lambda;\nu\rho}\gets\operatorname{Tr}[F(x)W_{\sqcup}F(y)W_{\sqcup}^{\dagger}]\)\\
\(\text{Combine: }O_g\gets M^{ti;ti}+M^{tj;tj}-2M^{ij;ij}\)\\
\(\text{Contract: }O(z,b)\gets\sum_{\boldsymbol x}O_g(x;z,b)\text{；按需平均时间或保留 }t\)\\
\(\text{Renormalize: }O_B\to h_R\text{，明确是 ratio、}Z_R\text{、hybrid 或 flow/SFTX 方案}\)\\
\(\text{Soft: }S_I(b)\gets R^2(b)/R^2(0)\text{，确认与 staple/流方案一致}\)\\
\(\text{CS: }K(b)\gets\log[h_R(P_{z1})/h_R(P_{z2})]/\log(P_{z1}/P_{z2})\text{，保存参考 }z_0\)\\
\(\text{Fourier: }h_R(z,b)\to x\widetilde g(x,b;P_z)\text{，用 Hermitian 半轴和统一单位}\)\\
\(\text{Matching: }Z\gets I+\alpha_sM\text{，加入 }S_I\text{ 和 }K\text{ 快度因子，求 }Z^{-1}\)\\
\(\text{Gate: }\text{检查 }b_T,L,\tau,P_z,\mu,\zeta,\text{软因子和 rapidity 方案是否全部有来源}\)\\
\(\text{Output: }O(z,b),h_R,x\widetilde g,K,S_I,\text{匹配 PDF、协方差和完整方案元数据}\)\\
\bottomrule
\end{tabular}
\end{table}
\src{../PyQCD/pyqcd/renorm/_tmd.py:140--415；../PyQCD/pyqcd/renorm/_tmdextract.py:30--304；refer/papers/TMD_Soft_Function_LaMET_latex/chapters/section02.tex；refer/papers/从准TMD波函数确定CS核_latex/chapters/section02.tex}

\section{TMD 的未验证边界}
梯度流确实提供 UV 平滑并能使小流时间展开有定义，但它不单独提供
TMD 所需的 soft factor、rapidity counterterm、cusp anomalous dimension、
Collins--Soper 核或 LaMET matching。当前接口和本报告的
\texttt{matching\_tmd\_boundary} 审计只支持公式骨架与一阶矩阵代数；
没有真实同几何软函数、多个 \(P_z\) 和完整非微扰卷积时，最终 TMD
必须标为 \unverified。

\subsection{边界矩阵与极限检查}
为了避免把“接口存在”误读成“物理对象已闭合”，可以把几条关键链的几何
支撑和极限写成一张边界矩阵：
\begin{table}[htbp]
\centering
\caption{OPE、准 PDF、准 TMD 与流重整化的边界矩阵}
\small
\begin{tabularx}{0.96\textwidth}{@{}P{3.0cm}P{4.2cm}P{3.3cm}Y@{}}
\toprule
对象 & 几何支撑 & 可接受极限 & 失败模式\\
\midrule
\codeid{gluon_ope_operator_z0} & straight Wilson line, \(z\)-方向 & \(z=0\) 退化为同点场强收缩 & 不能替代 finite-staple TMD 的横向几何\\
\texttt{quasi\_pdf\_gluon} & 单轴 \(z\) 半轴积分 & \(P_z\to\infty\) 才逼近光锥 PDF & \(x=0\) 节点没有主值处方\\
\texttt{quasi\_tmd\_pdf} & Hermitian 半轴 + \(b_\perp\) 依赖 & \(b_\perp\to0\) 只作交叉检查 & 不能与 sin 通道混用\\
\codeid{gluon_tmd_operator} & staple geometry + nonzero \(b_\perp\) & \(L\) 只影响几何延展 & 若无 soft/rapidity 处理，物理量不闭合\\
\texttt{wilson\_flow} & 流时间 \(t>0\) 的局域平滑 & \(t\to0^+\) 回到裸算符 & 不能单独替代 TMD renorm\\
\codeid{sftx_gluon_matching_coeff} & 小流时间窗口 & \(a^2\ll t\ll\Lambda_{\rm QCD}^{-2}\) & 无窗口时只剩形式匹配\\
\codeid{hR_PDF/tmd_matching_hybrid} & 离散卷积矩阵 & 网格细化后应稳定收敛 & 错用 \(x=0\) 或 \(y\) 权重会破坏主值结构\\
\bottomrule
\end{tabularx}
\end{table}

这些极限里最容易误判的是三件事。第一，\(b_\perp\to0\) 只是几何退化，
不是自动的 soft 因子消失。第二，\(t\to0^+\) 只是回到裸算符，不会替代
完整重整化。第三，\(P_z\to\infty\) 是因子化极限，不是有限动量下的数值
恒等式。换句话说，\texttt{gluon\_ope\_operator\_z0}、\texttt{quasi\_pdf\_gluon}、
\texttt{quasi\_tmd\_pdf} 和 \texttt{gluon\_tmd\_operator} 分别处在四条不同的
物理链上，只有共享部分几何和方案标签时才可交叉核对，不能直接互相替代。

\begin{equation}
 \lim_{b_\perp\to0}\mathcal O_g(z,b_\perp)\neq \mathcal O_{\rm straight}(z),
 \qquad
 \lim_{t\to0^+}\mathcal O_{\rm flow}(t)=\mathcal O_{\rm bare}+O(t),
\end{equation}
\begin{equation}
 \lim_{P_z\to\infty}\widetilde g(x,b_\perp,P_z)=g(x,b_\perp)
 \quad\text{only after matching and power corrections are controlled}.
\end{equation}

\iffalse
\section{两点和三点函数的谱分解}
设传递矩阵存在、Euclidean 反射正性成立，且算符能够与能量本征态重叠。
采用相对论归一化
\(\langle n,\boldsymbol p|n,\boldsymbol p\rangle=2E_nV\)，两点函数的无限时间近似为
\begin{align}
 C_2(t,\boldsymbol p)
 &=\sum_n\frac{|Z_n(\boldsymbol p)|^2}{2E_n(\boldsymbol p)}e^{-E_n t},
 \\[-2pt]
 C_2^{\rm periodic}(t)
 &=\sum_n\frac{|Z_n|^2}{2E_n}\left[e^{-E_nt}+e^{-E_n(N_t-t)}\right].
\end{align}
有限 (N_t)、反周期费米边界和算符的反向传播贡献会改变拟合模型，
不能把第一式在所有 (t) 上硬套。三点函数在源—插入—汇时间顺序下为
\begin{equation}
 C_3(t_{\rm sep},\tau)=
 \sum_{m,n}\frac{Z_m^f Z_n^{i*}}{4E_m^fE_n^i}
 e^{-E_m^f(t_{\rm sep}-\tau)}e^{-E_n^i\tau}
 \langle m,f|\mathcal O|n,i\rangle.
 \label{eq:c3-spectrum}
\end{equation}
当 (\tau)、(t_{\rm sep}-\tau) 都大于激发态时间尺度时，基态项支配；
否则每一个省略项都带有 (e^{-\Delta E\tau}) 或
\(e^{-\Delta E(t_{\rm sep}-\tau)}) 的系统误差。

对于胶子 OPE 的不相连因子化分析，工作树中的中间量是
\begin{equation}
 C_3^{\rm fact}(t,\tau,z)=C_2(t)\,O(\tau,z),
 \qquad
 C_3^{\rm disc}=C_3^{\rm fact}-C_2\,\langle O\rangle,
 \qquad
 R(z)=\left\langle\frac{C_3^{\rm disc}}{C_2}\right\rangle_{t_{\rm src}}.
 \label{eq:disconnected-ratio}
\end{equation}
式中的真空扣除是 disconnected 矩阵元的定义组成部分；它既不是普通
统计均值，也不能在重采样之后随意改成两个独立均值的商。代码将相对时间
通过 `roll` 对齐源点，因此周期边界是该操作正确性的必要假设。
\src{../PyQCD/pyqcd/analysis/_analyse.py:361--520；../PyQCD/pyqcd/analysis/_ratio2pt.py:80--194；../PyQCD/pyqcd/analysis/_disconnected.py:1--13；../PyQCD/docs/格点QCD中的3pt构造.tex}

\section{本章的定理、假设与边界}
\begin{table}[htbp]
\centering
\caption{Wick 与关联函数理论链的公理—结论对应}
\small
\begin{tabularx}{0.96\textwidth}{@{}P{2.8cm}Y Y@{}}
\toprule
公理/假设 & 数学结论 & 在代码中的验证边界\\
\midrule
Grassmann 反对易、高斯费米测度、(D) 可逆 & Wick 展开及费米置换号 & 可由有限索引枚举；不证明真实 (D^{-1}) 已求精确。\\
局域规范协变性、颜色 singlet 算符 & 颜色指标闭合、关联函数规范不变 & VVV/VDV 标签结构可审计；不验证输入 eigenvector 与 gauge 同组态。\\
反射正性和传递矩阵 & Euclidean 能谱非负、长时间基态支配 & 需要正确作用量和边界；代数脚本不证明反射正性。\\
(gamma_5Dgamma_5=D^\dagger) & 式 \eqref{eq:seq-peram} 的反向传播子关系 & gamma 矩阵代数通过；具体费米子核仍需实测。\\
平移不变性与周期/反周期边界 & 源时平均、动量投影、`roll` 对齐 & 仅检查公式；边界符号和数据布局需用真实配置回归。\\
\bottomrule
\end{tabularx}
\end{table}


\chapter{谱学提取、重采样与相关拟合}
\label{sec:spectrum}

\section{动量、能量和有效质量}
代码中的动量列表是整数单位 \(\boldsymbol n=(n_z,n_y,n_x)\)，在空间长度
\(N_xa\)（立方空间）下
\begin{equation}
 |\boldsymbol p|=\frac{2\pi}{N_xa}|\boldsymbol n|,
 \qquad
 E(\boldsymbol p)=\sqrt{M_0^2+|\boldsymbol p|^2}.
 \label{eq:mom2gev}
\end{equation}
\texttt{Mom2GeV} 的一个容易误读之处是：若 `Mom` 是标量，它被解释为
\(|\boldsymbol n|^2\)，而若是 `[pz,py,px]` 列表，函数先求平方和；
因此调用者不能把标量 \(n\) 当作向量模长而再次平方。量纲检查要求
从 fm 转成 \(\mathrm{GeV}^{-1}\)，使用
\(\hbar c=0.1973269804\,\mathrm{GeV\,fm}\)。

单指数关联函数给出对数有效质量
\begin{equation}
 am_{\rm eff}^{\log}(t)=\log\frac{C(t)}{C(t+1)},
 \qquad
 m_{\rm eff}[\mathrm{GeV}]
 =\frac{\hbar c}{a[\mathrm{fm}]}
 \log\frac{C(t)}{C(t+1)}.
\end{equation}
周期边界下两向传播的有效质量满足
\begin{equation}
 m_{\rm eff}^{\cosh}(t)=
 \operatorname{arccosh}\frac{C(t+2)+C(t)}{2C(t+1)},
 \qquad
 m[\mathrm{GeV}]=\frac{\hbar c}{a[\mathrm{fm}]},m_{\rm eff}^{\cosh}.
 \label{eq:cosh-meff}
\end{equation}
PyQCD 在复数或非浮点输入上先取绝对值；在噪声使 arccosh 自变量小于
1 时，把结果置零。这是数值保护和信号筛选，不是物理上“质量为零”的
证明；报告结果时必须保留该掩码/截断规则。

对于算符基底 (O_i)，广义特征值问题为
\begin{equation}
 C(t)v_n(t,t_0)=\lambda_n(t,t_0)C(t_0)v_n(t,t_0),
 \qquad
 \lambda_n(t,t_0)\simeq e^{-E_n(t-t_0)}
 \left[1+O(e^{-\Delta E_n(t-t_0)})\right].
 \label{eq:gevp}
\end{equation}
需要 (C(t_0)) Hermitian 正定、算符重叠矩阵在目标子空间满秩。代码先做
(C\leftarrow(C+C^\dagger)/2)，再调用广义 Hermitian 特征值求解器，并按
(t<t_0) 与 (t\ge t_0) 使用不同排序；这解决的是复数舍入和排序约定，
不替代对 (t_0)、基底维数和激发态系统误差的扫描。
\src{../PyQCD/pyqcd/analysis/_analyse.py:30--84,253--354,599--664；refer/books/INTRODUCTION_TO_LATTICE_QCD_latex/chapters/sec17_correlators.tex}

\section{Jackknife、bootstrap 与协方差}
设有 (N) 个配置观测 \(x_i\)，均值为 \(\bar x=N^{-1}\sum_i x_i\)。
删除第 (i) 个配置的 Jackknife 样本是
\begin{equation}
 x_{(-i)}=\frac{\sum_jx_j-x_i}{N-1},
 \qquad
 \widehat{\operatorname{Cov}}_J(\bar x)
 =\frac{N-1}{N}\sum_{i=1}^{N}
 (x_{(-i)}-\bar x)(x_{(-i)}-\bar x)^T.
 \label{eq:jack-cov}
\end{equation}
如果 `cov_axes` 指定多个时间/空间点，代码先把非协方差轴保留，再把
协方差轴展平，执行 (r_{ni}r_{nj}) 的 Einstein 求和，最后恢复两份
协方差轴。复杂观测量可用实部、虚部拼成联合向量，但必须在统计说明中
明确该选择。

Bootstrap 以有放回索引 (I_b=(i_1,\ldots,i_M)) 生成
\begin{equation}
 x^{(b)}=\frac1M\sum_{k=1}^{M}x_{i_k},
 \qquad
 \widehat{\operatorname{Cov}}_B
 =\frac1{N_B}\sum_b(x^{(b)}-\bar x_B)(x^{(b)}-\bar x_B)^T.
\end{equation}
PyQCD 的默认 (M=\max(N-5,1))、(N_B=4N)，并把全配置均值放在
第一个样本位置；随机种子、配置自相关和阻塞策略应由上层管线固定，不能
把默认参数误认为普适最优选择。删除一配置 Jackknife 在 (N<2) 时没有
可辨识协方差，当前实现会留下 NaN/状态标记而不强行拟合。

源时平移平均利用
\begin{equation}
 C(\Delta t)=\frac1{N_t}\sum_{t_{\rm src}}
 C(t_{\rm src},t_{\rm src}+\Delta t),
 \qquad
 \Delta t=(t_{\rm sink}-t_{\rm src})\bmod N_t.
\end{equation}
反周期边界在跨过时间边界时引入负号；三点函数还要根据 (t_{\rm sep})
确定尾部符号。`loop_tsrc` 采用 CPU 临时数组实现这一步，若输入来自 GPU，
返回时再送回当前后端。这个数据搬运行为应在性能报告中单独记录。

\section{相关拟合和可辨识性}
给定数据向量 (y)、模型 (f(\theta)) 和协方差 (C)，相关拟合的核心
目标是
\begin{equation}
 \chi^2(\theta)=
 [y-f(\theta)]^\dagger C^{-1}[y-f(\theta)].
 \label{eq:correlated-chi2}
\end{equation}
当 (C) 由有限 bootstrap/Jackknife 样本估计而秩不足时，直接求逆会放大
噪声。PyQCD 的拟合器先检查 Hermitian 半正定性，再对相关矩阵做 SVD
调节并报告样本秩、有效秩和条件数。若数据自由度不超过参数数目，或有效
秩不足，拟合状态为不可辨识；有 prior 时可以继续，但状态必须写成
``prior_constrained``，不能声称参数由数据独立决定。

\begin{table}[htbp]
\centering
\caption{谱学与统计分析的单列算法}
\small
\begin{tabular}{@{}P{0.94\textwidth}@{}}
\toprule
\(\text{Input: }C_{\rm raw}[N_{\rm conf},\ldots,t],\;a[\mathrm{fm}],\;\text{边界条件},\;\text{拟合窗口}\)\\
\(\text{Validate: }N_{\rm conf}\ge2\text{（Jackknife），时间轴和源汇轴不重复且长度相容}\)\\
\(\text{Average: }C(t_{\rm src},t_{\rm sink})\to C(\Delta t)\text{，按周期/反周期边界处理符号}\)\\
\(\text{Resample: }\{C_i\}\to\{C_{(-i)}\}\text{ 或 }\{C^{(b)}\}\text{，保留配置轴在已知位置}\)\\
\(\text{Covariance: }r_{ni}=x_{ni}-\bar x_i,\;C_{ij}\gets\sum_n r_{ni}r_{nj}\text{ 并使用相应归一化}\)\\
\(\text{Spectrum option 1: }m_{\rm eff}\gets\log[C(t)/C(t+1)]\cdot\hbar c/a\)\\
\(\text{Spectrum option 2: }m_{\rm eff}\gets\operatorname{arccosh}[(C(t+2)+C(t))/(2C(t+1))]\cdot\hbar c/a\)\\
\(\text{Spectrum option 3: }C(t)v=\lambda(t,t_0)C(t_0)v\text{，检查 }C(t_0)>0\text{ 并排序 }\lambda\)\\
\(\text{Fit preparation: }y\gets \text{选定 }(t,\tau,z,p)\text{ 点，}\;C\gets \text{由同一重采样方案得到的协方差}\)\\
\(\text{Rank gate: }\operatorname{rank}_{\rm sample},\operatorname{rank}_{\rm SVD},N_{\rm data},N_{\rm par}\text{ 均记录}\)\\
\(\text{Fit: }\min_\theta\chi^2(\theta)\text{，必要时输入 prior，但标记 prior 对不可辨识方向的约束}\)\\
\(\text{Propagate: }\text{逐样本重复 }m_{\rm eff}/\text{GEVP}/\text{fit，最终从样本分布取均值和误差}\)\\
\(\text{Scan: }t_0,\;t_{\rm min},t_{\rm max},N_{\rm ev},P_z,\text{拟合模型和 SVD cut}\)\\
\(\text{Output: }\text{谱量、误差、协方差、条件数、有效秩、fit status、窗口和全部约定元数据}\)\\
\bottomrule
\end{tabular}
\end{table}
\src{../PyQCD/pyqcd/analysis/_analyse.py:91--162,169--246,527--664；../PyQCD/pyqcd/analysis/_fitter.py:59--256；../PyQCD/docs/格点QCD中的关联函数.tex}

\section{本章的 SymPy 复现}
新增 \texttt{jackknife\_covariance} 审计以二维符号观测验证式
\eqref{eq:jack-cov}；核心 \texttt{three\_point\_ratio\_spectral\_limit}
审计以 (C_2=|Z|^2e^{-Et}/(2E))、式 \eqref{eq:c3-spectrum} 的基态项
精确消去 overlap 和指数，得到
\begin{equation}
 R\longrightarrow\frac{\langle f|O|i\rangle}{2\sqrt{E_fE_i}}.
\end{equation}
这两个检查确证的是归一化和有限维代数；它们不测量真实 (C_2,C_3)，也
不排除有限 (t_{\rm sep}) 的激发态污染。
\iffalse

\chapter{Grassmann 积分、Wick 收缩与关联函数}
\label{sec:wick-clean}

\section{从路径积分到传播子}
费米子部分在给定规范背景下是 Grassmann 高斯积分。将全部颜色、Dirac、
味和格点坐标压缩为复合指标 \(i,j\)，写成
\begin{align}
 Z_F[\eta,\bar\eta;U]
 &=\int\mathcal D\bar\psi\mathcal D\psi\,
 \exp\left[-\bar\psi D[U]\psi+\bar\eta\psi+\bar\psi\eta\right]
 \nonumber\\
 &=\det D[U]\,\exp\left(\bar\eta D^{-1}[U]\eta\right),
 \qquad S[U]=D^{-1}[U].
\end{align}
这里要求 \(D\) 在所选边界、质量和规范背景下可逆；零模或接近零模时，
数值求逆和统计估计都要另作处理。对 Grassmann 源求导得到 Wick 定理：
\begin{align}
 \langle\psi_{i_1}\cdots\psi_{i_n}
 \bar\psi_{j_1}\cdots\bar\psi_{j_n}\rangle_U
 &=\sum_{\pi\in S_n}(-1)^{\pi}
 \prod_{k=1}^{n}S_{i_kj_{\pi(k)}}[U],
 \\
 \langle\psi_i\bar\psi_j\rangle_U&=S_{ij}[U].
\end{align}
费米负号不是经验系数，而是把 Grassmann 场重新排序到成对收缩所需的
置换奇偶性。这是自动 Wick 引擎可以枚举而不能随意省略的第一性原理。
它依赖 Grassmann 反对易公理和高斯测度；相互作用通过 \(\det D\) 和规范
组态平均进入，而不是被上述有限维代数单独消除。

若费米子离散化满足
\begin{equation}
 D^\dagger=\gamma_5D\gamma_5,\qquad
 S^\dagger=\gamma_5S\gamma_5,
\end{equation}
则反向时间/源汇顺序的传播子可以从已有传播子重建。PyQCD 的
\texttt{seq\_peram} 实际执行
\begin{equation}
 \tau_{\rm seq}(t_{\rm sink},t_{\rm src})
 =\gamma_5\,\tau(t_{\rm src},t_{\rm sink})^\dagger\,\gamma_5,
 \label{eq:seq-peram-clean}
\end{equation}
其中厄米共轭同时交换 Dirac 和蒸馏指标。这个等式的前提是所用 \(D\)、
边界和 gamma 约定确实满足 \(\gamma_5\)-Hermiticity；代码接口本身不能
替调用者证明该前提。
\src{../PyQCD/pyqcd/contraction/_seqperam.py:20--59；refer/books/INTRODUCTION_TO_LATTICE_QCD_latex/chapters/sec17_correlators.tex}

\section{算符、味平衡和自动收缩}
一个双线性算符的典型形式为
\begin{equation}
 \mathcal O_M(x)=\bar q_a^f(x)\,\Gamma_{\alpha\beta}\,
 q_b^g(x)\,\mathcal C_{ab}^{fg},
\end{equation}
重子算符则以 \(\epsilon_{abc}\) 闭合颜色：
\begin{equation}
 \mathcal O_B(x)=\epsilon_{abc}(q_1^a q_2^b)q_3^c,
\end{equation}
并由 gamma 链、味结构和自旋投影决定具体态。PyQCD 的 DSL 用
\texttt{|} 分隔不同时间/粒子区间，用 \texttt{gamma\_5}、
\texttt{gamma\_7}、\texttt{gamma\_mu} 等字符串表示 Dirac 插入，用
\texttt{u}、\texttt{d} 和 \texttt{u\^d} 等标记夸克及反夸克。分隔符
不是物理场，而是解析器识别 VDV/VVV 区间和时间标签的元数据。

\begin{table}[htbp]
\centering
\caption{自动 Wick 收缩的单列算法：从算符 DSL 到可执行 contraction plan}
\small
\begin{tabular}{@{}P{0.94\textwidth}@{}}
\toprule
\(\text{Input: }\texttt{sink\_operators},\;\texttt{source\_operators},\;\texttt{curr\_operators},\;Cpt\in\{\text{bubble},2\text{pt},3\text{pt},4\text{pt}\}\)\\
\(\text{Normalize: }\text{若首尾没有 }\texttt{|}\text{，自动补上；2pt 情形丢弃 current 区间}\)\\
\(\text{Parse: }\text{把味标记分为 }q\text{ 与 }q^d，\text{把 gamma 和分隔符位置分别记录}\)\\
\(\text{Expand wildcard: }q\to\{u,d,s,c,b,t\},\quad l\to\{u,d\},\quad\text{逐组合代入}\)\\
\(\text{Balance test: }N_q(f)=N_{\bar q}(f)\text{；不平衡的味组合直接跳过}\)\\
\(\text{Label fields: }(q,\bar q)\text{ 按出现次序分配 }ab,cd,ef,\ldots\text{ 收缩标签}\)\\
\(\text{Enumerate: }\text{对每种味分别生成夸克到反夸克的所有排列，再做不同味匹配的笛卡尔积}\)\\
\(\text{Fermi sign: }\mathcal s=(-1)^{N_{\rm inv}},\quad N_{\rm inv}=\#\{(i,j):i<j,\,\ell_i>\ell_j\}\)\\
\(\text{Time parse: }|\cdots|\text{ 内含两个夸克给 VDV，含三个夸克给 VVV；区间标成 }t_{\rm sink},t_{\rm cur},t_{\rm src}\)\\
\(\text{Gamma parse: }\text{每个 gamma 读取相邻夸克标签，生成 gamma registry 的输入索引}\)\\
\(\text{Peram entry: }S_{q\bar q}[t_q,t_{\bar q}]\text{ 保存夸克位、反夸克位、味和 Dirac/ev 标签}\)\\
\(\text{Vertex entry: }\text{两夸克区间绑定 }V_{mn}\text{，三夸克区间绑定 }V_{mnl}\text{，并写入时间标签}\)\\
\(\text{Index projection: }\text{重复小写字母求和，自由小写字母保留；大写动量/链接标签保留一次}\)\\
\(\text{Diagram output: }\texttt{result\_indx},\texttt{result\_name},\texttt{result\_sign},\texttt{peram},\texttt{V},\texttt{gamma\_pos}\)\\
\(\text{For each diagram: }\mathcal C_D=\mathcal s_D\,\operatorname{einsum}(V_{\rm sink},S_1,\ldots,S_n,\Gamma,V_{\rm src})\)\\
\(\text{2pt: }\mathcal C_2(t)=\sum_D\mathcal C_D(t)\)\\
\(\text{3pt: }\mathcal C_3(t,\tau)=\sum_D\mathcal C_D(t,\tau;\Gamma_{\rm cur})\)\\
\(\text{4pt/bubble: }\text{保留所有 current/区间连通分支，并按 registry 解析每个中间量}\)\\
\(\text{Check: }\text{味平衡、颜色闭合、输出自由指标与预期轴序一致；否则停止而不猜测}\)\\
\(\text{Output: }\text{一个字典（无 wildcard）或字典列表（有 wildcard），供静态绘图和动态 contraction 使用}\)\\
\bottomrule
\end{tabular}
\end{table}
\src{../PyQCD/pyqcd/contraction/_autowick.py:39--59,208--445；../PyQCD/pyqcd/contraction/_dynamic.py:262--647；../PyQCD/docs/格点QCD中的维克收缩与傅里叶变换.tex}

该算法的“枚举所有排列”在 \(n_f\) 个相同味夸克/反夸克时规模含有
\(n_f!\) 因子；这是物理拓扑的完整性与计算复杂度之间的直接冲突。实际
高阶多强子问题应先利用同味交换对称性、颜色张量和等价图合并，再由
\texttt{dynamic\_contraction} 根据 einsum 路径和设备估算执行顺序。合并
必须保持相对费米号，不能只按字符串去重。

\section{从收缩图到平台区：比值、真空扣除与局域极限}
\label{sec:platform-bridge}
自动收缩一旦给出图的列表，只说明 Grassmann 代数是否闭合、费米号是否
正确；它还没有回答“这个图能不能在谱分解下变成一个稳定矩阵元”。要把
图变成物理量，必须再经过比值构造、平台判据和极限检查三道门。

对连接型三点函数，最常见的归一化比值写成
\begin{equation}
 R(t,\tau)
 =\frac{C_3(t,\tau)}
 {\sqrt{C_2^{\rm src}(t)\,C_2^{\rm snk}(t)}}
 =\mathcal M
 +A\,e^{-\Delta E\,\tau}
 +B\,e^{-\Delta E\,(t-\tau)}
 +O(e^{-\Delta E t}),
\end{equation}
其中 \(\mathcal M\) 是目标矩阵元，\(A,B\) 是两侧激发态污染。这个公式的
含义很具体：平台区不是“看起来平”，而是对 \(t\)、\(\tau\)、\(t-\tau\)
三者同时满足同一组误差控制。若窗口平移后 \(\mathcal M\) 漂移，或者
Jackknife/Bootstrap 的误差带对窗口选择极敏感，那么抽取出的不是矩阵元，
而只是某个特定窗口上的经验平均值。

对 disconnected 通道，真空扣除必须保留在同一条重采样链里：
\begin{equation}
 C_3^{\rm disc}(t,\tau)
 =\bigl\langle C_3(t,\tau)\bigr\rangle
 -\bigl\langle C_2(t)\bigr\rangle\,
  \bigl\langle O(\tau)\bigr\rangle.
\end{equation}
这里的减法不是后处理，而是定义的一部分。若把扣除放到另一套样本平均里，
协方差就被改写，最后得到的误差与原始估计量不再对应。因而
\texttt{analysis/\_disconnected.py} 的路径必须和 \texttt{analysis/\_analyse.py}
的源时平均、resample 与 covariance 约定同步，不能拆成两条彼此独立的统计链。

对 straight Wilson line 的 bilocal 算符，局域极限同样不是“把距离直接设为零”
就完事：
\begin{equation}
 \mathcal O(z)=\bar q(0)\Gamma W(0,z)q(z)
 =\mathcal O_{\rm local}+z^\mu\mathcal O_\mu^{(1)}+O(z^2).
\end{equation}
上式的真正可比对象是“经过重整化和匹配后的局域极限”，而不是裸算符的几何
退化。也因此，\texttt{gluon\_ope\_operator\_z0}、\texttt{gluon\_tmd\_operator}、
\texttt{quasi\_pdf\_gluon} 和 \texttt{quasi\_tmd\_pdf} 共享的是 Wilson 线语言，
不是同一个物理极限；前者是 straight bilocal，后者还带着 staple、\(b_\perp\)
与 soft/rapidity 因子。对应地，\texttt{matching\_tmd\_boundary} 只能做结构性
审计，因为完整 TMD 需要 soft factor、Collins--Soper 核和匹配核三者一起闭合。

\begin{table}[htbp]
\centering
\caption{从收缩输出到物理量的接口链与失效模式}
\small
\begin{tabularx}{0.96\textwidth}{@{}P{3.0cm}P{4.4cm}Y@{}}
\toprule
接口 & 物理前提 & 典型失效模式\\
\midrule
\texttt{seq\_peram} & \(\gamma_5\)-Hermiticity 与源汇顺序已知 & 反向传播子被错误当成独立对象，Dirac 共轭关系断裂。\\
\texttt{ratio\_3pt} & 三点/两点共享同一重采样链与源时平均 & 激发态污染未压到平台，ratio 只是局部平坦。\\
\texttt{run\_bare\_matrix} & 方向平均、多窗口拟合和 SVD cut 都稳定 & 裸矩阵元 \(c_0\) 对窗口漂移敏感，不能声称已收敛。\\
\texttt{gluon\_ope\_operator\_z0} & straight Wilson line 与场强插入的局域展开 & 把 \(z=0\) 误写成“无需匹配的局域等式”。\\
\texttt{gluon\_tmd\_operator} & staple 几何、\(b_\perp\) 和 soft factor 同时存在 & 只改横向距离不处理 rapidity，TMD 仍不闭合。\\
\texttt{matching\_tmd\_boundary} & OPE/LaMET/TMD 三条链的极限都被标注 & 把结构性审计误读成数值闭环，超出当前证据范围。\\
\bottomrule
\end{tabularx}
\end{table}

因此，真正的闭环顺序应当是：先确认 Grassmann 代数和 Wick 收缩闭合，再确认
比值与谱分解能给出稳定平台，最后才把平台量送入 OPE、LaMET、flow 或 TMD 的
匹配核。若这三个层次的边界没有分别标清，任何“最终结果”都只是把不同极限
混写在一起。  
\src{../PyQCD/pyqcd/contraction/_seqperam.py:20--59；../PyQCD/pyqcd/analysis/_ratio2pt.py:1--192；../PyQCD/pyqcd/analysis/_bare_matrix.py:1--160；../PyQCD/pyqcd/analysis/_disconnected.py:1--220；../PyQCD/pyqcd/operator/_gluon_ope.py:617--848；../PyQCD/pyqcd/renorm/_tmd.py:140--415；../PyQCD/pyqcd/renorm/_tmdextract.py:30--304；refer/books/INTRODUCTION_TO_LATTICE_QCD_latex/chapters/sec17_correlators.tex；refer/books/INTRODUCTION_TO_LATTICE_QCD_latex/chapters/sec18_hadron_spectrum.tex}
\fi

\chapter{Grassmann 积分、Wick 收缩与关联函数}
\label{sec:wick-clean}

\section{从路径积分到传播子}
费米子部分在给定规范背景下是 Grassmann 高斯积分。将全部颜色、Dirac、
味和格点坐标压缩为复合指标 \(i,j\)，写成
\begin{align}
 Z_F[\eta,\bar\eta;U]
 &=\int\mathcal D\bar\psi\mathcal D\psi\,
 \exp\left[-\bar\psi D[U]\psi+\bar\eta\psi+\bar\psi\eta\right]
 \nonumber\\
 &=\det D[U]\,\exp\left(\bar\eta D^{-1}[U]\eta\right),
 \qquad S[U]=D^{-1}[U].
\end{align}
这里要求 \(D\) 在所选边界、质量和规范背景下可逆；零模或接近零模时，
数值求逆和统计估计都要另作处理。对 Grassmann 源求导得到 Wick 定理：
\begin{align}
 \langle\psi_{i_1}\cdots\psi_{i_n}
 \bar\psi_{j_1}\cdots\bar\psi_{j_n}\rangle_U
 &=\sum_{\pi\in S_n}(-1)^{\pi}
 \prod_{k=1}^{n}S_{i_kj_{\pi(k)}}[U],
 \\
 \langle\psi_i\bar\psi_j\rangle_U&=S_{ij}[U].
\end{align}
费米负号不是经验系数，而是把 Grassmann 场重新排序到成对收缩所需的
置换奇偶性。这是自动 Wick 引擎可以枚举而不能随意省略的第一性原理。
它依赖 Grassmann 反对易公理和高斯测度；相互作用通过 \(\det D\) 和规范
组态平均进入，而不是被上述有限维代数单独消除。

若费米子离散化满足
\begin{equation}
 D^\dagger=\gamma_5D\gamma_5,\qquad
 S^\dagger=\gamma_5S\gamma_5,
\end{equation}
则反向时间/源汇顺序的传播子可以从已有传播子重建。PyQCD 的
\texttt{seq\_peram} 实际执行
\begin{equation}
 \tau_{\rm seq}(t_{\rm sink},t_{\rm src})
 =\gamma_5\,\tau(t_{\rm src},t_{\rm sink})^\dagger\,\gamma_5,
 \label{eq:seq-peram-clean}
\end{equation}
其中厄米共轭同时交换 Dirac 和蒸馏指标。这个等式的前提是所用 \(D\)、
边界和 gamma 约定确实满足 \(\gamma_5\)-Hermiticity；代码接口本身不能
替调用者证明该前提。
\src{../PyQCD/pyqcd/contraction/_seqperam.py:20--59；refer/books/INTRODUCTION_TO_LATTICE_QCD_latex/chapters/sec17_correlators.tex}

\section{算符、味平衡和自动收缩}
一个双线性算符的典型形式为
\begin{equation}
 \mathcal O_M(x)=\bar q_a^f(x)\,\Gamma_{\alpha\beta}\,
 q_b^g(x)\,\mathcal C_{ab}^{fg},
\end{equation}
重子算符则以 \(\epsilon_{abc}\) 闭合颜色：
\begin{equation}
 \mathcal O_B(x)=\epsilon_{abc}(q_1^a q_2^b)q_3^c,
\end{equation}
并由 gamma 链、味结构和自旋投影决定具体态。PyQCD 的 DSL 用
\texttt{|} 分隔不同时间/粒子区间，用 \texttt{gamma\_5}、
\texttt{gamma\_7}、\texttt{gamma\_mu} 等字符串表示 Dirac 插入，用
\texttt{u}、\texttt{d} 和 \texttt{u\^d} 等标记夸克及反夸克。分隔符
不是物理场，而是解析器识别 VDV/VVV 区间和时间标签的元数据。

\begin{table}[htbp]
\centering
\caption{自动 Wick 收缩的单列算法：从算符 DSL 到可执行 contraction plan}
\small
\begin{tabular}{@{}P{0.94\textwidth}@{}}
\toprule
\(\text{Input: }\texttt{sink\_operators},\;\texttt{source\_operators},\;\texttt{curr\_operators},\;Cpt\in\{\text{bubble},2\text{pt},3\text{pt},4\text{pt}\}\)\\
\(\text{Normalize: }\text{若首尾没有 }\texttt{|}\text{，自动补上；2pt 情形丢弃 current 区间}\)\\
\(\text{Parse: }\text{把味标记分为 }q\text{ 与 }q^d，\text{把 gamma 和分隔符位置分别记录}\)\\
\(\text{Expand wildcard: }q\to\{u,d,s,c,b,t\},\quad l\to\{u,d\},\quad\text{逐组合代入}\)\\
\(\text{Balance test: }N_q(f)=N_{\bar q}(f)\text{；不平衡的味组合直接跳过}\)\\
\(\text{Label fields: }(q,\bar q)\text{ 按出现次序分配 }ab,cd,ef,\ldots\text{ 收缩标签}\)\\
\(\text{Enumerate: }\text{对每种味分别生成夸克到反夸克的所有排列，再做不同味匹配的笛卡尔积}\)\\
\(\text{Fermi sign: }\mathcal s=(-1)^{N_{\rm inv}},\quad N_{\rm inv}=\#\{(i,j):i<j,\,\ell_i>\ell_j\}\)\\
\(\text{Time parse: }|\cdots|\text{ 内含两个夸克给 VDV，含三个夸克给 VVV；区间标成 }t_{\rm sink},t_{\rm cur},t_{\rm src}\)\\
\(\text{Gamma parse: }\text{每个 gamma 读取相邻夸克标签，生成 gamma registry 的输入索引}\)\\
\(\text{Peram entry: }S_{q\bar q}[t_q,t_{\bar q}]\text{ 保存夸克位、反夸克位、味和 Dirac/ev 标签}\)\\
\(\text{Vertex entry: }\text{两夸克区间绑定 }V_{mn}\text{，三夸克区间绑定 }V_{mnl}\text{，并写入时间标签}\)\\
\(\text{Index projection: }\text{重复小写字母求和，自由小写字母保留；大写动量/链接标签保留一次}\)\\
\(\text{Diagram output: }\texttt{result\_indx},\texttt{result\_name},\texttt{result\_sign},\texttt{peram},\texttt{V},\texttt{gamma\_pos}\)\\
\(\text{For each diagram: }\mathcal C_D=\mathcal s_D\,\operatorname{einsum}(V_{\rm sink},S_1,\ldots,S_n,\Gamma,V_{\rm src})\)\\
\(\text{2pt: }\mathcal C_2(t)=\sum_D\mathcal C_D(t)\)\\
\(\text{3pt: }\mathcal C_3(t,\tau)=\sum_D\mathcal C_D(t,\tau;\Gamma_{\rm cur})\)\\
\(\text{4pt/bubble: }\text{保留所有 current/区间连通分支，并按 registry 解析每个中间量}\)\\
\(\text{Check: }\text{味平衡、颜色闭合、输出自由指标与预期轴序一致；否则停止而不猜测}\)\\
\(\text{Output: }\text{一个字典（无 wildcard）或字典列表（有 wildcard），供静态绘图和动态 contraction 使用}\)\\
\bottomrule
\end{tabular}
\end{table}
\src{../PyQCD/pyqcd/contraction/_autowick.py:39--59,208--445；../PyQCD/pyqcd/contraction/_dynamic.py:262--647；../PyQCD/docs/格点QCD中的维克收缩与傅里叶变换.tex}

该算法的“枚举所有排列”在 \(n_f\) 个相同味夸克/反夸克时规模含有
\(n_f!\) 因子；这是物理拓扑的完整性与计算复杂度之间的直接冲突。实际
高阶多强子问题应先利用同味交换对称性、颜色张量和等价图合并，再由
\texttt{dynamic\_contraction} 根据 einsum 路径和设备估算执行顺序。合并
必须保持相对费米号，不能只按字符串去重。
\fi

\chapter{对称性、重整化与极限层次：从局域算符到非局域分布}
\label{sec:sym-renorm}

\section{欧氏对称性、离散对称与算符分类}
格点 QCD 的连续极限由欧氏 \(O(4)\) 对称控制，而有限晶格间距只保留
超立方群 \(H(4)\)；因此，连续洛伦兹表示下本来互不混合的张量分量，在
有限 \(a\) 时可能落入同一 \(H(4)\) 表示并发生算符混合。对局域双线性，
最基本的基底可写为
\begin{equation}
 O_S=\bar q q,\quad
 O_P=\bar q i\gamma_5 q,\quad
 O_V^\mu=\bar q\gamma^\mu q,\quad
 O_A^\mu=\bar q\gamma^\mu\gamma_5 q,\quad
 O_T^{\mu\nu}=\bar q\sigma^{\mu\nu}q,
\end{equation}
以及胶子局域算符
\begin{equation}
 O_{G^2}=F_{\mu\nu}^aF^{a\,\mu\nu},\qquad
 O_{G\widetilde G}=F_{\mu\nu}^a\widetilde F^{a\,\mu\nu},
 \qquad
 \widetilde F^{\mu\nu}=\frac12\epsilon^{\mu\nu\rho\sigma}F_{\rho\sigma}.
\end{equation}
这些算符的物理内容分别对应标量凝聚、赝标量拓扑、守恒流、轴流、
张量荷、胶子能动量密度和拓扑密度。对高阶 twist-2 算符，则需要对
多个协变导数做对称无迹化，例如
\begin{equation}
 O_{\{\mu_1\cdots\mu_n\}}
 =\bar q\,\gamma_{\{\mu_1} iD_{\mu_2}\cdots iD_{\mu_n\}}\,q
 -\text{traces},
\end{equation}
其矩阵元直接给出部分子分布的矩。有限 \(a\) 时的关键不是“有没有这个
算符”，而是“它属于哪一个表示、和谁混合、需要怎样的减迹与重整化”。
\src{refer/books/INTRODUCTION_TO_LATTICE_QCD_latex/chapters/sec06_formulation_gauge.tex；refer/books/INTRODUCTION_TO_LATTICE_QCD_latex/chapters/sec17_correlators.tex；refer/books/INTRODUCTION_TO_LATTICE_QCD_latex/chapters/sec18_hadron_spectrum.tex}

\section{Wilson 线、线性发散与乘法重整化}
非局域算符的核心对象是路径有序 Wilson 线
\begin{equation}
 W(0,z)=\mathcal P\exp\!\left(ig\int_0^1 d\lambda\, z^\mu A_\mu(\lambda z)\right),
\end{equation}
它保证双线性
\begin{equation}
 \mathcal O_\Gamma(z)=\bar q(0)\,\Gamma\,W(0,z)\,q(z)
\end{equation}
在规范变换下仍是规范不变量。对格点上带直 Wilson 线的裸矩阵元，最显著
的 UV 结构是长度成正比的线性发散，通常可参数化为
\begin{equation}
 \mathcal O_\Gamma^{\rm bare}(z,a)
 =e^{-\delta m(a)|z|}\,Z_\Gamma(z,a,\mu)\,
 \mathcal O_\Gamma^{\rm ren}(z,\mu)
 +\sum_{\Gamma'} Z_{\Gamma\Gamma'}^{\rm mix}(z,a,\mu)\,
 \mathcal O_{\Gamma'}^{\rm ren}(z,\mu).
\end{equation}
这里 \(\delta m\) 物理上对应 Wilson 线自能，和静态色源的质量重整化同源；
而混合矩阵 \(Z_{\Gamma\Gamma'}\) 则编码了晶格对称性降低后允许的算符混合。
若路径不是直线而是 staple，则还会出现 cusp 发散和端点发散，重整化因子
要额外带上 cusp 角与 rapidity 方案依赖。换句话说，Wilson 线不是“装饰性
的连线”，而是把 gauge invariant 非局域算符与其 UV 结构绑在一起的主体。
\src{refer/papers/General_Method_Nonpert_Renorm_latex/；refer/papers/Nonpert_Renorm_Quark_Bilinears_latex/；refer/papers/Complete_NP_Renorm_QuasiPDF_latex/；refer/papers/Properties_Uses_Wilson_Flow_latex/}

\section{小距离 OPE、LaMET 与小 \(b_T\) 展开}
在局域 OPE 中，双线性在小距离极限可展开为一串局域算符：
\begin{equation}
 \mathcal O_\Gamma(z,\mu)
 =\sum_{n=0}^{\infty}
 \frac{z^{\mu_1}\cdots z^{\mu_n}}{n!}\,
 C_n(z^2\mu^2,\alpha_s)\,
 O_{\mu_1\cdots\mu_n}(\mu)
 +O(z^2\Lambda_{\rm QCD}^2),
\end{equation}
其中 \(C_n\) 是短距 Wilson 系数，\(O_{\mu_1\cdots\mu_n}\) 是 twist-2 局域算符。
这就是为什么 \(z\to 0\) 的极限给出矩而不是直接给出全 \(x\) 分布。
对准分布，LaMET 因子化写成
\begin{equation}
 \widetilde f(x,P_z,\mu)
 =\int_{-1}^{1}\frac{dy}{|y|}\,
 C\!\left(\frac{x}{y},\frac{\mu}{|y|P_z}\right)\,
 f(y,\mu)
 +O\!\left(\frac{\Lambda_{\rm QCD}^2}{P_z^2}\right),
\end{equation}
说明有限 \(P_z\) 下的准分布是光锥分布的可匹配代理，而不是同一对象。
对 TMD，则在小横向距离处有
\begin{equation}
 F(x,b_T;\mu,\zeta)
 =\sum_j\int_x^1 \frac{d\xi}{\xi}\,
 C_{j}\!\left(\frac{x}{\xi},b_T;\mu,\zeta\right)
 f_j(\xi,\mu)
 +O(b_T^2\Lambda_{\rm QCD}^2).
\end{equation}
因此 \(z\to0\)、\(b_T\to0\) 和 \(P_z\to\infty\) 是三条不同极限；它们共享
短距/大动量的共同直觉，但对应的因子化结构、匹配核和误差账本并不相同。
\src{refer/papers/LaMET_RMP_Review_latex/；refer/papers/Quasi_Pseudo_PDF_Radyushkin_latex/；refer/papers/One_Loop_Matching_Nonsinglet_latex/；refer/papers/Locally_Smeared_OPE_latex/}

\section{Collins--Soper 演化、软因子与过程依赖}
对已减软的 TMD 分布，rapidity 演化由 Collins--Soper 方程控制：
\begin{equation}
 \frac{\partial \ln F(x,b_T;\mu,\zeta)}{\partial \ln\sqrt{\zeta}}
 =K(b_T,\mu),
 \qquad
 \mu\frac{d}{d\mu}K(b_T,\mu)
 =-\gamma_K(\alpha_s)=-2\Gamma_{\rm cusp}(\alpha_s).
\end{equation}
这里 \(K\) 是 Collins--Soper 核，\(\Gamma_{\rm cusp}\) 是 cusp anomalous
 dimension。软因子减除的物理作用，是把双线性与 Wilson 线自能中与硬过程
无关的长距离 rapidity 效应剥离出去，使最后的物理截面只依赖可匹配的
非扰动 TMD 和硬核。若过程改变，staple gauge link 的未来向/过去向也会
改变，T-odd 成分的符号关系随之变化；这不是代码约定，而是量子场论中的
路径与因果结构。
\src{refer/papers/TMD_Soft_Function_LaMET_latex/；refer/papers/从准TMD波函数确定CS核_latex/；refer/papers/LaMET计算TMD软函数_latex/}

\section{有限体积、离散误差与连续外推}
真实格点结果还要同时面对有限体积、非零晶格间距和有限动量三种系统误差。
一个常见的外推记号可写成
\begin{equation}
 Q(a,L,P_z)
 =Q_{\rm cont}
 +c_1 a^{\nu}
 +c_2 a^{\nu+1}
 +d_1 e^{-m_\pi L}
 +\frac{e_1}{P_z^2}
 +\cdots,
\end{equation}
其中 \(\nu\) 取决于所用作用量与改进程度：未改进 Wilson 型离散化常见
\(O(a)\) 项，而对称性更好的改进作用量可把 leading artifact 推到 \(O(a^2)\)。
对非局域算符，还需要额外满足物理窗口，例如小距离 OPE 要求
\(a\ll |z| \ll \Lambda_{\rm QCD}^{-1}\)，小横向距离 OPE 要求
\(a\ll b_T \ll \Lambda_{\rm QCD}^{-1}\)，梯度流原型要求
\(a^2\ll t\ll \Lambda_{\rm QCD}^{-2}\)。这些窗口一旦塌缩，剩下的就不是“可
外推的系统误差”，而是定义本身已经换了。
\begin{table}[htbp]
\centering
\caption{主要物理极限与它们真正控制的对象}
\small
\begin{tabularx}{0.96\textwidth}{@{}P{2.0cm}P{3.5cm}P{3.8cm}Y@{}}
\toprule
极限 & 物理对象 & 需要满足 & 不能替代的内容\\
\midrule
\(a\to0\) & 离散化误差 & 统一重整化方案、同一物理尺度 & 不能把不同改进级别直接相加。\\
\(L\to\infty\) & 有限体积效应 & \(m_\pi L\gg 4\) 且态质量已知 & 不能用 smearing 代替体积外推。\\
\(z\to0\) & 局域 OPE & \(a\ll |z|\ll\Lambda^{-1}\) & 不能把裸 bilocal 直接当局域算符。\\
\(b_T\to0\) & small-\(b_T\) OPE & 软因子与 rapidity 方案已定 & 不能把 TMD 直接改写成 collinear PDF。\\
\(P_z\to\infty\) & LaMET 匹配 & 多个动量点、幂修正可控 & 不能用单一 \(P_z\) 冒充光锥极限。\\
\(t\to0^+\) & flow window & \(a^2\ll t\ll\Lambda^{-2}\) & 不能把流时间当作自动重整化。\\
\bottomrule
\end{tabularx}
\end{table}
从物理上看，这一整套极限层次形成的是“定义—匹配—外推”链条：定义决定
算符几何与对称性，匹配决定不同方案之间的短距连接，外推决定能否回到
连续 QCD。PyQCD 中的实现只是在每一层给出一个可运行的有限维代理，用于
检查这些极限是否真的闭合，而不是替代它们。
\src{refer/books/INTRODUCTION_TO_LATTICE_QCD_latex/chapters/sec06_formulation_gauge.tex；refer/books/INTRODUCTION_TO_LATTICE_QCD_latex/chapters/sec17_correlators.tex；refer/books/INTRODUCTION_TO_LATTICE_QCD_latex/chapters/sec18_hadron_spectrum.tex；refer/papers/Quasi_PDF_Gradient_Flow_latex/；refer/papers/威尔逊流的性质与应用_latex/}

\chapter{接口映射、统一链路与方案边界}
\label{sec:interface-chain}

前文已经把局域双线性、非局域 Wilson 线、谱学提取与 TMD 重整化分开叙述。
这里再把它们压成一张物理接口图：什么输入对应什么对象，哪一步只是压缩
自由度，哪一步真正改变了方案，哪一步只能算匹配。这样才能区分“实现细节”
和“物理定义”的边界。

\section{从函数到物理对象：接口总表}
下面这张表不把代码当结论，而把代码看成物理对象的离散代理。概念名与
真实入口不必逐字一致，但它们的输入、输出和物理责任必须一一对应。

\begin{table}[htbp]
\centering
\caption{PyQCD 中若干关键入口与其物理对象的对应}
\small
\begin{tabularx}{0.98\textwidth}{@{}P{2.4cm}P{3.2cm}P{3.5cm}Y@{}}
\toprule
入口/概念 & 主要输入 & 中间量 & 输出与物理含义\\
\midrule
HYP/Stout smearing & 原始链接 \(U_\mu(x)\)、涂抹参数 & 迭代投影后的平滑链接 & 只过滤 UV 噪声，不能替代连续极限。\\
\codeid{momsmear_phase} / momentum smearing & 源、目标动量 \(\vec P\)、相位因子 & 动量增强的源向量 & 改善 boosted hadron 的重叠因子。\\
eigvec / distillation & 拉普拉斯本征模数 \(N_{\rm ev}\) & \(V\)、\(P=VV^\dagger\)、perambulator \(\tau\) & 低秩投影子空间，便于收缩和谱学。\\
\texttt{VdV}, \texttt{VVV} & 本征模与链接相位 & 蒸馏顶点张量 & 把空间几何压缩成有限维顶点。\\
2pt & 插值算符、动量投影、时距 & \(C_2(t,\vec P)\)、\(m_{\rm eff}\) & 由谱分解提取能量和基态占比。\\
3pt & 两点函数、插入算符、源汇时序 & \(C_3(t,\tau)\) & 用于矩阵元与比值提取。\\
ope / staple operator & \(F_{\mu\nu}\)、Wilson 线、方向 \(\pm z\)、\(b_T\) & 规范不变双局域算符 & 生成 quasi/PDF/TMD 所需的非局域对象。\\
ratio & \(C_3/C_2\)、投影与真空扣除 & 平台区比值 \(R\) & 消去重叠因子，隔离裸矩阵元。\\
barematrix & 逐窗拟合、协方差 & \(c_0\) & 裸矩阵元或其统计样本。\\
renorm & \(Z_R\)、RI/MOM、hybrid、flow & 重整化矩阵元 & 消除线性/对数发散并定义方案。\\
matching & \(h_R\)、\(F\)、\(P_z\)、\(\mu\)、\(\zeta\) & 卷积核 \(C\) & 把欧氏格点对象接到光锥分布。\\
\bottomrule
\end{tabularx}
\end{table}
\src{../PyQCD/pyqcd/AGENTS.md；../PyQCD/pyqcd/pipeline/_runner.py；../PyQCD/pyqcd/pipeline/_steps.py；../PyQCD/pyqcd/vertex/AGENTS.md；../PyQCD/pyqcd/analysis/AGENTS.md；../PyQCD/pyqcd/operator/AGENTS.md；../PyQCD/pyqcd/renorm/AGENTS.md}

\section{统一算法链路：从组态到可匹配分布}
把整个物理链写成“一个长算法”时，最重要的是记住每一步只回答一个问题：
是否把某个自由度压缩掉，或把某个方案变量改写掉，而不是一步把所有问题都
解决掉。总链条可概括为
\begin{align}
 U_\mu(x)&\xrightarrow{\text{UV 预处理}}U_\mu^{\rm sm}(x)
 \xrightarrow{\text{eigen/distillation}}V,\tau
 \xrightarrow{\text{vertex}}VdV,VVV \nonumber\\
 &\xrightarrow{C_2,C_3,\mathcal O}R,c_0
 \xrightarrow{Z_R/\mathrm{RI\text{-}MOM}/\mathrm{hybrid}}h_R
 \xrightarrow{\text{Fourier / matching}}f_{\rm LC}.
 \label{eq:unified-theory-chain}
\end{align}

\begin{table}[htbp]
\centering
\caption{统一算法链路：每一步对应的物理门与可停条件}
\small
\begin{tabularx}{0.98\textwidth}{@{}P{0.7cm}P{3.1cm}Y P{3.2cm}@{}}
\toprule
步 & 输入 & 更新 / 输出 & 停止条件与物理解释\\
\midrule
0 & \(U_\mu(x),a,L,m_q\)；\(P_z,z,b_T,\mu,\zeta\) & 识别目标属于 2pt/3pt、准 PDF、伪 PDF 还是 TMD。 & 若 \(m_\pi L\) 太小、\(a\) 太粗或窗口不闭合，则只停在原型。\\
1 & 原始规范场 & 可选 HYP/Stout/flow 平滑。 & 只在目标算符定义允许时使用；不能把 smearing 代替重整化。\\
2 & 平滑后的 \(U_\mu\) & 拉普拉斯本征模、\(P=VV^\dagger\)、perambulator。 & \(V^\dagger V=\mathbf 1\) 必须成立；否则蒸馏子空间失效。\\
3 & \(V\)、links、phases & \(VdV\)、\(VVV\)。 & 顶点张量必须满足颜色和动量闭合。\\
4 & 顶点、传播子、源汇几何 & \(C_2\)、\(C_3\)、disconnected loops、OPE staple。 & 若味指标或自由 Lorentz 指标未闭合，则停止。\\
5 & \(C_2, C_3\) & 比值 \(R\)、\(m_{\rm eff}\)、plateau。 & 平台应对拟合窗稳定；否则只能报告窗口依赖。\\
6 & \(R\) 或 disconnected ratio & 裸矩阵元 \(c_0\) 与系统误差条目。 & 若激发态污染未压低，则不能升级为物理矩阵元。\\
7 & \(c_0\)、scheme choice & \(Z_R\)、RI/MOM、hybrid 或 flow 结果。 & 方案必须与几何、端点和软减除一致。\\
8 & \(h_R(z)\) 或 \(\widetilde F(x,b_T)\) & Fourier / small-\(z\) / small-\(b_T\) / LaMET 匹配。 & 必须保留卷积核和幂修正，不可直接等同光锥分布。\\
9 & 多个 \(a,L,P_z,t\) 点 & 连续、无限体积、无限动量极限。 & 至少两个方向的外推参数必须可约束；单点只能算示意。\\
10 & 全部门通过 & 物理结果 \(\in\) `verified`. & 若任何门未闭合，结果只到 `structural` 或 `inferred`.\\
\bottomrule
\end{tabularx}
\end{table}

这条链的关键在于：每一层都只压缩一个自由度。smearing 只压 UV，distillation
只压空间低模，ratio 只压重叠因子，\(Z_R\)/RI/MOM/hybrid 只压 UV 发散，
matching 只把短距系数重新分配到光锥分布；若某一步同时声称解决了所有问题，
通常就是把定义、重整化和外推混成了一步。
\src{../PyQCD/pyqcd/pipeline/_runner.py；../PyQCD/pyqcd/pipeline/_steps.py；../PyQCD/pyqcd/analysis/_analyse.py；../PyQCD/pyqcd/analysis/_ratio2pt.py；../PyQCD/pyqcd/analysis/_disconnected.py；../PyQCD/pyqcd/renorm/AGENTS.md}

\section{重整化方案比较：ratio、\(Z_R\)、RI/MOM、hybrid 与 flow}
Wilson 线的线性发散不是“格点上多出来的一点噪声”，而是线状静色源自能
带来的线张力项，通常可写成
\begin{equation}
 \mathcal O_\Gamma^{\rm bare}(z,a)
 =\exp[-\delta m(a)|z|]\,
 Z_{\Gamma}^{\log}(z,\mu)\,
 \mathcal O_\Gamma^{\rm ren}(z,\mu)
 +\sum_{\Gamma'} Z_{\Gamma\Gamma'}^{\rm mix}(z,a,\mu)\,
 \mathcal O_{\Gamma'}^{\rm ren}(z,\mu).
\end{equation}
若路径是 staple，则还会多出 cusp、端点和 rapidity 相关的重整化因子。
因此，重整化不是单一乘法，而是“线性发散 + 对数发散 + 可能的混合 + 方案
定义”四件事的联合账本。

\begin{table}[htbp]
\centering
\caption{常用重整化/减除方案的物理职责}
\small
\begin{tabularx}{0.98\textwidth}{@{}P{1.8cm}P{2.6cm}P{2.7cm}P{3.1cm}Y@{}}
\toprule
方案 & 主要减除对象 & 典型输入 & 优势 & 主要局限\\
\midrule
ratio & 共同重叠因子与部分短距噪声 & 同几何 2pt/3pt 或 \(z=0\) 参考点 & 简单、稳定、短距鲁棒。 & 长距离和方案依赖不能靠它单独解决。\\
\(Z_R\) self-renorm & Wilson 线自能与参考长度归一 & 裸 bilocal、参考点或零动量比值 & 直接针对线性发散，适合长 \(z\) 轨道。 & 需要清楚的 renorm window 和参考定义。\\
RI/MOM & 指定 off-shell Green 函数的投影条件 & Landau 规约后的外部夸克态 & 方案定义明确，便于与连续方案对接。 & 对含 Wilson 线的双局域算符易受混合与窗口污染。\\
hybrid & 短距用 ratio，长距用显式 \(Z_R\) & 短长距离拼接点 \(z_s\) & 能利用短距稳定区和长距显式减除各自优点。 & 需要在 \(z_s\) 附近做连续拼接。\\
flow & 流时间平滑后的 UV 模式 & \(t>0\) 的 flowed 场或 flowed 算符 & 局域原型最清楚，数值噪声较低。 & 必须满足 \(a^2\ll t\ll\Lambda_{\rm QCD}^{-2}\)。\\
soft subtraction & soft factor 与 rapidity 发散 & non-local staple / TMD 几何 & 这是 TMD 物理不可缺的减除。 & 不能拿来代替 Wilson 线自能重整化。\\
\bottomrule
\end{tabularx}
\end{table}

一个实用但不失真的写法是把拼接方案写成分段函数：
\begin{equation}
 h_R^{\rm hybrid}(z,\mu)=
 \theta(z_s-|z|)\,\frac{h_{\rm bare}(z,P_z)}{h_{\rm bare}(z,0)}
 +\theta(|z|-z_s)\,\eta_s\frac{h_{\rm bare}(z,P_z)}{Z_R(z,\mu)},
 \qquad \eta_s=\frac{Z_R(z_s,\mu)}{h_{\rm bare}(z_s,0)},
\end{equation}
并在 \(z_s\) 附近用光滑插值保证连续性。这个表达式的要点不是“公式长得
漂亮”，而是把短距离的统计稳定区和长距离的显式减除区分开。RI/MOM 的
定义则可写成投影条件
\begin{equation}
 Z_q^{-1} Z_\Gamma^{\rm RI/MOM}(z,\mu,p^2)\,
 \Tr\!\left[P\,\Lambda_\Gamma(p,z)\right]_{p^2=\mu^2}
 =\Tr\!\left[P\,\Lambda_\Gamma^{\rm tree}(p,z)\right],
\end{equation}
它的物理意义是“在选定的离壳点上固定方案”，而不是把所有双局域算符都
变成局域算符。
\src{refer/papers/General_Method_Nonpert_Renorm_latex/；refer/papers/Nonpert_Renorm_Quark_Bilinears_latex/；refer/papers/Complete_NP_Renorm_QuasiPDF_latex/；refer/papers/威尔逊线重正化改进准分布_latex/；refer/papers/准光前关联的混合重正化方案_latex/；refer/papers/Quasi_PDF_Gradient_Flow_latex/；refer/papers/TMD_Soft_Function_LaMET_latex/}

\section{准 PDF、伪 PDF 与准 TMD 的短距语言}
这三类对象共享“小距离/大动量”的直觉，但它们的坐标与卷积变量并不相同。
准 PDF 更偏向固定 \(P_z\) 下的空间 Fourier 变换；伪 PDF 更偏向 Ioffe 时间
\(\nu=zP_z\) 的坐标空间表述；TMD 则保留横向分离 \(b_T\) 并必须做软因子减除。
\begin{equation}
 \mathcal M(\nu,z^2)
 =\int_{-1}^{1}\! \dd x\, \e^{\ii x\nu}\,\mathcal P(x,z^2),
 \qquad
 \widetilde q(x,P_z,\mu)
 =\int\frac{\dd z}{2\pi}\,\e^{-\ii x z P_z}\,\widetilde h(z,P_z,\mu).
\end{equation}
局域 OPE 给出的是 Ioffe 时间或短距展开：
\begin{equation}
 \mathcal M(\nu,z^2)
 =\sum_{n=0}^{\infty}\frac{(\ii\nu)^n}{n!}\,
 C_n(z^2\mu^2)\,a_n(\mu)
 +O(z^2\Lambda_{\rm QCD}^2),
\end{equation}
其中 \(a_n\) 是局域 twist-2 矩，\(C_n\) 是短距 Wilson 系数。
TMD 则保留横向距离并在软减除后满足
\begin{equation}
 F(x,b_T;\mu,\zeta)
 =\frac{Z_F(\mu,\zeta)}{\sqrt{S(b_T;\mu,\delta)}}\,
 F_{\rm unsub}(x,b_T;\delta)
 =\sum_j C_j(b_T,\mu,\zeta)\otimes f_j(x,\mu)
 +O(b_T^2\Lambda_{\rm QCD}^2).
\end{equation}
三者的差别可以概括为：
\begin{itemize}
  \item \(z\to0\) 是局域 OPE 的极限，目标是矩而不是完整分布；
  \item \(P_z\to\infty\) 是 LaMET 的极限，目标是把准 PDF 匹配到光锥 PDF；
  \item \(b_T\to0\) 是 TMD 的短横向距离极限，目标是把 TMD 匹配到 collinear PDF。
\end{itemize}
因此，不能把“都用了 Fourier”误写成“它们是同一个对象的不同记号”。
\src{refer/papers/LaMET_RMP_Review_latex/；refer/papers/Quasi_Pseudo_PDF_Radyushkin_latex/；refer/papers/Locally_Smeared_OPE_latex/；refer/papers/TMD_Soft_Function_LaMET_latex/；refer/papers/从准TMD波函数确定CS核_latex/；refer/papers/One_Loop_Matching_Nonsinglet_latex/}

\section{本章的最小验证门}
本章多出的内容并没有引入新的可执行结论，而是把已有结论的边界写得更硬。
最小验证门仍然是：
\begin{itemize}
  \item \(V^\dagger V=\mathbf 1\) 与 \(P^2=P\)：蒸馏投影必须闭合；
  \item \(C_2\) 与 \(C_3\) 的平台区在拟合窗内稳定：比值才可视为裸矩阵元代理；
  \item \(Z_R\)、RI/MOM、hybrid 三者在重叠窗口内应给出相容的短距极限；
  \item \(a\ll |z|, b_T\) 与 \(a^2\ll t\) 必须同时成立：否则只能报告原型，不可报告物理结果；
  \item TMD 必须保留 soft subtraction 与 Collins--Soper 演化：少一个就不是完整 TMD。
\end{itemize}
从这个意义上说，物理图像、公式链和代码接口三者必须同时闭合，报告才算
真正把“能算”与“可解释”连到了一起。
\src{../PyQCD/pyqcd/vertex/AGENTS.md；../PyQCD/pyqcd/analysis/AGENTS.md；../PyQCD/pyqcd/renorm/AGENTS.md；../PyQCD/pyqcd/operator/AGENTS.md；../PyQCD/pyqcd/AGENTS.md}

\section{从生成泛函、传递矩阵到分布函数的统一推导}
格点 QCD 的整条物理链条，最核心的不是某个接口函数，而是 Euclidean 生成
泛函、传递矩阵谱分解、局域/非局域算符插入、重整化和因子化这五个层次的
连续衔接。起点是带源路径积分
\begin{equation}
 Z[J,\eta,\bar\eta]
 =\int\mathcal D U\,\mathcal D\bar\psi\,\mathcal D\psi\,
 \exp\!\left[-S_E[U,\bar\psi,\psi]+J\!\cdot\!A+\bar\eta\psi+\bar\psi\eta\right],
\end{equation}
它的物理意义只有一个：所有相关函数都是对源的变分导数。对格点理论，
Euclidean 时间演化由传递矩阵 \(T=\e^{-aH}\) 代替连续 \(e^{-Ht}\)，于是
两点函数自动具有谱分解
\begin{equation}
 C_2(t,\vec p)
 = \sum_{\vec x}e^{-i\vec p\cdot\vec x}
 \langle 0|O(t,\vec x)O^\dagger(0)|0\rangle
 =\sum_n \frac{|Z_n(\vec p)|^2}{2E_n(\vec p)}e^{-E_n(\vec p)t}.
\end{equation}
这里 \(Z_n=\langle 0|\widehat O|n\rangle\) 是插值算符与本征态的重叠因子，
而 \(2E_n\) 来自相对论态归一。长时间极限下最低能态支配，因此有效质量
和平台区不是经验巧合，而是谱分解的必然后果。三点函数在插入算符
\(\mathcal O\) 后满足
\begin{align}
 C_3(t_{\rm sep},\tau)
 &=\sum_{m,n}\frac{Z_m^f Z_n^{i*}}{4E_m^fE_n^i}
 e^{-E_m^f(t_{\rm sep}-\tau)}e^{-E_n^i\tau}
 \langle m,f|\mathcal O|n,i\rangle,\\
 R(t_{\rm sep},\tau)
 &\longrightarrow \frac{\langle f|\mathcal O|i\rangle}{2\sqrt{E_fE_i}}
 \qquad (0\ll\tau\ll t_{\rm sep}).
\end{align}
所以 ratio 方案并不是“取个比值”，而是利用相同谱因子在二点/三点中
的共同出现，将激发态污染、插值算符重叠和部分指数项消去，留下带相对论
归一的矩阵元。

若插入的是非局域 Wilson 线双线性，则物理对象改写为
\begin{equation}
 \mathcal O_\Gamma(z)=\bar q(0)\Gamma W(0,z)q(z),\qquad
 W(0,z)=\mathcal P\exp\!\left[ig\int_0^1\!d\lambda\,z^\mu A_\mu(\lambda z)\right].
\end{equation}
Wilson 线不是附加装饰，而是把分离端点重新缝合成规范不变量的主体。
在短距离极限下它满足局域 OPE：
\begin{equation}
 \mathcal O_\Gamma(z)
 =\sum_{n=0}^{\infty}\frac{z^{\mu_1}\cdots z^{\mu_n}}{n!}
 \bar q(0)\Gamma (iD_{\mu_1})\cdots(iD_{\mu_n})q(0)
 +O(z^2\Lambda_{\rm QCD}^2),
\end{equation}
因此 \(z\to0\) 给出的是局域矩，而不是完整分布。若再通过固定
空间动量 \(P_z\) 做 Fourier 变换，就得到准 PDF；这一步的正确物理内容是
LaMET 因子化：
\begin{equation}
 \widetilde f(x,P_z,\mu)
 =\int_{-1}^{1}\frac{dy}{|y|}\,
 C\!\left(\frac{x}{y},\frac{\mu}{|y|P_z}\right)f(y,\mu)
 +O\!\left(\frac{\Lambda_{\rm QCD}^2}{P_z^2}\right).
\end{equation}
所以准 PDF 与光锥 PDF 不是同一个对象的两种写法，而是通过短距匹配核
连接的不同定义。对伪 PDF，则主变量是 Ioffe 时间 \(\nu=zP_z\)：
\begin{equation}
 \mathcal M(\nu,z^2)=\int_{-1}^{1}\dd x\,e^{ix\nu}\,\mathcal P(x,z^2)
 =\sum_{n=0}^{\infty}\frac{(i\nu)^n}{n!}C_n(z^2\mu^2)a_n(\mu)+O(z^2\Lambda_{\rm QCD}^2),
\end{equation}
它的物理内容是“短距矩展开 + 非局域 Fourier 表述”的统一。

TMD 再向上保留横向距离 \(b_T\)，并必须做 soft subtraction：
\begin{equation}
 F(x,b_T;\mu,\zeta)
 =\frac{Z_F(\mu,\zeta)}{\sqrt{S(b_T;\mu,\delta)}}\,F_{\rm unsub}(x,b_T;\delta)
 =\sum_j C_j(b_T,\mu,\zeta)\otimes f_j(x,\mu)
 +O(b_T^2\Lambda_{\rm QCD}^2).
\end{equation}
这里 \(S\) 不是任意归一化常数，而是把 Wilson 线自能、soft 远场和
rapidity 发散单独剥离出来的软因子。它随 Collins--Soper 演化：
\begin{equation}
 \frac{\partial\ln F(x,b_T;\mu,\zeta)}{\partial\ln\sqrt{\zeta}}=K(b_T,\mu),
 \qquad
 \mu\frac{d}{d\mu}K(b_T,\mu)=-\gamma_K(\alpha_s)
 =-2\Gamma_{\rm cusp}(\alpha_s).
\end{equation}
这条方程说明：TMD 的本质不是“多一个变量”，而是一个额外的快度演化
维度。若过程改为未来向或过去向 staple，T-odd 结构的符号关系也要随之
改变；这不是实现选择，而是因果结构与路径有序性的结果。

从这五个层次合起来看，整条理论链可以概括成
\[
 \text{生成泛函}\to\text{谱分解}\to\text{局域/非局域算符}\to
 \text{重整化}\to\text{匹配与演化}.
\]
SymPy 在这里的作用很有限：它可以验证有限维代理、指数展开、Hermitian
投影和部分极限，但不能证明真实系综上的非微扰闭合。真正的物理结论仍然
要靠同几何、同方案、同数据源、同外推门槛的数值链条完成。

\begin{table}[htbp]
\centering
\caption{从 Euclidean 相关函数到分布函数的主导关系}
\small
\begin{tabularx}{0.98\textwidth}{@{}P{2.2cm}P{3.1cm}P{2.8cm}Y@{}}
\toprule
对象 & 主导公式 & 物理极限 & 失败时的正确表述\\
\midrule
\(C_2\) & \(\sum_n |Z_n|^2 e^{-Et}/(2E)\) & \(t\to\infty\) & 只能报告多态混合或平台不足。\\
\(C_3\) & 双指数谱分解 & \(0\ll\tau\ll t\) & 只能报告激发态污染和窗依赖。\\
bilocal & \( \bar q\Gamma W q \) & \(z\to0\) & 只能报告局域矩或裸双局域。\\
quasi PDF & LaMET 卷积 & \(P_z\to\infty\) & 只能报告有限动量代理。\\
pseudo PDF & Ioffe 时间傅里叶 & \(\nu\) 短距展开 & 只能报告坐标空间矩。\\
TMD & soft-subtracted \(F(x,b_T)\) & \(b_T\to0\) 与 CS 演化 & 只能报告 finite-staple 原型。\\
\bottomrule
\end{tabularx}
\end{table}
\src{refer/books/INTRODUCTION_TO_LATTICE_QCD_latex/chapters/sec17_correlators.tex；refer/books/INTRODUCTION_TO_LATTICE_QCD_latex/chapters/sec18_hadron_spectrum.tex；refer/books/Quantum_Chromodynamics_on_the_Lattice_latex/chapters/Chapter05.tex；refer/papers/LaMET_RMP_Review_latex/；refer/papers/TMD_Soft_Function_LaMET_latex/；refer/papers/从准TMD波函数确定CS核_latex/；refer/papers/Complete_NP_Renorm_QuasiPDF_latex/}

\chapter{MyQCD 的完整公式链、SymPy 复现与课程证据}
\label{sec:myqcd-evidence}

\noindent 前一章讲的是对照对象 \fileline{PyQCD} 的物理接口。这里必须把视角
拉回工作对象本身：\fileline{MyQCD} 不是“围着 PyQCD 写说明”，而是把格点
QCD 理论拆成可执行、可审计、可追溯的完整公式链、SymPy 复现和课程模块。
另一个必须明确的负面约束是：按你指出的 \texttt{tag-bug0} 标注，先前的
\texttt{LQCD course PPT} 只是质量反例，不作为正向来源，也不进入后文的
公式链或参考源。

\section{证明链与来源闭环}
\fileline{myqcd/derivations.py} 中的推导函数与本报告正文共同构成 MyQCD 的
理论审计链；\fileline{myqcd/formula_registry.py} 仅是机器审计用的索引，
不是本报告的理论正文，也不能替代这里逐步写出的定义、假设、推导和适用边界。
正文的主线从场论公理、格点离散、谱学提取、统计推断、重整化推进到
LaMET/TMD 匹配；每一层都必须回到对应的参考源和可检验条件。

\begin{table}[htbp]
\centering
\caption{MyQCD 的三层证据}
\small
\begin{tabularx}{0.98\textwidth}{@{}P{2.3cm}P{2.5cm}P{3.1cm}Y@{}}
\toprule
层 & 规模 & 物理重心 & 结论\\
\midrule
\fileline{myqcd/} & 推导函数与审计索引 & 生成泛函、链接变量、涂抹、谱学、重整化、LaMET、TMD、流与匹配。 & 只证明已声明的代数子问题；不等于真实系综或非微扰结果。\\
\fileline{docs/lattice_qcd_gluon_tmd_course/myqcd/} & 辅助课程模块 & 群论/QFT、格点谱学、梯度流、重整化、软因子与 Collins--Soper。 & 只作可运行的局部复现；不替代正文推导和学术来源。\\
\fileline{refer/books/} + \fileline{refer/papers/} & 50 篇论文 + 若干教材 & 原文公式、定理与编号。 & 来源事实源；负责锚定，不负责代替推导。\\
\bottomrule
\end{tabularx}
\end{table}

\subsection{从生成泛函到格点链接}
最底层的连续场论基石是生成泛函
\begin{equation}
 Z[J]=\int\mathcal D\phi\,\exp[-S_E[\phi]+J\phi],
\end{equation}
它在有限维高斯类比里可被精确验证。格点化后，连续规范场变成链接变量
\begin{equation}
 U_\mu(x)=\exp\!\left[iagA_\mu\!\left(x+\frac{a\hat\mu}{2}\right)\right],
\end{equation}
并在 \(a\to0\) 时恢复连续联络。对非阿贝尔情形，这里不是标量指数而是矩阵
指数；可执行检查只负责验证单值性和展开，不替代真正的非微扰动力学。
\src{myqcd/derivations.py；refer/books/INTRODUCTION_TO_LATTICE_QCD_latex/chapters/sec06_formulation_gauge.tex}

\subsection{涂抹、动量平移与蒸馏}
MyQCD 把 UV 过滤、动量增强和低秩投影分成三类不同对象。夸克 Gaussian/Jacobi
涂抹写作
\begin{equation}
 J_{\sigma,n_\sigma}
 =\left(1+\frac{\sigma\nabla^2}{n_\sigma}\right)^{n_\sigma}
 \longrightarrow e^{\sigma\nabla^2},
\end{equation}
其作用是压低局域高频模式，而不是改写物理定义。动量涂抹则对应 Fourier
空间平移
\begin{equation}
 \widetilde f(p)\mapsto \widetilde f(p-k),\qquad f_z\mapsto e^{ikz}f_z.
\end{equation}
HYP/Stout 涂抹对应群值投影与局域更新；在 SU(2) 代理中，stout 更新可写为
\begin{equation}
 Q=\frac{i}{2}(\Omega^\dagger-\Omega)-\frac{i}{2N}\mathrm{Tr}(\Omega^\dagger-\Omega),
 \qquad
 U'=e^{iQ}U.
\end{equation}
蒸馏则用拉普拉斯本征模构造低秩投影
\begin{equation}
 \Box=VV^\dagger,\qquad V^\dagger V=\mathbf 1,\qquad P=VV^\dagger,
\end{equation}
从而把空间传播子压缩到可控子空间，后续的 \(VdV\)、\(VVV\) 顶点与 perambulator
都在这个低秩几何里闭合。
\src{myqcd/derivations.py；docs/lattice_qcd_gluon_tmd_course/myqcd/group_qft.py}

\subsection{谱学提取、协方差与比值}
格点两点函数和三点函数的物理结构是这条链的中段：
\begin{equation}
 C_2(t,\mathbf p)=\sum_n \frac{|Z_n(\mathbf p)|^2}{2E_n(\mathbf p)}e^{-E_n t},
\end{equation}
\begin{equation}
 C_3(t_{\rm sep},\tau)=
 \sum_{m,n}\frac{Z_m^f Z_n^{i*}}{4E_m^fE_n^i}
 e^{-E_m^f(t_{\rm sep}-\tau)}e^{-E_n^i\tau}
 \langle m,f|\mathcal O|n,i\rangle.
\end{equation}
当 \(\tau\) 与 \(t_{\rm sep}-\tau\) 都进入平台区时，裸矩阵元的基态极限为
\begin{equation}
 R(\tau)\longrightarrow \frac{\langle f|\mathcal O|i\rangle}{2\sqrt{E_fE_i}}.
\end{equation}
对应的统计链路还包括 jackknife、bootstrap、协方差和 \(\chi^2\) 拟合；MyQCD
把它们写成可重复执行的有限样本公式，而不是靠经验口头约定。
\src{myqcd/derivations.py；docs/lattice_qcd_gluon_tmd_course/myqcd/lattice_spectroscopy.py}

\subsection{流、重整化与匹配}
梯度流满足
\begin{equation}
 \partial_t B_\mu=D_\nu G_{\nu\mu},
\end{equation}
并在流时间 \(t\) 的窗口内产生平滑半径 \(\sqrt{8t}\) 与 flowed 观测量。
Wilson 线双局域算符的裸矩阵元则必须先处理线性发散与混合：
\begin{equation}
 \mathcal O_\Gamma^{\rm bare}(z,a)
 =e^{-\delta m(a)|z|}\,
 Z_{\Gamma}^{\log}(z,\mu)\,
 \mathcal O_\Gamma^{\rm ren}(z,\mu)
 +\sum_{\Gamma'} Z_{\Gamma\Gamma'}^{\rm mix}(z,a,\mu)\,
 \mathcal O_{\Gamma'}^{\rm ren}(z,\mu).
\end{equation}
RI/MOM 条件、比值方案、混合方案与自重整化都只是这条账本的不同切面：
\begin{equation}
 Z_q^{-1} Z_\Gamma^{\rm RI/MOM}(z,\mu,p^2)\,
 \Tr\!\left[P\,\Lambda_\Gamma(p,z)\right]_{p^2=\mu^2}
 =\Tr\!\left[P\,\Lambda_\Gamma^{\rm tree}(p,z)\right],
\end{equation}
\begin{equation}
 h_R^{\rm hybrid}(z,\mu)=
 \begin{cases}
 h_{\rm bare}(z,P_z)/h_{\rm bare}(z,0),&|z|<z_s,\\
 \eta_s h_{\rm bare}(z,P_z)/Z_R(z,\mu),&|z|\ge z_s,
 \end{cases}
 \quad \eta_s=\frac{Z_R(z_s,\mu)}{h_{\rm bare}(z_s,0)} .
\end{equation}
从短距到长距的匹配则对应
\begin{equation}
 \mathcal M(\nu,z^2)
 =\int_{-1}^{1}\! \dd x\, \e^{\ii x\nu}\,\mathcal P(x,z^2),
 \qquad
 \widetilde q(x,P_z,\mu)
 =\int\frac{\dd z}{2\pi}\,\e^{-\ii x z P_z}\,\widetilde h(z,P_z,\mu),
\end{equation}
以及 TMD 的软减除与 Collins--Soper 演化
\begin{equation}
 F(x,b_T;\mu,\zeta)
 =\frac{Z_F(\mu,\zeta)}{\sqrt{S(b_T;\mu,\delta)}}\,
 F_{\rm unsub}(x,b_T;\delta)
 =\sum_j C_j(b_T,\mu,\zeta)\otimes f_j(x,\mu)
 +O(b_T^2\Lambda_{\rm QCD}^2),
\end{equation}
\begin{equation}
 \frac{\partial \ln F(x,b_T;\mu,\zeta)}{\partial \ln\sqrt{\zeta}}
 =K(b_T,\mu),
 \qquad
 \mu\frac{d}{d\mu}K(b_T,\mu)
 =-\gamma_K(\alpha_s)=-2\Gamma_{\rm cusp}(\alpha_s)
 \end{equation}
这条链的物理含义很硬：smearing 只压 UV，蒸馏只压空间低模，ratio 只压重叠
因子，\(Z_R\)/RI-MOM/hybrid 只处理重整化账，matching 才把欧氏对象接到
光锥或 TMD 语言。
\src{myqcd/derivations.py；docs/lattice_qcd_gluon_tmd_course/myqcd/renormalization_tmd.py；docs/lattice_qcd_gluon_tmd_course/SOURCES.md}

\section{SymPy 辅助推导：只验证代数骨架}
根目录 \texttt{derivations.py} 的 93 项核心检查全部通过。它们覆盖了从
生成泛函、链接变量、SU(3) 代数、谱学、统计、流、EMT，到重整化与
TMD 匹配的完整骨架。具体地说，\texttt{generating\_functional}、
\texttt{lattice\_link}、\texttt{stout\_smearing\_su2}、\texttt{momentum\_smearing\_shift}、
\texttt{quark\_gaussian\_smearing}、\texttt{su3\_generator\_identities}、
\texttt{generalized\_eigenvalue\_problem}、\texttt{delete\_one\_jackknife\_identity}、
\texttt{gradient\_flow\_heat\_kernel}、\texttt{wilson\_line\_linear\_counterterm}、
\texttt{ri\_mom\_and\_ratio\_schemes}、\texttt{soft\_subtraction\_identity}、
\texttt{collins\_soper\_evolution} 和 \texttt{two\_momentum\_cs\_estimator}
等推导，都在有限维或已声明模型下闭合。

这层证据只说明：公式链的代数骨架是自洽的。它不说明真实系综、有限动量、
统计涨落、方案截断和非微扰动力学已经自然消失；那些必须由独立数值数据、
误差预算和外推窗口来处理。
\src{myqcd/README.md；myqcd/derivations.py}

\section{辅助模块的边界：群论、谱学、重整化与 TMD}
课程配套代码只能作为局部代数复现和接口审计，不能提升为系综或物理结果。
其内容可按三组理解：
\begin{itemize}
  \item 群论与 QFT：SU(2)、SU(3)、BCH、U(1) 规范不变性、生成泛函、横向投影；
  \item 格点谱学：色散关系、有限时间关联函数、反宇称伙伴、有效能量、GEVP、
  Lüscher 阈值能移、有限温度、jackknife；
  \item 重整化与 TMD：梯度流热核、平滑半径、Wilson 线线性反项、RI/MOM、
  比值、混合方案、operator mixing、soft subtraction、Collins--Soper 演化、
  two-momentum CS 核估计。
\end{itemize}
若模块返回统一的 \texttt{SymbolicExample}，其字段包括 \texttt{equations}、
\texttt{checks}、\texttt{assumptions}、\texttt{boundary}、\texttt{course\_refs}、
\texttt{source\_refs}。课程层强调讲解顺序和最小闭环，根目录强调全局公式链
和来源边界；这只构成辅助的“可解释 + 可执行”证据，不是对真实格点
系综、连续极限或 TMD 因子化的证明。
\src{docs/lattice_qcd_gluon_tmd_course/myqcd/README.md；docs/lattice_qcd_gluon_tmd_course/myqcd/run_all.py；docs/lattice_qcd_gluon_tmd_course/myqcd/group_qft.py；docs/lattice_qcd_gluon_tmd_course/myqcd/lattice_spectroscopy.py；docs/lattice_qcd_gluon_tmd_course/myqcd/renormalization_tmd.py}

\section{范围修正}
因此，本报告的正确读法应是：\fileline{MyQCD} 是工作对象，\fileline{PyQCD}
是对照对象；前者给出完整公式链、SymPy 验证和课程模块，后者给出物理接口、
管线和实现边界。参考源锚定在 \fileline{refer/books/} 与
\fileline{refer/papers/}，而不是历史上那个被你明确标记为反例的旧
\texttt{LQCD course PPT}。换句话说，这份报告的目标不是“只理解 PyQCD”，
而是把 \fileline{MyQCD} 的理论主线、课程模块、参考源和对照实现统一到同一套
物理语言里。
\src{myqcd/README.md；docs/lattice_qcd_gluon_tmd_course/SOURCES.md；refer/books/AGENTS.md；refer/papers/AGENTS.md}
\end{document}
